% Môn: Vật lý
% Lớp: 10
% Chương: Chương I. Mở đầu
% Bài: L10C1 Bài 1. Làm quen với Vật lí
% Số câu: 281
% App đọc file này trực tiếp — sửa rồi Commit trên GitHub, bấm Đồng bộ GitHub .tex

% L10C1 Bài 1. Làm quen với Vật lí

% ===== Câu 1 =====
% ID: L10C1B1-01-DS
% Mức: NB
\dangbt{Nhận diện và sắp xếp các bước của phương pháp thực nghiệm}
\begin{ex}
% ID: L10C1B1-01-DS
% Mức: NB
Hãy đánh giá tính đúng sai của các nhận định dưới đây về quy trình của phương pháp thực nghiệm:
\choiceTF
{Phương pháp thực nghiệm là phương pháp chỉ dựa hoàn toàn vào suy luận logic toán học mà không cần làm bất kì thí nghiệm thực tế nào.}
{\True Thí nghiệm kiểm tra dự đoán là bước mang tính quyết định để chứng minh hoặc bác bỏ một giả thuyết khoa học.}
{\True Sơ đồ tiến trình tìm hiểu thế giới tự nhiên dưới góc độ vật lí gồm 5 bước cơ bản được thực hiện tuần tự.}
{Quan sát hằng ngày về hạt mưa rơi là bước cuối cùng trong tiến trình thực nghiệm của Galilei về sự rơi tự do.}
\loigiai{
\begin{itemchoice}
\item (Sai) Phương pháp thực nghiệm là phương pháp chỉ dựa hoàn toàn vào suy luận logic toán học mà không cần làm bất kì thí nghiệm thực tế nào. (Vì): Phương pháp thực nghiệm dựa trên việc tiến hành các thí nghiệm thực tế để kiểm chứng các giả thuyết, không chỉ dựa vào suy luận logic toán học
\item (Đúng) Thí nghiệm kiểm tra dự đoán là bước mang tính quyết định để chứng minh hoặc bác bỏ một giả thuyết khoa học. (Vì): Thí nghiệm kiểm tra dự đoán là bước quan trọng nhất để xác nhận hoặc bác bỏ một giả thuyết khoa học đã đề ra
\item (Đúng) Sơ đồ tiến trình tìm hiểu thế giới tự nhiên dưới góc độ vật lí gồm 5 bước cơ bản được thực hiện tuần tự. (Vì): Sơ đồ tiến trình tìm hiểu thế giới tự nhiên dưới góc độ vật lí bao gồm 5 bước cơ bản được thực hiện theo một trình tự nhất định
\item (Sai) Quan sát hằng ngày về hạt mưa rơi là bước cuối cùng trong tiến trình thực nghiệm của Galilei về sự rơi tự do. (Vì): Quan sát hằng ngày về hạt mưa rơi là bước quan sát ban đầu, là tiền đề để xây dựng giả thuyết, không phải là bước cuối cùng trong tiến trình thực nghiệm của Galilei về sự rơi tự do
\end{itemchoice}
}
\end{ex}

% ===== Câu 2 =====
% ID: L10C1B1-01-TL
% Mức: NB
\dangbt{Phân biệt các loại mô hình trong nghiên cứu Vật lí}
\begin{bt}
% ID: L10C1B1-01-TL
% Mức: NB
Trong các loại mô hình vật lí (mô hình vật chất, mô hình lí thuyết, mô hình toán học), việc sử dụng các công thức, phương trình hay đồ thị để mô tả các đặc điểm của đối tượng nghiên cứu thuộc loại mô hình nào?
\loigiai{
Mô hình toán học là mô hình sử dụng các công thức, phương trình, đồ thị, kí hiệu của Toán học để mô tả các đặc điểm của đối tượng nghiên cứu trong Vật lí.
}
\end{bt}

% ===== Câu 3 =====
% ID: L10C1B1-01-TLN
% Mức: NB
\dangbt{Phân tích vai trò và ảnh hưởng hai mặt của Vật lí đối với đời sống, khoa học và công nghệ}
\begin{ex}
% ID: L10C1B1-01-TLN
% Mức: NB
Phát hiện ra hiện tượng cảm ứng điện từ của nhà vật lí nào đã mở đầu cho kỉ nguyên sử dụng điện năng trong sản xuất và đời sống?
\shortans{Faraday}
\loigiai{
Nhà vật lí Michael Faraday đã phát hiện ra hiện tượng cảm ứng điện từ, làm cơ sở cho sự ra đời của máy phát điện và động cơ điện.
}
\end{ex}

% ===== Câu 4 =====
% ID: L10C1B1-01-TN
% Mức: NB
\dangbt{Nhận diện và sắp xếp các bước của phương pháp thực nghiệm}
\begin{ex}
% ID: L10C1B1-01-TN
% Mức: NB
Bước đầu tiên trong phương pháp nghiên cứu vật lí là gì?
\choice
{Đề xuất vấn đề}
{Hình thành giả thuyết}
{Kiểm tra giả thuyết}
{\True Quan sát, suy luận}
\loigiai{
Theo tiến trình nghiên cứu vật lí, bước đầu tiên là \textbf{quan sát, suy luận} để nhận ra các hiện tượng và vấn đề cần nghiên cứu.
}
\end{ex}

% ===== Câu 5 =====
% ID: L10C1B1-02-DS
% Mức: NB
\begin{ex}
% ID: L10C1B1-02-DS
% Mức: NB
Trình tự các bước trong phương pháp nghiên cứu vật lí gồm:
\choiceTF
{\True Quan sát, suy luận}
{\True Đề xuất vấn đề}
{\True Hình thành giả thuyết}
{Rút ra kết luận trước khi kiểm tra giả thuyết}
\loigiai{
\begin{itemchoice}
\item (Đúng) Quan sát, suy luận. (Vì): Là bước khởi đầu quan trọng trong phương pháp nghiên cứu
\item (Đúng) Đề xuất vấn đề. (Vì): Sau quan sát, cần đề xuất vấn đề cần nghiên cứu
\item (Đúng) Hình thành giả thuyết. (Vì): Sau khi có vấn đề, hình thành giả thuyết để kiểm chứng
\item (Sai) Rút ra kết luận trước khi kiểm tra giả thuyết. (Vì): Phải kiểm tra giả thuyết rồi mới được phép rút ra kết luận
\end{itemchoice}
}
\end{ex}

% ===== Câu 6 =====
% ID: L10C1B1-02-TL
% Mức: NB
\dangbt{Phương pháp mô hình trong nghiên cứu Vật lí}
\begin{bt}
% ID: L10C1B1-02-TL
% Mức: NB
Đề xuất cách vận dụng nội dung “Phương pháp mô hình trong nghiên cứu Vật lí” trong một tình huống học tập hoặc thực tiễn.
\loigiai{
Cần xác định tình huống, chọn kiến thức phù hợp, thực hiện theo quy trình, kiểm tra điều kiện và đánh giá kết quả. Nhận diện mô hình vật chất, mô hình toán học và mô hình lí thuyết; xác định phạm vi áp dụng và kiểm chứng.
}
\end{bt}

% ===== Câu 7 =====
% ID: L10C1B1-02-TLN
% Mức: NB
\begin{ex}
% ID: L10C1B1-02-TLN
% Mức: NB
Cho bốn nhận định: (1) Mô hình vật chất có thể là mẫu thu nhỏ. (2) Chất điểm luôn là vật có kích thước bằng không trong thực tế. (3) Có thể điều chỉnh mô hình khi bằng chứng mới xuất hiện. (4) Mô hình không có khả năng dự đoán. Có bao nhiêu nhận định đúng? Chỉ ghi số nguyên.
\shortans{2}
\loigiai{
Nhận định (1), (3) đúng; (2), (4) sai. Vậy có 2 nhận định đúng.
}
\end{ex}

% ===== Câu 8 =====
% ID: L10C1B1-02-TN
% Mức: NB
\dangbt{Nhận diện và sắp xếp các bước của phương pháp thực nghiệm}
\begin{ex}
% ID: L10C1B1-02-TN
% Mức: NB
Trong tiến trình tìm hiểu thế giới tự nhiên dưới góc độ vật lí theo phương pháp thực nghiệm, bước đầu tiên cần thực hiện là gì?
\choice
{Đưa ra dự đoán}
{Thí nghiệm kiểm tra dự đoán}
{\True Xác định vấn đề cần nghiên cứu}
{Quan sát, thu thập thông tin}
\loigiai{
Bước đầu tiên trong tiến trình của phương pháp thực nghiệm luôn luôn là xác định vấn đề cần nghiên cứu.
}
\end{ex}

% ===== Câu 9 =====
% ID: L10C1B1-03-DS
% Mức: NB
\begin{ex}
% ID: L10C1B1-03-DS
% Mức: NB
Trong sơ đồ tiến trình phương pháp thực nghiệm, hãy xác định tính đúng sai của các mệnh đề sau:
\choiceTF
{\True Bước "Quan sát, thu thập thông tin" đứng trước bước "Đưa ra dự đoán".}
{\True Bước "Đưa ra dự đoán" đứng trước bước "Thí nghiệm kiểm tra dự đoán".}
{\True Bước "Thí nghiệm kiểm tra dự đoán" đứng trước bước "Kết luận".}
{Bước "Quan sát, thu thập thông tin" đứng sau bước "Thí nghiệm kiểm tra dự đoán".}
\loigiai{
\begin{itemchoice}
\item (Đúng) Bước "Quan sát, thu thập thông tin" đứng trước bước "Đưa ra dự đoán". (Vì): Mệnh đề
\item (Đúng) Bước "Đưa ra dự đoán" đứng trước bước "Thí nghiệm kiểm tra dự đoán". (Vì): Mệnh đề
\item (Đúng) Bước "Thí nghiệm kiểm tra dự đoán" đứng trước bước "Kết luận". (Vì): Mệnh đề
\item (Sai) Bước "Quan sát, thu thập thông tin" đứng sau bước "Thí nghiệm kiểm tra dự đoán". (Vì): Mệnh đề
\end{itemchoice}
}
\end{ex}

% ===== Câu 10 =====
% ID: L10C1B1-03-TL
% Mức: TH
\dangbt{Phương pháp mô hình trong nghiên cứu Vật lí}
\begin{bt}
% ID: L10C1B1-03-TL
% Mức: TH
Trình bày kiến thức cốt lõi của dạng “Phương pháp mô hình trong nghiên cứu Vật lí” và minh họa bằng một ví dụ.
\loigiai{
Nhận diện mô hình vật chất, mô hình toán học và mô hình lí thuyết; xác định phạm vi áp dụng và kiểm chứng. Ví dụ phù hợp cần nêu rõ đối tượng, dữ kiện và kết luận theo kiến thức SGK.
}
\end{bt}

% ===== Câu 11 =====
% ID: L10C1B1-03-TLN
% Mức: NB
\begin{ex}
% ID: L10C1B1-03-TLN
% Mức: NB
Cho bốn nhận định: (1) Kết quả từ mô hình cần được đối chiếu thực nghiệm. (2) Không được sửa mô hình khi có dữ liệu mới. (3) Mô hình là sự biểu diễn đơn giản hóa đối tượng hoặc quá trình. (4) Mô hình toán học không dùng phương trình. Có bao nhiêu nhận định đúng? Chỉ ghi số nguyên.
\shortans{2}
\loigiai{
Nhận định (1), (3) đúng; (2), (4) sai. Vậy có 2 nhận định đúng.
}
\end{ex}

% ===== Câu 12 =====
% ID: L10C1B1-03-TN
% Mức: NB
\dangbt{Nhận diện và sắp xếp các bước của phương pháp thực nghiệm}
\begin{ex}
% ID: L10C1B1-03-TN
% Mức: NB
Bước nào giúp đánh giá giả thuyết có phù hợp hay không với thực tế?
\choice
{Đề xuất vấn đề}
{Rút ra kết luận}
{Hình thành giả thuyết}
{\True Kiểm tra giả thuyết}
\loigiai{
Kiểm tra giả thuyết là bước giúp xác định giả thuyết có đúng với hiện tượng thực tế hay không.
}
\end{ex}

% ===== Câu 13 =====
% ID: L10C1B1-04-DS
% Mức: NB
\begin{ex}
% ID: L10C1B1-04-DS
% Mức: NB
Phương pháp thực nghiệm gồm:
\choiceTF
{Tính lực tác dụng bằng công thức Newton}
{\True Đo điện áp bằng vôn kế}
{\True Đo độ dài con lắc}
{Mô hình hoá bằng toán}
\loigiai{
\begin{itemchoice}
\item (Sai) Tính lực tác dụng bằng công thức Newton. (Vì): Là lí thuyết
\item (Đúng) Đo điện áp bằng vôn kế. (Vì): Là phép đo cụ thể
\item (Đúng) Đo độ dài con lắc. (Vì): Là hoạt động thực nghiệm
\item (Sai) Mô hình hoá bằng toán. (Vì): Là phương pháp lí thuyết
\end{itemchoice}
}
\end{ex}

% ===== Câu 14 =====
% ID: L10C1B1-04-TL
% Mức: TH
\dangbt{Phương pháp mô hình trong nghiên cứu Vật lí}
\begin{bt}
% ID: L10C1B1-04-TL
% Mức: TH
So sánh hai nhận định: “Mỗi mô hình chỉ phù hợp trong phạm vi nhất định.” và “Mô hình không cần kiểm chứng.”.
\loigiai{
Nhận định thứ nhất đúng: Mỗi mô hình chỉ phù hợp trong phạm vi nhất định. Nhận định thứ hai sai: Mô hình không cần kiểm chứng. Sự khác nhau nằm ở điều kiện và bản chất kiến thức được áp dụng.
}
\end{bt}

% ===== Câu 15 =====
% ID: L10C1B1-04-TLN
% Mức: TH
\begin{ex}
% ID: L10C1B1-04-TLN
% Mức: TH
Cho bốn nhận định: (1) Mô hình toán học dùng phương trình mô tả quan hệ giữa các đại lượng. (2) Mô hình không cần kiểm chứng. (3) Sơ đồ tia sáng là một dạng mô hình. (4) Một đối tượng chỉ có duy nhất một mô hình. Có bao nhiêu nhận định đúng? Chỉ ghi số nguyên.
\shortans{2}
\loigiai{
Nhận định (1), (3) đúng; (2), (4) sai. Vậy có 2 nhận định đúng.
}
\end{ex}

% ===== Câu 16 =====
% ID: L10C1B1-04-TN
% Mức: NB
\dangbt{Nhận diện và sắp xếp các bước của phương pháp thực nghiệm}
\begin{ex}
% ID: L10C1B1-04-TN
% Mức: NB
Bước nào trong tiến trình nghiên cứu khoa học giúp xác định vấn đề cần tìm hiểu?
\choice
{Kiểm tra giả thuyết}
{\True Đề xuất vấn đề}
{Rút ra kết luận}
{Hình thành giả thuyết}
\loigiai{
Sau khi quan sát, cần đề xuất vấn đề — đây là bước xác định rõ chủ đề cần nghiên cứu.
}
\end{ex}

% ===== Câu 17 =====
% ID: L10C1B1-05-DS
% Mức: TH
\begin{ex}
% ID: L10C1B1-05-DS
% Mức: TH
Trong nghiên cứu vật lí:
\choiceTF
{\True Phương pháp lí thuyết cần được kiểm chứng bằng thực nghiệm}
{Phương pháp thực nghiệm không cần phương pháp lí thuyết}
{\True Phương pháp thực nghiệm có vai trò quyết định}
{Chỉ cần thực nghiệm là đủ để chứng minh giả thuyết}
\loigiai{
\begin{itemchoice}
\item (Đúng) Phương pháp lí thuyết cần được kiểm chứng bằng thực nghiệm. (Vì): Lí thuyết cần được thực nghiệm xác minh
\item (Sai) Phương pháp thực nghiệm không cần phương pháp lí thuyết. (Vì): Hai phương pháp luôn hỗ trợ nhau
\item (Đúng) Phương pháp thực nghiệm có vai trò quyết định. (Vì): Thực nghiệm là căn cứ kiểm chứng giả thuyết
\item (Sai) Chỉ cần thực nghiệm là đủ để chứng minh giả thuyết. (Vì): Cần có lí thuyết hỗ trợ thiết kế thí nghiệm hợp lí
\end{itemchoice}
}
\end{ex}

% ===== Câu 18 =====
% ID: L10C1B1-05-TL
% Mức: VD
\dangbt{Phương pháp mô hình trong nghiên cứu Vật lí}
\begin{bt}
% ID: L10C1B1-05-TL
% Mức: VD
Phân tích phát biểu “Mô hình vật chất có thể là mẫu thu nhỏ.” và giải thích vì sao phát biểu đó đúng.
\loigiai{
Phát biểu đúng. Mô hình vật chất có thể là mẫu thu nhỏ. Nội dung này phù hợp với kiến thức trọng tâm của dạng bài.
}
\end{bt}

% ===== Câu 19 =====
% ID: L10C1B1-05-TLN
% Mức: VD
\begin{ex}
% ID: L10C1B1-05-TLN
% Mức: VD
Cho bốn nhận định: (1) Mô hình lí thuyết giúp giải thích và dự đoán. (2) Sơ đồ không thể là mô hình. (3) Chất điểm là mô hình dùng khi kích thước vật có thể bỏ qua. (4) Mô hình phải giống hoàn toàn vật thật ở mọi chi tiết. Có bao nhiêu nhận định đúng? Chỉ ghi số nguyên.
\shortans{2}
\loigiai{
Nhận định (1), (3) đúng; (2), (4) sai. Vậy có 2 nhận định đúng.
}
\end{ex}

% ===== Câu 20 =====
% ID: L10C1B1-05-TN
% Mức: NB
\dangbt{Nhận diện và sắp xếp các bước của phương pháp thực nghiệm}
\begin{ex}
% ID: L10C1B1-05-TN
% Mức: NB
Khi một dự đoán khoa học bị kết quả thí nghiệm kiểm tra bác bỏ, nhà nghiên cứu cần thực hiện bước nào tiếp theo theo tiến trình phương pháp thực nghiệm?
\choice
{Rút ra kết luận ngay lập tức}
{Dừng việc nghiên cứu và không đưa ra kết luận}
{\True Đưa ra dự đoán mới hoặc xác định lại vấn đề cần nghiên cứu}
{Giữ nguyên dự đoán cũ và công bố kết quả}
\loigiai{
Nếu thí nghiệm chứng tỏ dự đoán là sai thì người ta phải đưa ra dự đoán mới và làm thí nghiệm kiểm tra dự đoán này, hoặc xác định lại vấn đề cần nghiên cứu.
}
\end{ex}

% ===== Câu 21 =====
% ID: L10C1B1-06-DS
% Mức: TH
\begin{ex}
% ID: L10C1B1-06-DS
% Mức: TH
Nhận xét các nhận định sau về phương pháp nghiên cứu Vật lí
\choiceTF
{Nghiên cứu Vật lí chỉ sử dụng phương pháp thực nghiệm.}
{\True Dự đoán khoa học có cơ sở dựa trên các yếu tố quan sát, trải nghiệm thực tế, các kiến thức đã có liên quan đến dự đoán.}
{Phương pháp thực nghiệm có 5 bước chính bao gồm xác định vấn đề cần nghiên cứu, quan sát, thí nghiệm, kết luận, dự đoán}
{\True Galileo Galilei là cha đẻ của phương pháp thực nghiệm nhờ thí nghiệm bác bỏ nhận định “vật nặng rơi nhanh hơn vật nhẹ”.}
\loigiai{
\begin{itemchoice}
\item (Sai) Nghiên cứu Vật lí chỉ sử dụng phương pháp thực nghiệm.
\item (Đúng) Dự đoán khoa học có cơ sở dựa trên các yếu tố quan sát, trải nghiệm thực tế, các kiến thức đã có liên quan đến dự đoán.
\item (Sai) Phương pháp thực nghiệm có 5 bước chính bao gồm xác định vấn đề cần nghiên cứu, quan sát, thí nghiệm, kết luận, dự đoán.
\item (Đúng) Galileo Galilei là cha đẻ của phương pháp thực nghiệm nhờ thí nghiệm bác bỏ nhận định “vật nặng rơi nhanh hơn vật nhẹ”.
\end{itemchoice}
}
\end{ex}

% ===== Câu 22 =====
% ID: L10C1B1-06-TL
% Mức: VD
\dangbt{Phương pháp mô hình trong nghiên cứu Vật lí}
\begin{bt}
% ID: L10C1B1-06-TL
% Mức: VD
Xây dựng một quy trình ngắn để giải quyết nhiệm vụ thuộc dạng “Phương pháp mô hình trong nghiên cứu Vật lí”.
\loigiai{
Quy trình: đọc yêu cầu; xác định khái niệm hoặc đại lượng; chọn nguyên tắc/công thức; xử lí dữ kiện; kiểm tra đơn vị và điều kiện; kết luận. Nhận diện mô hình vật chất, mô hình toán học và mô hình lí thuyết; xác định phạm vi áp dụng và kiểm chứng.
}
\end{bt}

% ===== Câu 23 =====
% ID: L10C1B1-06-TLN
% Mức: VDC
\begin{ex}
% ID: L10C1B1-06-TLN
% Mức: VDC
Cho bốn nhận định: (1) Mô hình là sự biểu diễn đơn giản hóa đối tượng hoặc quá trình. (2) Mô hình toán học không dùng phương trình. (3) Kết quả từ mô hình cần được đối chiếu thực nghiệm. (4) Không được sửa mô hình khi có dữ liệu mới. Có bao nhiêu nhận định đúng? Chỉ ghi số nguyên.
\shortans{2}
\loigiai{
Nhận định (1), (3) đúng; (2), (4) sai. Vậy có 2 nhận định đúng.
}
\end{ex}

% ===== Câu 24 =====
% ID: L10C1B1-06-TN
% Mức: NB
\dangbt{Nhận diện và sắp xếp các bước của phương pháp thực nghiệm}
\begin{ex}
% ID: L10C1B1-06-TN
% Mức: NB
Các nhà vật lí dùng phương pháp thực nghiệm để tìm hiểu thế giới tự nhiên trong khoảng thời gian nào?
\choice
{\True Từ thế kỉ XVII đến cuối thế kỉ XIX}
{Từ năm $1900$ đến nay}
{Từ năm $350$ TCN đến năm $1831$}
{Từ cuối thế kỉ XIX đến nay}
\end{ex}

% ===== Câu 25 =====
% ID: L10C1B1-07-DS
% Mức: TH
\begin{ex}
% ID: L10C1B1-07-DS
% Mức: TH
Khi nghiên cứu một hiện tượng vật lí, việc áp dụng phương pháp nghiên cứu khoa học thường bao gồm các bước cơ bản.
\choiceTF
{\True Quan sát hiện tượng và đặt ra giả thuyết là những bước đầu tiên trong phương pháp nghiên cứu khoa học.}
{Việc kiểm tra giả thuyết bằng thực nghiệm thường được bỏ qua nếu giả thuyết đó có vẻ hợp lý.}
{\True Phân tích kết quả thực nghiệm và rút ra kết luận là bước cần thiết để xác nhận hoặc bác bỏ giả thuyết ban đầu.}
{Một lý thuyết khoa học đã được công nhận là chân lý tuyệt đối và không thể bị thay đổi bởi các bằng chứng mới.}
\loigiai{
\begin{itemchoice}
\item (Đúng) Quan sát hiện tượng và đặt ra giả thuyết là những bước đầu tiên trong phương pháp nghiên cứu khoa học. (Vì): Quan sát và thu thập dữ liệu là nền tảng để nhận diện vấn đề và hình thành câu hỏi nghiên cứu
\item (Sai) Việc kiểm tra giả thuyết bằng thực nghiệm thường được bỏ qua nếu giả thuyết đó có vẻ hợp lý. (Vì): Giả thuyết khoa học cần có cơ sở từ quan sát, kiến thức đã biết và phải có khả năng kiểm chứng được bằng thực nghiệm
\item (Đúng) Phân tích kết quả thực nghiệm và rút ra kết luận là bước cần thiết để xác nhận hoặc bác bỏ giả thuyết ban đầu. (Vì): Thực nghiệm giúp kiểm tra giả thuyết trong điều kiện kiểm soát, cung cấp bằng chứng cụ thể để xác nhận hoặc bác bỏ giả thuyết
\item (Sai) Một lý thuyết khoa học đã được công nhận là chân lý tuyệt đối và không thể bị thay đổi bởi các bằng chứng mới. (Vì): Một lí thuyết vật lí chỉ được chấp nhận rộng rãi khi nó đã được kiểm chứng nhiều lần bằng thực nghiệm và quan sát, đồng thời có khả năng giải thích và dự đoán các hiện tượng
\end{itemchoice}
}
\end{ex}

% ===== Câu 26 =====
% ID: L10C1B1-07-TL
% Mức: VDC
\dangbt{Phương pháp mô hình trong nghiên cứu Vật lí}
\begin{bt}
% ID: L10C1B1-07-TL
% Mức: VDC
Phân tích phát biểu “Mô hình toán học không dùng phương trình.”, chỉ ra sai sót và sửa lại.
\loigiai{
Phát biểu sai: Mô hình toán học không dùng phương trình. Cần sửa thành: Mô hình toán học dùng phương trình mô tả quan hệ giữa các đại lượng. Nhận diện mô hình vật chất, mô hình toán học và mô hình lí thuyết; xác định phạm vi áp dụng và kiểm chứng.
}
\end{bt}

% ===== Câu 27 =====
% ID: L10C1B1-07-TLN
% Mức: VDC
\begin{ex}
% ID: L10C1B1-07-TLN
% Mức: VDC
Cho bốn nhận định: (1) Mỗi mô hình chỉ phù hợp trong phạm vi nhất định. (2) Mô hình càng nhiều chi tiết thừa càng tốt. (3) Mô hình giúp tập trung vào đặc trưng quan trọng. (4) Mọi mô hình đều đúng trong mọi điều kiện. Có bao nhiêu nhận định đúng? Chỉ ghi số nguyên.
\shortans{2}
\loigiai{
Nhận định (1), (3) đúng; (2), (4) sai. Vậy có 2 nhận định đúng.
}
\end{ex}

% ===== Câu 28 =====
% ID: L10C1B1-07-TN
% Mức: NB
\dangbt{Nhận diện và sắp xếp các bước của phương pháp thực nghiệm}
\begin{ex}
% ID: L10C1B1-07-TN
% Mức: NB
Trình tự nào sau đây thể hiện đúng tiến trình phương pháp nghiên cứu khoa học trong Vật lí?
\choice
{Đề xuất vấn đề → Quan sát → Kiểm tra giả thuyết → Hình thành giả thuyết → Rút ra kết luận}
{Quan sát → Kiểm tra giả thuyết → Hình thành giả thuyết → Đề xuất vấn đề → Rút ra kết luận}
{\True Quan sát → Đề xuất vấn đề → Hình thành giả thuyết → Kiểm tra giả thuyết → Rút ra kết luận}
{Hình thành giả thuyết → Đề xuất vấn đề → Quan sát → Kiểm tra giả thuyết → Rút ra kết luận}
\loigiai{
Tiến trình đúng là:
 
 
• Quan sát, suy luận
 
• Đề xuất vấn đề
 
• Hình thành giả thuyết
 
• Kiểm tra giả thuyết
 
• Rút ra kết luận
}
\end{ex}

% ===== Câu 29 =====
% ID: L10C1B1-08-DS
% Mức: NB
\dangbt{Phân biệt các loại mô hình trong nghiên cứu Vật lí}
\begin{ex}
% ID: L10C1B1-08-DS
% Mức: NB
Phương pháp lí thuyết trong Vật lí:
\choiceTF
{\True Tính toán quỹ đạo bằng công thức}
{Đo nhiệt độ bằng nhiệt kế}
{\True Mô hình hoá va chạm bằng toán học}
{Đo chiều dài bằng thước}
\loigiai{
\begin{itemchoice}
\item (Đúng) Tính toán quỹ đạo bằng công thức. (Vì): Là phương pháp lí thuyết
\item (Sai) Đo nhiệt độ bằng nhiệt kế. (Vì): Là thực nghiệm
\item (Đúng) Mô hình hoá va chạm bằng toán học. (Vì): ùng mô hình toán
\item (Sai) Đo chiều dài bằng thước. (Vì): Là phép đo thực nghiệm
\end{itemchoice}
}
\end{ex}

% ===== Câu 30 =====
% ID: L10C1B1-08-TL
% Mức: NB
\dangbt{Phương pháp thực nghiệm trong nghiên cứu Vật lí}
\begin{bt}
% ID: L10C1B1-08-TL
% Mức: NB
Đề xuất cách vận dụng nội dung “Phương pháp thực nghiệm trong nghiên cứu Vật lí” trong một tình huống học tập hoặc thực tiễn.
\loigiai{
Cần xác định tình huống, chọn kiến thức phù hợp, thực hiện theo quy trình, kiểm tra điều kiện và đánh giá kết quả. Tiến trình cơ bản: xác định vấn đề, dự đoán, thiết kế và thực hiện thí nghiệm, xử lí dữ liệu, kết luận và điều chỉnh.
}
\end{bt}

% ===== Câu 31 =====
% ID: L10C1B1-08-TLN
% Mức: NB
\begin{ex}
% ID: L10C1B1-08-TLN
% Mức: NB
Cho bốn nhận định: (1) Bước đầu cần xác định vấn đề nghiên cứu. (2) Không cần xác định vấn đề trước khi đo. (3) Biến độc lập là đại lượng người nghiên cứu chủ động thay đổi. (4) Giả thuyết không cần kiểm tra. Có bao nhiêu nhận định đúng? Chỉ ghi số nguyên.
\shortans{2}
\loigiai{
Nhận định (1), (3) đúng; (2), (4) sai. Vậy có 2 nhận định đúng.
}
\end{ex}

% ===== Câu 32 =====
% ID: L10C1B1-08-TN
% Mức: NB
\dangbt{Nhận diện và sắp xếp các bước của phương pháp thực nghiệm}
\begin{ex}
% ID: L10C1B1-08-TN
% Mức: NB
Trong phương pháp nghiên cứu khoa học, sau bước "Hình thành giả thuyết" là bước nào?
\choice
{Đề xuất vấn đề}
{Rút ra kết luận}
{\True Kiểm tra giả thuyết}
{Quan sát, suy luận}
\loigiai{
Theo đúng tiến trình, sau khi hình thành giả thuyết là bước kiểm tra giả thuyết.
}
\end{ex}

% ===== Câu 33 =====
% ID: L10C1B1-09-DS
% Mức: NB
\dangbt{Phân biệt các loại mô hình trong nghiên cứu Vật lí}
\begin{ex}
% ID: L10C1B1-09-DS
% Mức: NB
Khi nghiên cứu về các loại mô hình được sử dụng trong học tập và nghiên cứu Vật lí, hãy đánh giá tính đúng sai của các nhận định dưới đây:
\choiceTF
{\True Công thức tính quãng đường chuyển động thẳng đều $s = vt$ là một ví dụ về mô hình toán học.}
{Khái niệm tia sáng dùng để biểu diễn đường truyền ánh sáng là mô hình vật chất vì mắt người có thể nhìn thấy tia sáng.}
{\True Mô hình cấu trúc hạt nhân nguyên tử đặt trong phòng thí nghiệm thuộc loại mô hình vật chất.}
{Mọi mô hình lí thuyết đều phải giữ lại toàn bộ các chi tiết cấu tạo phức tạp của vật thật để đảm bảo tính chính xác.}
\loigiai{
\begin{itemchoice}
\item (Đúng) Công thức tính quãng đường chuyển động thẳng đều $s = vt$ là một ví dụ về mô hình toán học. (Vì): Công thức tính quãng đường chuyển động thẳng đều $s = vt$ sử dụng các kí hiệu và phương trình toán học để mô tả đặc điểm của chuyển động, đây là một ví dụ điển hình của mô hình toán học
\item (Sai) Khái niệm tia sáng dùng để biểu diễn đường truyền ánh sáng là mô hình vật chất vì mắt người có thể nhìn thấy tia sáng. (Vì): Khái niệm tia sáng là một sự lí tưởng hóa, dùng để biểu diễn đường truyền ánh sáng trong các bài toán quang học, nó không phải là mô hình vật chất vì mắt người không thể nhìn thấy tia sáng
\item (Đúng) Mô hình cấu trúc hạt nhân nguyên tử đặt trong phòng thí nghiệm thuộc loại mô hình vật chất. (Vì): Mô hình cấu trúc hạt nhân nguyên tử đặt trong phòng thí nghiệm, ví dụ như các mô hình lắp ghép, là một vật thể vật chất được tạo ra để mô phỏng cấu trúc của hạt nhân nguyên tử, do đó thuộc loại mô hình vật chất
\item (Sai) Mọi mô hình lí thuyết đều phải giữ lại toàn bộ các chi tiết cấu tạo phức tạp của vật thật để đảm bảo tính chính xác. (Vì): Mô hình lí thuyết là sự trừu tượng hóa và đơn giản hóa đối tượng nghiên cứu, nó thường bỏ qua các chi tiết phức tạp không cần thiết để làm nổi bật các đặc điểm cốt lõi, chứ không giữ lại toàn bộ
\end{itemchoice}
}
\end{ex}

% ===== Câu 34 =====
% ID: L10C1B1-09-TL
% Mức: TH
\dangbt{Phương pháp thực nghiệm trong nghiên cứu Vật lí}
\begin{bt}
% ID: L10C1B1-09-TL
% Mức: TH
Trình bày kiến thức cốt lõi của dạng “Phương pháp thực nghiệm trong nghiên cứu Vật lí” và minh họa bằng một ví dụ.
\loigiai{
Tiến trình cơ bản: xác định vấn đề, dự đoán, thiết kế và thực hiện thí nghiệm, xử lí dữ liệu, kết luận và điều chỉnh. Ví dụ phù hợp cần nêu rõ đối tượng, dữ kiện và kết luận theo kiến thức SGK.
}
\end{bt}

% ===== Câu 35 =====
% ID: L10C1B1-09-TLN
% Mức: NB
\begin{ex}
% ID: L10C1B1-09-TLN
% Mức: NB
Cho bốn nhận định: (1) Nếu kết quả không phù hợp, cần xem lại giả thuyết hoặc phương án. (2) Không cần ghi chép số liệu. (3) Phương pháp thực nghiệm sử dụng thí nghiệm để kiểm tra dự đoán. (4) Kết quả trái dự đoán phải bị bỏ đi. Có bao nhiêu nhận định đúng? Chỉ ghi số nguyên.
\shortans{2}
\loigiai{
Nhận định (1), (3) đúng; (2), (4) sai. Vậy có 2 nhận định đúng.
}
\end{ex}

% ===== Câu 36 =====
% ID: L10C1B1-09-TN
% Mức: NB
\dangbt{Nhận diện và sắp xếp các bước của phương pháp thực nghiệm}
\begin{ex}
% ID: L10C1B1-09-TN
% Mức: NB
Bước đầu tiên và cơ bản nhất trong phương pháp nghiên cứu thực nghiệm vật lí là gì?
\choice
{\True Quan sát hiện tượng, thu thập thông tin và đặt câu hỏi.}
{Đề xuất giả thuyết để giải thích hiện tượng.}
{Tiến hành thí nghiệm để kiểm chứng giả thuyết.}
{Rút ra kết luận và xây dựng lí thuyết.}
\loigiai{
Phương pháp nghiên cứu thực nghiệm vật lí thường bắt đầu bằng việc quan sát kĩ lưỡng một hiện tượng tự nhiên hoặc một quá trình nào đó, sau đó thu thập các thông tin liên quan và đặt ra những câu hỏi về hiện tượng đó. Đây là bước khởi đầu để định hình vấn đề cần nghiên cứu.
}
\end{ex}

% ===== Câu 37 =====
% ID: L10C1B1-10-DS
% Mức: NB
\dangbt{Phân tích vai trò và ảnh hưởng hai mặt của Vật lí đối với đời sống, khoa học và công nghệ}
\begin{ex}
% ID: L10C1B1-10-DS
% Mức: NB
Nhận xét các nhận định sau về quá trình phát triển của Vật lí đối với các cuộc cách mạng công nghiệp.
\choiceTF
{Đặc trưng của cuộc cách mạng công nghiệp lần thứ nhất là sự xuất hiện các thiết bị dùng điện trong mọi lĩnh vực sản xuất và đời sống con người}
{\True Nghiên cứu về hiện tượng cảm ứng điện từ mở đầu cho kỉ nguyên sử dụng điện năng là cơ sở cho sự ra đời cuộc cách mạng công nghiệp lần thứ hai}
{\True Qua cuộc cách mạng công nghiệp lần thứ ba với đặc trưng là tự động hóa các quá trình sản xuất kéo theo các thành tựu về cách mạng kỹ thuật số}
{Cuộc cách mạng công nghiệp lần thứ tư hay còn gọi là cách mạng công nghệ 4.0 chưa có nhiều đóng góp vượt bậc so với các cuộc cánh mạng trước đó}
\loigiai{
\begin{itemchoice}
\item (Sai) Đặc trưng của cuộc cách mạng công nghiệp lần thứ nhất là sự xuất hiện các thiết bị dùng điện trong mọi lĩnh vực sản xuất và đời sống con người.
\item (Đúng) Nghiên cứu về hiện tượng cảm ứng điện từ mở đầu cho kỉ nguyên sử dụng điện năng là cơ sở cho sự ra đời cuộc cách mạng công nghiệp lần thứ hai.
\item (Đúng) Qua cuộc cách mạng công nghiệp lần thứ ba với đặc trưng là tự động hóa các quá trình sản xuất kéo theo các thành tựu về cách mạng kỹ thuật số.
\item (Sai) Cuộc cách mạng công nghiệp lần thứ tư hay còn gọi là cách mạng công nghệ 4.0 chưa có nhiều đóng góp vượt bậc so với các cuộc cánh mạng trước đó.
\end{itemchoice}
}
\end{ex}

% ===== Câu 38 =====
% ID: L10C1B1-10-TL
% Mức: TH
\dangbt{Phương pháp thực nghiệm trong nghiên cứu Vật lí}
\begin{bt}
% ID: L10C1B1-10-TL
% Mức: TH
So sánh hai nhận định: “Kết quả đo cần được xử lí trước khi kết luận.” và “Biến phụ thuộc luôn do người làm tùy ý đặt giá trị.”.
\loigiai{
Nhận định thứ nhất đúng: Kết quả đo cần được xử lí trước khi kết luận. Nhận định thứ hai sai: Biến phụ thuộc luôn do người làm tùy ý đặt giá trị. Sự khác nhau nằm ở điều kiện và bản chất kiến thức được áp dụng.
}
\end{bt}

% ===== Câu 39 =====
% ID: L10C1B1-10-TLN
% Mức: TH
\begin{ex}
% ID: L10C1B1-10-TLN
% Mức: TH
Cho bốn nhận định: (1) Sau khi nêu vấn đề có thể đề xuất giả thuyết hoặc dự đoán. (2) Biến phụ thuộc luôn do người làm tùy ý đặt giá trị. (3) Biến phụ thuộc là đại lượng được theo dõi theo biến độc lập. (4) Thí nghiệm không cần lặp lại. Có bao nhiêu nhận định đúng? Chỉ ghi số nguyên.
\shortans{2}
\loigiai{
Nhận định (1), (3) đúng; (2), (4) sai. Vậy có 2 nhận định đúng.
}
\end{ex}

% ===== Câu 40 =====
% ID: L10C1B1-10-TN
% Mức: TH
\dangbt{Nhận diện và sắp xếp các bước của phương pháp thực nghiệm}
\begin{ex}
% ID: L10C1B1-10-TN
% Mức: TH
Các hiện tượng nào sau đây có liên quan đến phương pháp lí thuyết:
\choice
{\True Tính toán quỹ đạo chuyển động của Thiên vương tinh dựa vào toán học}
{Thả rơi 1 vật từ trên cao xuống}
{Kiểm tra sự thay đổi nhiệt độ trong quá trình nóng chảy hoặc bay hơi của một chất}
{Ném một quả bóng lên cao}
\loigiai{
Phương pháp lí thuyết sử dụng công cụ toán học để mô hình hóa và dự đoán các hiện tượng vật lí.
}
\end{ex}

% ===== Câu 41 =====
% ID: L10C1B1-100-TN
% Mức: NB
\dangbt{Vận dụng tiến trình của phương pháp mô hình}
\begin{ex}
% ID: L10C1B1-100-TN
% Mức: NB
Trong phương pháp mô hình, khi kiểm tra sự phù hợp của mô hình với thực tế mà cho kết quả chưa phù hợp thì nhà nghiên cứu cần thực hiện hành động nào tiếp theo?
\choice
{Kết luận ngay mô hình đó là chính xác}
{\True Điều chỉnh mô hình và kiểm tra lại}
{Hủy bỏ toàn bộ quá trình nghiên cứu}
{Giữ nguyên mô hình và không cần thay đổi gì}
\loigiai{
Nếu kiểm tra tính phù hợp cho thấy mô hình chưa khớp với thực tế khách quan, nhà nghiên cứu cần thực hiện điều chỉnh mô hình và quay lại bước kiểm tra tính phù hợp cho đến khi đạt độ tương thích tốt.
}
\end{ex}

% ===== Câu 42 =====
% ID: L10C1B1-101-TN
% Mức: NB
\begin{ex}
% ID: L10C1B1-101-TN
% Mức: NB
Mục tiêu chính của việc nghiên cứu Vật lí là gì?
\choice
{\True Tìm hiểu thế giới tự nhiên}
{Phát minh ra các loại máy móc}
{Chế tạo vũ khí hiện đại}
{Tạo ra các sản phẩm tiêu dùng}
\loigiai{
Mục tiêu cốt lõi của Vật lí là hiểu rõ bản chất của thế giới tự nhiên xung quanh chúng ta thông qua việc quan sát, mô tả, giải thích các hiện tượng và tìm ra các quy luật vật lí phổ quát. Việc chế tạo thiết bị công nghệ là ứng dụng của Vật lí, phân loại vật chất có thể liên quan đến Hóa học hoặc Thiên văn học, còn nghiên cứu lịch sử văn minh thuộc lĩnh vực Lịch sử.
}
\end{ex}

% ===== Câu 43 =====
% ID: L10C1B1-102-TN
% Mức: NB
\begin{ex}
% ID: L10C1B1-102-TN
% Mức: NB
Có bao nhiêu bước trong phương pháp thực nghiệm ?
\choice
{\True 5}
{4}
{3}
{1}
\loigiai{
Phương pháp thực nghiệm thường gồm năm bước: xác định vấn đề, quan sát và thu thập thông tin, đưa ra dự đoán, thí nghiệm kiểm tra, rút ra kết luận.
}
\end{ex}

% ===== Câu 44 =====
% ID: L10C1B1-103-TN
% Mức: NB
\begin{ex}
% ID: L10C1B1-103-TN
% Mức: NB
Hoạt động y tế nào dưới đây không sử dụng các thành tựu của vật lí ?
\choice
{\True Lấy thuốc theo đơn}
{Xạ trị}
{Chụp $X$ - quang}
{Chữa tật khúc xạ của mắt bằng laze}
\loigiai{
Xạ trị, chụp $X$ - quang và chữa tật khúc xạ bằng laze đều ứng dụng các thành tựu của Vật lí. Lấy thuốc theo đơn không phải hoạt động sử dụng trực tiếp thành tựu vật lí.
}
\end{ex}

% ===== Câu 45 =====
% ID: L10C1B1-104-TN
% Mức: NB
\begin{ex}
% ID: L10C1B1-104-TN
% Mức: NB
Trong quá trình phát triển của vật lí học, vai trò của toán học được thể hiện rõ nhất ở khía cạnh nào?
\choice
{Giúp các nhà vật lí ghi chép lại kết quả thí nghiệm một cách nhanh chóng và gọn gàng}
{Cung cấp ngôn ngữ để mô tả các hiện tượng tự nhiên một cách định tính, không cần số liệu cụ thể}
{\True Là công cụ thiết yếu để xây dựng các mô hình lí thuyết, dự đoán các hiện tượng mới và kiểm chứng định luật một cách định lượng}
{Chỉ được sử dụng để thực hiện các phép tính đơn giản sau khi đã có kết quả cuối cùng từ thực nghiệm}
\loigiai{
Toán học đóng vai trò cực kì quan trọng trong vật lí học. Nó không chỉ là ngôn ngữ để mô tả các quy luật tự nhiên một cách chính xác và định lượng, mà còn là công cụ mạnh mẽ để xây dựng các mô hình lí thuyết, từ đó đưa ra các dự đoán về các hiện tượng chưa được quan sát. Những dự đoán này sau đó sẽ được kiểm chứng bằng thực nghiệm. Sự phù hợp giữa dự đoán toán học và kết quả thực nghiệm là yếu tố then chốt để xác nhận tính đúng đắn của một lí thuyết vật lí.
}
\end{ex}

% ===== Câu 46 =====
% ID: L10C1B1-105-TN
% Mức: NB
\begin{ex}
% ID: L10C1B1-105-TN
% Mức: NB
Máy hơi nước ra đời trong cuộc cách mạng công nghiệp lần thứ mấy ?
\choice
{\True Lần thứ nhất}
{Lần thứ ba}
{Lần thứ hai}
{Lần thứ tư}
\loigiai{
Máy hơi nước là thành tựu tiêu biểu của cuộc cách mạng công nghiệp lần thứ nhất, giúp thay thế sức lao động thủ công bằng máy móc.
}
\end{ex}

% ===== Câu 47 =====
% ID: L10C1B1-106-TN
% Mức: NB
\dangbt{Vận dụng tiến trình nghiên cứu Vật lí vào tình huống thực tiễn}
\begin{ex}
% ID: L10C1B1-106-TN
% Mức: NB
Trong dạng “Vận dụng tiến trình nghiên cứu Vật lí vào tình huống thực tiễn”, phát biểu nào sau đây đúng?
\choice
{\True Câu hỏi nghiên cứu cần rõ ràng và có thể kiểm tra.}
{Câu hỏi càng mơ hồ càng dễ nghiên cứu.}
{Không cần ghi đơn vị trong bảng số liệu.}
{Sai số không cần được xem xét.}
\loigiai{
Phát biểu đúng là: Câu hỏi nghiên cứu cần rõ ràng và có thể kiểm tra. Vận dụng đầy đủ chu trình nghiên cứu vào vấn đề gần gũi, từ câu hỏi đến đánh giá và cải tiến.
}
\end{ex}

% ===== Câu 48 =====
% ID: L10C1B1-107-TN
% Mức: NB
\begin{ex}
% ID: L10C1B1-107-TN
% Mức: NB
Trong dạng “Vận dụng tiến trình nghiên cứu Vật lí vào tình huống thực tiễn”, phát biểu nào sau đây đúng?
\choice
{\True Đồ thị có thể làm rõ quan hệ giữa hai đại lượng.}
{Đồ thị không thể hiện quan hệ đại lượng.}
{Có thể chọn bỏ dữ liệu tùy ý.}
{Câu hỏi càng mơ hồ càng dễ nghiên cứu.}
\loigiai{
Phát biểu đúng là: Đồ thị có thể làm rõ quan hệ giữa hai đại lượng. Vận dụng đầy đủ chu trình nghiên cứu vào vấn đề gần gũi, từ câu hỏi đến đánh giá và cải tiến.
}
\end{ex}

% ===== Câu 49 =====
% ID: L10C1B1-108-TN
% Mức: NB
\begin{ex}
% ID: L10C1B1-108-TN
% Mức: NB
Trong dạng “Vận dụng tiến trình nghiên cứu Vật lí vào tình huống thực tiễn”, phát biểu nào sau đây đúng?
\choice
{\True Điều kiện an toàn phải được đưa vào kế hoạch.}
{Kế hoạch không cần dụng cụ.}
{Kết luận có thể trái dữ liệu mà không cần giải thích.}
{Đồ thị không thể hiện quan hệ đại lượng.}
\loigiai{
Phát biểu đúng là: Điều kiện an toàn phải được đưa vào kế hoạch. Vận dụng đầy đủ chu trình nghiên cứu vào vấn đề gần gũi, từ câu hỏi đến đánh giá và cải tiến.
}
\end{ex}

% ===== Câu 50 =====
% ID: L10C1B1-109-TN
% Mức: TH
\begin{ex}
% ID: L10C1B1-109-TN
% Mức: TH
Trong dạng “Vận dụng tiến trình nghiên cứu Vật lí vào tình huống thực tiễn”, phát biểu nào sau đây đúng?
\choice
{Kết luận có thể trái dữ liệu mà không cần giải thích.}
{\True Dự đoán nên dựa trên kiến thức hoặc quan sát ban đầu.}
{Đồ thị không thể hiện quan hệ đại lượng.}
{Có thể chọn bỏ dữ liệu tùy ý.}
\loigiai{
Phát biểu đúng là: Dự đoán nên dựa trên kiến thức hoặc quan sát ban đầu. Vận dụng đầy đủ chu trình nghiên cứu vào vấn đề gần gũi, từ câu hỏi đến đánh giá và cải tiến.
}
\end{ex}

% ===== Câu 51 =====
% ID: L10C1B1-11-DS
% Mức: NB
\dangbt{Phân tích vai trò và ảnh hưởng hai mặt của Vật lí đối với đời sống, khoa học và công nghệ}
\begin{ex}
% ID: L10C1B1-11-DS
% Mức: NB
Hãy đánh giá tính đúng sai của các nhận định dưới đây về vai trò của Vật lí đối với các cuộc cách mạng công nghiệp:
\choiceTF
{\True Cuộc cách mạng công nghiệp lần thứ hai gắn liền với việc phát hiện ra hiện tượng cảm ứng điện từ của Faraday.}
{Tự động hóa sản xuất trong cuộc cách mạng công nghiệp lần thứ ba dựa trên các nghiên cứu về động cơ hơi nước.}
{\True Trí tuệ nhân tạo, robot và internet vạn vật là những công nghệ đặc trưng của cuộc cách mạng công nghiệp lần thứ tư.}
{Mọi phát minh kĩ thuật của Vật lí đều tuyệt đối an toàn và không bao giờ gây ra nguy cơ hủy hoại hệ sinh thái tự nhiên.}
\loigiai{
\begin{itemchoice}
\item (Đúng) Cuộc cách mạng công nghiệp lần thứ hai gắn liền với việc phát hiện ra hiện tượng cảm ứng điện từ của Faraday. (Vì): Mệnh đề một đúng vì phát hiện hiện tượng cảm ứng điện từ dẫn đến sự ra đời của máy phát điện và động cơ điện, mở đầu kỉ nguyên sử dụng điện năng
\item (Sai) Tự động hóa sản xuất trong cuộc cách mạng công nghiệp lần thứ ba dựa trên các nghiên cứu về động cơ hơi nước. (Vì): Mệnh đề hai sai vì tự động hóa trong cách mạng công nghiệp lần thứ ba dựa trên các nghiên cứu về điện tử, chất bán dẫn và vi mạch chứ không phải động cơ hơi nước
\item (Đúng) Trí tuệ nhân tạo, robot và internet vạn vật là những công nghệ đặc trưng của cuộc cách mạng công nghiệp lần thứ tư. (Vì): Mệnh đề ba đúng vì cách mạng công nghiệp lần thứ tư dựa trên các công nghệ cao như robot, trí tuệ nhân tạo, internet vạn vật
\item (Sai) Mọi phát minh kĩ thuật của Vật lí đều tuyệt đối an toàn và không bao giờ gây ra nguy cơ hủy hoại hệ sinh thái tự nhiên. (Vì): Mệnh đề bốn sai vì các phát minh vật lí nếu ứng dụng sai mục đích có thể gây ô nhiễm môi trường, rò rỉ phóng xạ hoặc thảm họa hạt nhân
\end{itemchoice}
}
\end{ex}

% ===== Câu 52 =====
% ID: L10C1B1-11-TL
% Mức: VD
\dangbt{Phương pháp thực nghiệm trong nghiên cứu Vật lí}
\begin{bt}
% ID: L10C1B1-11-TL
% Mức: VD
Phân tích phát biểu “Bước đầu cần xác định vấn đề nghiên cứu.” và giải thích vì sao phát biểu đó đúng.
\loigiai{
Phát biểu đúng. Bước đầu cần xác định vấn đề nghiên cứu. Nội dung này phù hợp với kiến thức trọng tâm của dạng bài.
}
\end{bt}

% ===== Câu 53 =====
% ID: L10C1B1-11-TLN
% Mức: VD
\begin{ex}
% ID: L10C1B1-11-TLN
% Mức: VD
Cho bốn nhận định: (1) Phương án thí nghiệm phải xác định đại lượng cần đo. (2) Không cần kiểm soát điều kiện. (3) Các yếu tố khác nên được kiểm soát. (4) Thực nghiệm chỉ cần kết luận, không cần dữ liệu. Có bao nhiêu nhận định đúng? Chỉ ghi số nguyên.
\shortans{2}
\loigiai{
Nhận định (1), (3) đúng; (2), (4) sai. Vậy có 2 nhận định đúng.
}
\end{ex}

% ===== Câu 54 =====
% ID: L10C1B1-11-TN
% Mức: TH
\dangbt{Nhận diện và sắp xếp các bước của phương pháp thực nghiệm}
\begin{ex}
% ID: L10C1B1-11-TN
% Mức: TH
Các hiện tượng vật lí nào sau đây không liên quan đến phương pháp thực nghiệm ?
\choice
{Thả rơi một vật từ trên cao xuống mặt đất}
{\True Tính toán quỹ đạo chuyển động của Mặt trăng dựa vào toán học}
{Ném một quả bóng lên trên cao}
{Kiểm tra sự thay đổi nhiệt độ trong quá trình nóng chảy hoặc bay hơi của một chất}
\end{ex}

% ===== Câu 55 =====
% ID: L10C1B1-110-TN
% Mức: TH
\dangbt{Vận dụng tiến trình nghiên cứu Vật lí vào tình huống thực tiễn}
\begin{ex}
% ID: L10C1B1-110-TN
% Mức: TH
Trong dạng “Vận dụng tiến trình nghiên cứu Vật lí vào tình huống thực tiễn”, phát biểu nào sau đây đúng?
\choice
{Lặp đo luôn vô ích.}
{\True Kết luận phải dựa trên dữ liệu.}
{Kế hoạch không cần dụng cụ.}
{Kết luận có thể trái dữ liệu mà không cần giải thích.}
\loigiai{
Phát biểu đúng là: Kết luận phải dựa trên dữ liệu. Vận dụng đầy đủ chu trình nghiên cứu vào vấn đề gần gũi, từ câu hỏi đến đánh giá và cải tiến.
}
\end{ex}

% ===== Câu 56 =====
% ID: L10C1B1-111-TN
% Mức: TH
\begin{ex}
% ID: L10C1B1-111-TN
% Mức: TH
Trong dạng “Vận dụng tiến trình nghiên cứu Vật lí vào tình huống thực tiễn”, phát biểu nào sau đây đúng?
\choice
{Không được cải tiến phương án.}
{\True Có thể đề xuất cải tiến sau khi đánh giá kết quả.}
{Không cần quan tâm an toàn.}
{Lặp đo luôn vô ích.}
\loigiai{
Phát biểu đúng là: Có thể đề xuất cải tiến sau khi đánh giá kết quả. Vận dụng đầy đủ chu trình nghiên cứu vào vấn đề gần gũi, từ câu hỏi đến đánh giá và cải tiến.
}
\end{ex}

% ===== Câu 57 =====
% ID: L10C1B1-112-TN
% Mức: VD
\begin{ex}
% ID: L10C1B1-112-TN
% Mức: VD
Trong dạng “Vận dụng tiến trình nghiên cứu Vật lí vào tình huống thực tiễn”, phát biểu nào sau đây đúng?
\choice
{Không cần quan tâm an toàn.}
{Lặp đo luôn vô ích.}
{\True Kế hoạch cần nêu dụng cụ và cách thu thập dữ liệu.}
{Kế hoạch không cần dụng cụ.}
\loigiai{
Phát biểu đúng là: Kế hoạch cần nêu dụng cụ và cách thu thập dữ liệu. Vận dụng đầy đủ chu trình nghiên cứu vào vấn đề gần gũi, từ câu hỏi đến đánh giá và cải tiến.
}
\end{ex}

% ===== Câu 58 =====
% ID: L10C1B1-113-TN
% Mức: VD
\begin{ex}
% ID: L10C1B1-113-TN
% Mức: VD
Trong dạng “Vận dụng tiến trình nghiên cứu Vật lí vào tình huống thực tiễn”, phát biểu nào sau đây đúng?
\choice
{Sai số không cần được xem xét.}
{Không được cải tiến phương án.}
{\True Cần nêu nguồn sai số và giới hạn của kết quả.}
{Không cần quan tâm an toàn.}
\loigiai{
Phát biểu đúng là: Cần nêu nguồn sai số và giới hạn của kết quả. Vận dụng đầy đủ chu trình nghiên cứu vào vấn đề gần gũi, từ câu hỏi đến đánh giá và cải tiến.
}
\end{ex}

% ===== Câu 59 =====
% ID: L10C1B1-114-TN
% Mức: VDC
\begin{ex}
% ID: L10C1B1-114-TN
% Mức: VDC
Trong dạng “Vận dụng tiến trình nghiên cứu Vật lí vào tình huống thực tiễn”, phát biểu nào sau đây đúng?
\choice
{Không cần ghi đơn vị trong bảng số liệu.}
{Sai số không cần được xem xét.}
{Không được cải tiến phương án.}
{\True Bảng số liệu giúp tổ chức kết quả đo.}
\loigiai{
Phát biểu đúng là: Bảng số liệu giúp tổ chức kết quả đo. Vận dụng đầy đủ chu trình nghiên cứu vào vấn đề gần gũi, từ câu hỏi đến đánh giá và cải tiến.
}
\end{ex}

% ===== Câu 60 =====
% ID: L10C1B1-115-TN
% Mức: VDC
\begin{ex}
% ID: L10C1B1-115-TN
% Mức: VDC
Trong dạng “Vận dụng tiến trình nghiên cứu Vật lí vào tình huống thực tiễn”, phát biểu nào sau đây đúng?
\choice
{Có thể chọn bỏ dữ liệu tùy ý.}
{Câu hỏi càng mơ hồ càng dễ nghiên cứu.}
{Không cần ghi đơn vị trong bảng số liệu.}
{\True Lặp phép đo giúp đánh giá độ ổn định.}
\loigiai{
Phát biểu đúng là: Lặp phép đo giúp đánh giá độ ổn định. Vận dụng đầy đủ chu trình nghiên cứu vào vấn đề gần gũi, từ câu hỏi đến đánh giá và cải tiến.
}
\end{ex}

% ===== Câu 61 =====
% ID: L10C1B1-116-TN
% Mức: NB
\dangbt{Xác định đối tượng nghiên cứu và mục tiêu của môn Vật lí}
\begin{ex}
% ID: L10C1B1-116-TN
% Mức: NB
Đối tượng nghiên cứu của Vật lí gồm
\choice
{Vật chất và năng lượng}
{Các chuyển động cơ học và năng lượng}
{\True các dạng vận động của vật chất và năng lượng}
{Các hiện tượng tự nhiên}
\loigiai{
Đối tượng nghiên cứu của Vật lí là các dạng vận động của vật chất và năng lượng. Đây là quan điểm bao quát trong khoa học hiện đại.
}
\end{ex}

% ===== Câu 62 =====
% ID: L10C1B1-117-TN
% Mức: NB
\begin{ex}
% ID: L10C1B1-117-TN
% Mức: NB
Lĩnh vực nghiên cứu nào sau đây của Vật lí ?
\choice
{Nghiên cứu về sự thay đổi của các chất khi kết hợp với nhau}
{\True Nghiên cứu về các dạng chuyển động và các dạng năng lượng khác nhau}
{Nghiên cứu sự phát minh và phát triển của các vi khuẩn}
{Nghiên cứu về sự hình thành và phát triển của các tầng lớp, giai cấp trong xã hội}
\end{ex}

% ===== Câu 63 =====
% ID: L10C1B1-118-TN
% Mức: NB
\begin{ex}
% ID: L10C1B1-118-TN
% Mức: NB
Các nhà triết học tìm hiểu thế giới tự nhiên dựa trên quan sát và suy luận chủ quan trong khoảng thời gian nào ?
\choice
{Từ năm $1900$ đến nay}
{Từ cuối thế kỉ XIX đến nay}
{Từ thế kỉ XVII đến cuối thế kỉ XIX}
{\True Từ năm $350$ TCN đến thế kỉ XVI}
\end{ex}

% ===== Câu 64 =====
% ID: L10C1B1-119-TN
% Mức: NB
\begin{ex}
% ID: L10C1B1-119-TN
% Mức: NB
Đối tượng nghiên cứu của Vật lí là gì?
\choice
{\True Các dạng vận động của vật chất và năng lượng}
{Qui luật tương tác của các dạng năng lượng}
{Các dạng vận động và tương tác của vật chất}
{Nghiên cứu về nhiệt động lực học}
\end{ex}

% ===== Câu 65 =====
% ID: L10C1B1-12-DS
% Mức: NB
\dangbt{Phân tích vai trò và ảnh hưởng hai mặt của Vật lí đối với đời sống, khoa học và công nghệ}
\begin{ex}
% ID: L10C1B1-12-DS
% Mức: NB
Một số ứng dụng của vật lí trong đời sống:
\choiceTF
{\True Chụp cộng hưởng từ (MRI) trong y học}
{\True Ứng dụng cảm biến không dây trong nông nghiệp}
{Xây dựng luật giao thông đường bộ}
{\True Tạo giống cây bằng chiếu xạ đột biến}
\loigiai{
\begin{itemchoice}
\item (Đúng) Chụp cộng hưởng từ (MRI) trong y học. (Vì): MRI là thiết bị hiện đại sử dụng kiến thức vật lí
\item (Đúng) Ứng dụng cảm biến không dây trong nông nghiệp. (Vì): Cảm biến là ứng dụng điện tử - vật lí
\item (Sai) Xây dựng luật giao thông đường bộ. (Vì): Không phải ứng dụng trực tiếp của vật lí
\item (Đúng) Tạo giống cây bằng chiếu xạ đột biến. (Vì): Chiếu xạ là ứng dụng của vật lí hạt nhân
\end{itemchoice}
}
\end{ex}

% ===== Câu 66 =====
% ID: L10C1B1-12-TL
% Mức: VD
\dangbt{Phương pháp thực nghiệm trong nghiên cứu Vật lí}
\begin{bt}
% ID: L10C1B1-12-TL
% Mức: VD
Xây dựng một quy trình ngắn để giải quyết nhiệm vụ thuộc dạng “Phương pháp thực nghiệm trong nghiên cứu Vật lí”.
\loigiai{
Quy trình: đọc yêu cầu; xác định khái niệm hoặc đại lượng; chọn nguyên tắc/công thức; xử lí dữ kiện; kiểm tra đơn vị và điều kiện; kết luận. Tiến trình cơ bản: xác định vấn đề, dự đoán, thiết kế và thực hiện thí nghiệm, xử lí dữ liệu, kết luận và điều chỉnh.
}
\end{bt}

% ===== Câu 67 =====
% ID: L10C1B1-12-TLN
% Mức: VDC
\begin{ex}
% ID: L10C1B1-12-TLN
% Mức: VDC
Cho bốn nhận định: (1) Phương pháp thực nghiệm sử dụng thí nghiệm để kiểm tra dự đoán. (2) Kết quả trái dự đoán phải bị bỏ đi. (3) Nếu kết quả không phù hợp, cần xem lại giả thuyết hoặc phương án. (4) Không cần ghi chép số liệu. Có bao nhiêu nhận định đúng? Chỉ ghi số nguyên.
\shortans{2}
\loigiai{
Nhận định (1), (3) đúng; (2), (4) sai. Vậy có 2 nhận định đúng.
}
\end{ex}

% ===== Câu 68 =====
% ID: L10C1B1-12-TN
% Mức: TH
\dangbt{Nhận diện và sắp xếp các bước của phương pháp thực nghiệm}
\begin{ex}
% ID: L10C1B1-12-TN
% Mức: TH
Các hiện tượng vật lí nào sau đây liên quan đến phương pháp lí thuyết ?
\choice
{\True Ô tô khi chạy đường dài có thể xem ô tô như là một chất điểm}
{Thả rơi một vật từ trên cao xuống mặt đất}
{Kiểm tra sự thay đổi nhiệt độ trong quá trình nóng chảy hoặc bay hơi của một chất}
{Ném một quả bóng lên trên cao}
\end{ex}

% ===== Câu 69 =====
% ID: L10C1B1-120-TN
% Mức: NB
\dangbt{Xác định đối tượng nghiên cứu và mục tiêu của môn Vật lí}
\begin{ex}
% ID: L10C1B1-120-TN
% Mức: NB
Kính thiên văn Hubble giúp nghiên cứu:
\choice
{Vi khuẩn, virus}
{Thành phần hóa học của thuốc}
{\True Vũ trụ, thiên hà xa xôi}
{Cấu trúc DNA}
\loigiai{
Kính thiên văn không gian Hubble chụp được hình ảnh của các thiên hà cách Trái Đất hơn 13 tỉ năm ánh sáng, giúp nghiên cứu vũ trụ.
}
\end{ex}

% ===== Câu 70 =====
% ID: L10C1B1-121-TN
% Mức: NB
\begin{ex}
% ID: L10C1B1-121-TN
% Mức: NB
Chọn cụm từ thích hợp điền vào chỗ trống: Đối tượng nghiên cứu của Vật lí gồm các dạng … … … của vật chất và năng lượng.
\choice
{trường}
{chất}
{năng lượng}
{\True vận động}
\loigiai{
Đoạn hoàn chỉnh là: “Đối tượng nghiên cứu của Vật lí gồm các dạng \textbf{vận động} của vật chất và năng lượng.”
}
\end{ex}

% ===== Câu 71 =====
% ID: L10C1B1-122-TN
% Mức: NB
\begin{ex}
% ID: L10C1B1-122-TN
% Mức: NB
Các nhà vật lí tập trung vào các mô hình lí thuyết tìm hiểu thế giới vi mô và sử dụng thí nghiệm để kiểm chứng trong khoảng thời gian nào ?
\choice
{Từ năm $1900$ đến nay}
{Từ năm $350$ TCN đến năm $1831$}
{\True Từ cuối thế kỉ XIX đến nay}
{Từ thế kỉ XVII đến cuối thế kỉ XIX}
\end{ex}

% ===== Câu 72 =====
% ID: L10C1B1-123-TN
% Mức: NB
\begin{ex}
% ID: L10C1B1-123-TN
% Mức: NB
Câu trả lời đúng nhất. Mục tiêu của Vật lí là
\choice
{khám phá ra qui luật chi phối sự vận động của vật chất}
{khám phá ra các qui luật chuyển động}
{\True khám phá ra qui luật tổng quát nhất chi phối sự vận động của vật chất và năng lượng, cũng như tương tác giữa chúng ở cấp độ vi mô và vĩ mô}
{khám phá năng lượng của vật chất ở nhiều cấp độ}
\end{ex}

% ===== Câu 73 =====
% ID: L10C1B1-124-TN
% Mức: NB
\begin{ex}
% ID: L10C1B1-124-TN
% Mức: NB
Đối tượng nghiên cứu nào sau đây không thuộc phạm vi nghiên cứu chủ yếu của Vật lí học?
\choice
{Các dạng vận động của cơ học và nhiệt học}
{Sự chuyển hóa năng lượng giữa các dạng điện và quang}
{\True Quá trình sinh trưởng và phát triển của các loài sinh vật}
{Cấu tạo chất và liên kết của các hạt ở cấp độ vi mô}
\loigiai{
Đối tượng nghiên cứu chủ yếu của Vật lí là các dạng vận động của vật chất và năng lượng. Quá trình sinh trưởng và phát triển của các loài sinh vật thuộc đối tượng nghiên cứu của Sinh học.
}
\end{ex}

% ===== Câu 74 =====
% ID: L10C1B1-125-TN
% Mức: NB
\begin{ex}
% ID: L10C1B1-125-TN
% Mức: NB
Mục tiêu của môn Vật lí là:
\choice
{\True khám phá ra quy luật tổng quát nhất chi phối sự vận động của vật chất và năng lượng, cũng như tương tác giữa chúng ở mọi cấp độ: vi mô, vĩ mô}
{khám phá ra quy luật tổng quát nhất chi phối sự vận động của vật chất và năng lượng}
{khảo sát sự tương tác của vật chất ở mọi cấp độ: vi mô, vĩ mô}
{khám phá ra quy luật vận động cũng như tương tác của vật chất ở mọi cấp độ: vi mô, vĩ mô}
\loigiai{
Mục tiêu của Vật lí không chỉ là phát hiện quy luật vận động mà còn phải xem xét sự tương tác giữa vật chất và năng lượng ở mọi cấp độ.
}
\end{ex}

% ===== Câu 75 =====
% ID: L10C1B1-126-TN
% Mức: NB
\begin{ex}
% ID: L10C1B1-126-TN
% Mức: NB
Các phương pháp nghiên cứu nào sau đây thường dùng trong lĩnh vực Vật lí ?
\choice
{Phương pháp thực nghiệm và phương pháp mô hình}
{Phương pháp thực nghiệm và phương pháp quan sát - suy luận}
{Phương pháp mô hình và phương pháp quan sát - suy luận}
{\True Phương pháp thực nghiệm, phương pháp mô hình và phương pháp quan sát - suy luận}
\end{ex}

% ===== Câu 76 =====
% ID: L10C1B1-127-TN
% Mức: NB
\begin{ex}
% ID: L10C1B1-127-TN
% Mức: NB
Vật lý nghiên cứu các hệ thống ở các quy mô:
\choice
{Chỉ ở quy mô vĩ mô (các vật thể lớn)}
{Chỉ ở quy mô vi mô (hạt nhân nguyên tử, hạt cơ bản)}
{\True Cả ở quy mô vi mô và vĩ mô}
{Chỉ ở quy mô trung gian}
\loigiai{
Vật lý là một ngành khoa học tự nhiên nghiên cứu các hiện tượng và quy luật của tự nhiên. Nó bao gồm việc nghiên cứu các hệ thống ở nhiều quy mô khác nhau, từ vi mô đến vĩ mô. Ở quy mô vi mô, vật lý nghiên cứu các hạt cơ bản như electron, proton, neutron, và các hạt trong hạt nhân nguyên tử. Ở quy mô vĩ mô, vật lý nghiên cứu các vật thể lớn như hành tinh, sao, và thiên hà. Do đó, vật lý không chỉ giới hạn ở một quy mô nhất định mà bao gồm cả hai quy mô vi mô và vĩ mô.
 
 Phương án $C$ là đúng vì vật lý nghiên cứu cả các hệ thống ở quy mô vi mô và vĩ mô. Điều này bao gồm việc nghiên cứu các hạt cơ bản và các hiện tượng xảy ra ở cấp độ nguyên tử, cũng như các hiện tượng xảy ra ở cấp độ vũ trụ như chuyển động của các hành tinh và sự hình thành của các thiên hà.
}
\end{ex}

% ===== Câu 77 =====
% ID: L10C1B1-128-TN
% Mức: NB
\begin{ex}
% ID: L10C1B1-128-TN
% Mức: NB
Mục tiêu chính của Vật lý là:
\choice
{Mô tả các hiện tượng tự nhiên}
{Giải thích nguyên nhân của các hiện tượng tự nhiên}
{Dự đoán các hiện tượng tự nhiên}
{\True Tất cả các đáp án trên}
\loigiai{
Mục tiêu của Vật lý là mô tả, giải thích và dự đoán các hiện tượng tự nhiên, từ đó giúp con người hiểu và ứng dụng các quy luật tự nhiên.
}
\end{ex}

% ===== Câu 78 =====
% ID: L10C1B1-129-TN
% Mức: TH
\begin{ex}
% ID: L10C1B1-129-TN
% Mức: TH
Vào thời tiền Vật lí, các nhà triết học tìm hiểu thế giới tự nhiên dựa trên quan sát và suy luận chủ quan thể hiện ở nội dung nào sau đây ?
\choice
{Cái lông chim và hòn bi rơi nhanh như nhau trong ống hút hết không khí}
{Hiện tượng ánh sáng làm bật các electron ra khỏi kim loại}
{Hiện tượng cầu vồng}
{\True Vật nặng bao giờ cũng rơi nhanh hơn vật nhẹ}
\end{ex}

% ===== Câu 79 =====
% ID: L10C1B1-13-DS
% Mức: NB
\dangbt{Phân tích vai trò và ảnh hưởng hai mặt của Vật lí đối với đời sống, khoa học và công nghệ}
\begin{ex}
% ID: L10C1B1-13-DS
% Mức: NB
Ảnh hưởng của vật lí đến cuộc sống:
\choiceTF
{Vật lí không gây ảnh hưởng tiêu cực}
{\True Vật lí là cơ sở của khoa học và công nghệ}
{\True Vật lí kết hợp các ngành để tối ưu kết quả}
{Vật lí chỉ áp dụng cho năng lượng}
\loigiai{
\begin{itemchoice}
\item (Sai) Vật lí không gây ảnh hưởng tiêu cực. (Vì): Vật lí có thể gây tác hại nếu lạm dụng
\item (Đúng) Vật lí là cơ sở của khoa học và công nghệ. (Vì): Là nền tảng của các ngành khác
\item (Đúng) Vật lí kết hợp các ngành để tối ưu kết quả. (Vì): Liên ngành giúp nâng cao hiệu quả
\item (Sai) Vật lí chỉ áp dụng cho năng lượng. (Vì): Ứng dụng rất rộng
\end{itemchoice}
}
\end{ex}

% ===== Câu 80 =====
% ID: L10C1B1-13-TL
% Mức: VDC
\dangbt{Phương pháp thực nghiệm trong nghiên cứu Vật lí}
\begin{bt}
% ID: L10C1B1-13-TL
% Mức: VDC
Phân tích phát biểu “Kết quả trái dự đoán phải bị bỏ đi.”, chỉ ra sai sót và sửa lại.
\loigiai{
Phát biểu sai: Kết quả trái dự đoán phải bị bỏ đi. Cần sửa thành: Sau khi nêu vấn đề có thể đề xuất giả thuyết hoặc dự đoán. Tiến trình cơ bản: xác định vấn đề, dự đoán, thiết kế và thực hiện thí nghiệm, xử lí dữ liệu, kết luận và điều chỉnh.
}
\end{bt}

% ===== Câu 81 =====
% ID: L10C1B1-13-TLN
% Mức: VDC
\begin{ex}
% ID: L10C1B1-13-TLN
% Mức: VDC
Cho bốn nhận định: (1) Kết quả đo cần được xử lí trước khi kết luận. (2) Một phép đo duy nhất luôn tuyệt đối chính xác. (3) Thí nghiệm cần có khả năng lặp lại. (4) Có thể thay đổi đồng thời mọi yếu tố mà vẫn xác định chắc chắn nguyên nhân. Có bao nhiêu nhận định đúng? Chỉ ghi số nguyên.
\shortans{2}
\loigiai{
Nhận định (1), (3) đúng; (2), (4) sai. Vậy có 2 nhận định đúng.
}
\end{ex}

% ===== Câu 82 =====
% ID: L10C1B1-13-TN
% Mức: TH
\dangbt{Nhận diện và sắp xếp các bước của phương pháp thực nghiệm}
\begin{ex}
% ID: L10C1B1-13-TN
% Mức: TH
Các hiện tượng nào sau đây có liên quan đến phương pháp thực nghiệm:
\choice
{Ô tô khi chạy đường dài có thể xem ô tô như là một chất điểm}
{Quả địa cầu là mô hình thu nhỏ của Trái Đất}
{\True Kiểm tra sự thay đổi nhiệt độ trong quá trình nóng chảy hoặc bay hơi của một chất}
{Tính toán quỹ đạo chuyển động của Thiên vương tinh dựa vào toán học}
\loigiai{
Thực nghiệm là việc tiến hành thí nghiệm và đo đạc để kiểm chứng hiện tượng hoặc giả thuyết.
}
\end{ex}

% ===== Câu 83 =====
% ID: L10C1B1-130-TN
% Mức: TH
\dangbt{Xác định đối tượng nghiên cứu và mục tiêu của môn Vật lí}
\begin{ex}
% ID: L10C1B1-130-TN
% Mức: TH
Hoạt động nào sau đây là hoạt động nghiên cứu khoa học?
\choice
{Trồng hoa trong nhà kính}
{\True Tìm vaccine phòng chống virus trong phòng thí nghiệm}
{Sản xuất muối ăn từ nước biển}
{Vận hành nhà máy thủy điện để sản xuất điện}
\loigiai{
Tìm vaccine là hoạt động nghiên cứu bản chất của virus và tìm ra phương án ngăn chặn – đó là khoa học.
}
\end{ex}

% ===== Câu 84 =====
% ID: L10C1B1-131-TN
% Mức: TH
\begin{ex}
% ID: L10C1B1-131-TN
% Mức: TH
Đối tượng nghiên cứu nào sau đây thuộc lĩnh vực Vật lí ?
\choice
{\True Dòng điện không đổi}
{Sự sinh trưởng và phát triển của các loài trong thế giới tự nhiên}
{Sự cấu tạo và biến đổi các chất}
{Hiện tượng quang hợp}
\end{ex}

% ===== Câu 85 =====
% ID: L10C1B1-132-TN
% Mức: TH
\begin{ex}
% ID: L10C1B1-132-TN
% Mức: TH
Đối tượng nghiên cứu nào sau đây không thuộc lĩnh vực Vật lí?
\choice
{Vật chất và sự vận động, năng lượng}
{Vũ trụ (các hành tinh, ngôi sao...)}
{\True Các chất và sự biến đổi các chất, phương trình phản ứng của các chất trong tự nhiên}
{Trái Đất}
\end{ex}

% ===== Câu 86 =====
% ID: L10C1B1-133-TN
% Mức: TH
\begin{ex}
% ID: L10C1B1-133-TN
% Mức: TH
Một nhà vật lí đang nghiên cứu về sự thay đổi màu sắc của một chất lỏng khi nhiệt độ tăng lên. Đây là một ví dụ về việc nghiên cứu Vật lí ở cấp độ nào?
\choice
{\True Vi mô}
{Vĩ mô}
{Nguyên tử và phân tử}
{Lượng tử}
\loigiai{
Sự thay đổi màu sắc của chất lỏng là một hiện tượng có thể quan sát trực tiếp bằng mắt thường, không cần đến các thiết bị chuyên dụng để quan sát các hạt cấu tạo nên chất lỏng. Do đó, đây là nghiên cứu ở cấp độ vĩ mô. Nghiên cứu cấp độ vi mô, nguyên tử, hạ nguyên tử thường liên quan đến cấu trúc bên trong của vật chất.
}
\end{ex}

% ===== Câu 87 =====
% ID: L10C1B1-134-TN
% Mức: NB
\dangbt{Đối tượng nghiên cứu và mục tiêu của môn Vật lí}
\begin{ex}
% ID: L10C1B1-134-TN
% Mức: NB
Trong dạng “Đối tượng nghiên cứu và mục tiêu của môn Vật lí”, phát biểu nào sau đây đúng?
\choice
{\True Vật lí nghiên cứu các dạng vận động của vật chất và năng lượng.}
{Vật lí chỉ nghiên cứu vật thể nhìn thấy bằng mắt.}
{Vật lí chỉ nghiên cứu cấp độ vĩ mô.}
{Phản ứng tạo chất mới là đối tượng trung tâm duy nhất của Vật lí.}
\loigiai{
Phát biểu đúng là: Vật lí nghiên cứu các dạng vận động của vật chất và năng lượng. Xác định đúng đối tượng, phạm vi và mục tiêu của Vật lí; phân biệt với đối tượng trung tâm của các môn khoa học khác.
}
\end{ex}

% ===== Câu 88 =====
% ID: L10C1B1-135-TN
% Mức: NB
\begin{ex}
% ID: L10C1B1-135-TN
% Mức: NB
Trong dạng “Đối tượng nghiên cứu và mục tiêu của môn Vật lí”, phát biểu nào sau đây đúng?
\choice
{\True Chuyển động của hành tinh thuộc cấp độ vĩ mô.}
{Sự chuyển động của electron không thuộc Vật lí.}
{Vật lí không thể dự đoán hiện tượng.}
{Vật lí chỉ nghiên cứu vật thể nhìn thấy bằng mắt.}
\loigiai{
Phát biểu đúng là: Chuyển động của hành tinh thuộc cấp độ vĩ mô. Xác định đúng đối tượng, phạm vi và mục tiêu của Vật lí; phân biệt với đối tượng trung tâm của các môn khoa học khác.
}
\end{ex}

% ===== Câu 89 =====
% ID: L10C1B1-136-TN
% Mức: NB
\begin{ex}
% ID: L10C1B1-136-TN
% Mức: NB
Trong dạng “Đối tượng nghiên cứu và mục tiêu của môn Vật lí”, phát biểu nào sau đây đúng?
\choice
{\True Nghiên cứu sự truyền nhiệt thuộc lĩnh vực Vật lí.}
{Mọi vấn đề sinh học đều chỉ thuộc Vật lí.}
{Vật lí không nghiên cứu năng lượng.}
{Sự chuyển động của electron không thuộc Vật lí.}
\loigiai{
Phát biểu đúng là: Nghiên cứu sự truyền nhiệt thuộc lĩnh vực Vật lí. Xác định đúng đối tượng, phạm vi và mục tiêu của Vật lí; phân biệt với đối tượng trung tâm của các môn khoa học khác.
}
\end{ex}

% ===== Câu 90 =====
% ID: L10C1B1-137-TN
% Mức: TH
\begin{ex}
% ID: L10C1B1-137-TN
% Mức: TH
Trong dạng “Đối tượng nghiên cứu và mục tiêu của môn Vật lí”, phát biểu nào sau đây đúng?
\choice
{Vật lí không nghiên cứu năng lượng.}
{\True Mục tiêu của Vật lí là tìm các quy luật tổng quát chi phối vật chất và năng lượng.}
{Sự chuyển động của electron không thuộc Vật lí.}
{Vật lí không thể dự đoán hiện tượng.}
\loigiai{
Phát biểu đúng là: Mục tiêu của Vật lí là tìm các quy luật tổng quát chi phối vật chất và năng lượng. Xác định đúng đối tượng, phạm vi và mục tiêu của Vật lí; phân biệt với đối tượng trung tâm của các môn khoa học khác.
}
\end{ex}

% ===== Câu 91 =====
% ID: L10C1B1-138-TN
% Mức: TH
\begin{ex}
% ID: L10C1B1-138-TN
% Mức: TH
Trong dạng “Đối tượng nghiên cứu và mục tiêu của môn Vật lí”, phát biểu nào sau đây đúng?
\choice
{Vật lí hoàn toàn tách rời công nghệ.}
{\True Vật lí là cơ sở của nhiều ngành khoa học tự nhiên.}
{Mọi vấn đề sinh học đều chỉ thuộc Vật lí.}
{Vật lí không nghiên cứu năng lượng.}
\loigiai{
Phát biểu đúng là: Vật lí là cơ sở của nhiều ngành khoa học tự nhiên. Xác định đúng đối tượng, phạm vi và mục tiêu của Vật lí; phân biệt với đối tượng trung tâm của các môn khoa học khác.
}
\end{ex}

% ===== Câu 92 =====
% ID: L10C1B1-139-TN
% Mức: TH
\begin{ex}
% ID: L10C1B1-139-TN
% Mức: TH
Trong dạng “Đối tượng nghiên cứu và mục tiêu của môn Vật lí”, phát biểu nào sau đây đúng?
\choice
{Vật lí không có ứng dụng trong đời sống.}
{\True Kiến thức Vật lí có thể dùng để giải quyết vấn đề thực tiễn.}
{Mục tiêu của Vật lí chỉ là ghi chép hiện tượng.}
{Vật lí hoàn toàn tách rời công nghệ.}
\loigiai{
Phát biểu đúng là: Kiến thức Vật lí có thể dùng để giải quyết vấn đề thực tiễn. Xác định đúng đối tượng, phạm vi và mục tiêu của Vật lí; phân biệt với đối tượng trung tâm của các môn khoa học khác.
}
\end{ex}

% ===== Câu 93 =====
% ID: L10C1B1-14-DS
% Mức: TH
\dangbt{Phân tích vai trò và ảnh hưởng hai mặt của Vật lí đối với đời sống, khoa học và công nghệ}
\begin{ex}
% ID: L10C1B1-14-DS
% Mức: TH
Nhận xét các nhận định sau về vai trò của Vật lí đối với khoa học, kĩ thuật và công nghệ
\choiceTF
{Vật lí là ngành khoa học độc lập và không có mối liên hệ với các ngành khoa học khác.}
{Thành tựu nghiên cứu máy hơi nước do Jams Watt (Giêm Oát) sáng chế năm $1765$ dựa trên những kết quả nghiên cứu về Cơ học của Vật lí.}
{Sinh học lượng tử và Hóa học lượng tử là ngành nghiên cứu không liên quan đến Vật lí.}
{\True Ứng dụng các thành tựu Vật lí có thể hủy loại hệ sinh thái, gây ô nhiễm môi trường sống nếu không được sử dụng đúng phương pháp.}
\loigiai{
\begin{itemchoice}
\item (Sai) Vật lí là ngành khoa học độc lập và không có mối liên hệ với các ngành khoa học khác.
\item (Sai) Thành tựu nghiên cứu máy hơi nước do Jams Watt (Giêm Oát) sáng chế năm $1765$ dựa trên những kết quả nghiên cứu về Cơ học của Vật lí.
\item (Sai) Sinh học lượng tử và Hóa học lượng tử là ngành nghiên cứu không liên quan đến Vật lí.
\item (Đúng) Ứng dụng các thành tựu Vật lí có thể hủy loại hệ sinh thái, gây ô nhiễm môi trường sống nếu không được sử dụng đúng phương pháp.
\end{itemchoice}
}
\end{ex}

% ===== Câu 94 =====
% ID: L10C1B1-14-TL
% Mức: NB
\dangbt{Quá trình phát triển của Vật lí}
\begin{bt}
% ID: L10C1B1-14-TL
% Mức: NB
Đề xuất cách vận dụng nội dung “Quá trình phát triển của Vật lí” trong một tình huống học tập hoặc thực tiễn.
\loigiai{
Cần xác định tình huống, chọn kiến thức phù hợp, thực hiện theo quy trình, kiểm tra điều kiện và đánh giá kết quả. Nhận biết các giai đoạn và đặc trưng phát triển của Vật lí, nhấn mạnh quan hệ giữa lí thuyết, thực nghiệm và thiết bị.
}
\end{bt}

% ===== Câu 95 =====
% ID: L10C1B1-14-TLN
% Mức: NB
\begin{ex}
% ID: L10C1B1-14-TLN
% Mức: NB
Cho bốn nhận định: (1) Galileo góp phần đặt nền móng cho phương pháp thực nghiệm. (2) Lí thuyết khoa học không cần bằng chứng. (3) Quan sát và đo lường thúc đẩy sự phát triển của Vật lí. (4) Khoa học không bao giờ điều chỉnh lí thuyết. Có bao nhiêu nhận định đúng? Chỉ ghi số nguyên.
\shortans{2}
\loigiai{
Nhận định (1), (3) đúng; (2), (4) sai. Vậy có 2 nhận định đúng.
}
\end{ex}

% ===== Câu 96 =====
% ID: L10C1B1-14-TN
% Mức: TH
\dangbt{Nhận diện và sắp xếp các bước của phương pháp thực nghiệm}
\begin{ex}
% ID: L10C1B1-14-TN
% Mức: TH
Cho các dữ kiện sau. 1. Kiểm tra giả thuyết 2. Hình thành giả thuyết 3. Rút ra kết luận 4. Đề xuất vấn đề 5. Quan sát hiện tượng, suy luận Sắp xếp lại đúng các bước tìm hiểu thế giới tự nhiên dưới góc độ vật lí.
\choice
{\True 5 - 4 - 2 - 1 - 3}
{1 - 2 - 3 - 4 - 5}
{2 - 1 - 5 - 4 - 3}
{5 - 2 - 1 - 4 - 3}
\end{ex}

% ===== Câu 97 =====
% ID: L10C1B1-140-TN
% Mức: VD
\dangbt{Đối tượng nghiên cứu và mục tiêu của môn Vật lí}
\begin{ex}
% ID: L10C1B1-140-TN
% Mức: VD
Trong dạng “Đối tượng nghiên cứu và mục tiêu của môn Vật lí”, phát biểu nào sau đây đúng?
\choice
{Mục tiêu của Vật lí chỉ là ghi chép hiện tượng.}
{Vật lí hoàn toàn tách rời công nghệ.}
{\True Vật lí nghiên cứu cả cấp độ vi mô và vĩ mô.}
{Mọi vấn đề sinh học đều chỉ thuộc Vật lí.}
\loigiai{
Phát biểu đúng là: Vật lí nghiên cứu cả cấp độ vi mô và vĩ mô. Xác định đúng đối tượng, phạm vi và mục tiêu của Vật lí; phân biệt với đối tượng trung tâm của các môn khoa học khác.
}
\end{ex}

% ===== Câu 98 =====
% ID: L10C1B1-141-TN
% Mức: VD
\begin{ex}
% ID: L10C1B1-141-TN
% Mức: VD
Trong dạng “Đối tượng nghiên cứu và mục tiêu của môn Vật lí”, phát biểu nào sau đây đúng?
\choice
{Phản ứng tạo chất mới là đối tượng trung tâm duy nhất của Vật lí.}
{Vật lí không có ứng dụng trong đời sống.}
{\True Nghiên cứu dòng điện thuộc lĩnh vực Vật lí.}
{Mục tiêu của Vật lí chỉ là ghi chép hiện tượng.}
\loigiai{
Phát biểu đúng là: Nghiên cứu dòng điện thuộc lĩnh vực Vật lí. Xác định đúng đối tượng, phạm vi và mục tiêu của Vật lí; phân biệt với đối tượng trung tâm của các môn khoa học khác.
}
\end{ex}

% ===== Câu 99 =====
% ID: L10C1B1-142-TN
% Mức: VDC
\begin{ex}
% ID: L10C1B1-142-TN
% Mức: VDC
Trong dạng “Đối tượng nghiên cứu và mục tiêu của môn Vật lí”, phát biểu nào sau đây đúng?
\choice
{Vật lí chỉ nghiên cứu cấp độ vĩ mô.}
{Phản ứng tạo chất mới là đối tượng trung tâm duy nhất của Vật lí.}
{Vật lí không có ứng dụng trong đời sống.}
{\True Electron thuộc đối tượng ở cấp độ vi mô.}
\loigiai{
Phát biểu đúng là: Electron thuộc đối tượng ở cấp độ vi mô. Xác định đúng đối tượng, phạm vi và mục tiêu của Vật lí; phân biệt với đối tượng trung tâm của các môn khoa học khác.
}
\end{ex}

% ===== Câu 100 =====
% ID: L10C1B1-143-TN
% Mức: VDC
\begin{ex}
% ID: L10C1B1-143-TN
% Mức: VDC
Trong dạng “Đối tượng nghiên cứu và mục tiêu của môn Vật lí”, phát biểu nào sau đây đúng?
\choice
{Vật lí không thể dự đoán hiện tượng.}
{Vật lí chỉ nghiên cứu vật thể nhìn thấy bằng mắt.}
{Vật lí chỉ nghiên cứu cấp độ vĩ mô.}
{\True Nghiên cứu ánh sáng thuộc lĩnh vực Vật lí.}
\loigiai{
Phát biểu đúng là: Nghiên cứu ánh sáng thuộc lĩnh vực Vật lí. Xác định đúng đối tượng, phạm vi và mục tiêu của Vật lí; phân biệt với đối tượng trung tâm của các môn khoa học khác.
}
\end{ex}

% ===== Câu 101 =====
% ID: L10C1B1-15-DS
% Mức: TH
\dangbt{Phân tích vai trò và ảnh hưởng hai mặt của Vật lí đối với đời sống, khoa học và công nghệ}
\begin{ex}
% ID: L10C1B1-15-DS
% Mức: TH
Về vai trò và ảnh hưởng của vật lí học đối với các lĩnh vực khác và đời sống, nhận định nào sau đây là đúng?
\choiceTF
{Vật lí học là một môn khoa học độc lập hoàn toàn, không có mối liên hệ hay ảnh hưởng nào đến các ngành khoa học tự nhiên khác như hóa học hay sinh học.}
{Vật lí học chủ yếu tập trung nghiên cứu các vật thể có kích thước lớn và tốc độ thấp, không áp dụng cho thế giới vi mô hay các hiện tượng tốc độ cao.}
{\True Các phương pháp và công cụ nghiên cứu của vật lí thường được ứng dụng rộng rãi trong các ngành khoa học khác để giải quyết các vấn đề phức tạp.}
{\True Những khám phá trong vật lí đã và đang tạo ra những thay đổi sâu rộng trong công nghệ, y học và nhiều khía cạnh của đời sống con người.}
\loigiai{
\begin{itemchoice}
\item (Sai) Vật lí học là một môn khoa học độc lập hoàn toàn, không có mối liên hệ hay ảnh hưởng nào đến các ngành khoa học tự nhiên khác như hóa học hay sinh học. (Vì): Vật lí học là nền tảng cho nhiều ngành khoa học tự nhiên khác. Ví dụ, hóa học dựa trên các nguyên lí vật lí về cấu trúc nguyên tử, còn sinh học sử dụng các nguyên lí vật lí để giải thích các quá trình sống
\item (Sai) Vật lí học chủ yếu tập trung nghiên cứu các vật thể có kích thước lớn và tốc độ thấp, không áp dụng cho thế giới vi mô hay các hiện tượng tốc độ cao. (Vì): Vật lí học nghiên cứu các hiện tượng ở mọi quy mô, từ hạt hạ nguyên tử (vật lí hạt nhân, vật lí lượng tử) đến các thiên hà và vũ trụ (vật lí thiên văn), và ở mọi dải tốc độ, bao gồm cả tốc độ gần ánh sáng (thuyết tương đối
\item (Đúng) Các phương pháp và công cụ nghiên cứu của vật lí thường được ứng dụng rộng rãi trong các ngành khoa học khác để giải quyết các vấn đề phức tạp. (Vì): Vật lí cung cấp nhiều phương pháp và công cụ quan trọng như kính hiển vi, kính thiên văn, máy chụp X-quang, cộng hưởng từ (MRI) được sử dụng rộng rãi trong sinh học, y học, hóa học và khoa học vật liệu
\item (Đúng) Những khám phá trong vật lí đã và đang tạo ra những thay đổi sâu rộng trong công nghệ, y học và nhiều khía cạnh của đời sống con người. (Vì): Từ việc phát minh ra điện, radio, máy tính đến công nghệ laser, năng lượng hạt nhân và hình ảnh y tế, vật lí đã đóng góp to lớn vào sự phát triển của công nghệ và cải thiện chất lượng cuộc sống con người
\end{itemchoice}
}
\end{ex}

% ===== Câu 102 =====
% ID: L10C1B1-15-TL
% Mức: TH
\dangbt{Quá trình phát triển của Vật lí}
\begin{bt}
% ID: L10C1B1-15-TL
% Mức: TH
Trình bày kiến thức cốt lõi của dạng “Quá trình phát triển của Vật lí” và minh họa bằng một ví dụ.
\loigiai{
Nhận biết các giai đoạn và đặc trưng phát triển của Vật lí, nhấn mạnh quan hệ giữa lí thuyết, thực nghiệm và thiết bị. Ví dụ phù hợp cần nêu rõ đối tượng, dữ kiện và kết luận theo kiến thức SGK.
}
\end{bt}

% ===== Câu 103 =====
% ID: L10C1B1-15-TLN
% Mức: NB
\begin{ex}
% ID: L10C1B1-15-TLN
% Mức: NB
Cho bốn nhận định: (1) Kết quả mới có thể làm lí thuyết được bổ sung hoặc điều chỉnh. (2) Thiết bị đo cản trở mọi phát hiện. (3) Thực nghiệm giúp kiểm tra các dự đoán khoa học. (4) Vật lí hiện đại không nghiên cứu vi mô. Có bao nhiêu nhận định đúng? Chỉ ghi số nguyên.
\shortans{2}
\loigiai{
Nhận định (1), (3) đúng; (2), (4) sai. Vậy có 2 nhận định đúng.
}
\end{ex}

% ===== Câu 104 =====
% ID: L10C1B1-15-TN
% Mức: TH
\dangbt{Nhận diện và sắp xếp các bước của phương pháp thực nghiệm}
\begin{ex}
% ID: L10C1B1-15-TN
% Mức: TH
Các hiện tượng vật lí nào sau đây không liên quan đến phương pháp lí thuyết ?
\choice
{\True Ném một quả bóng lên trên cao}
{Biểu diễn đường truyền ánh sáng qua thấu kính}
{Tính toán quỹ đạo chuyển động của Sao Hỏa dựa vào toán học}
{Quả địa cầu là mô hình thu nhỏ của Trái đất}
\end{ex}

% ===== Câu 105 =====
% ID: L10C1B1-16-DS
% Mức: TH
\dangbt{Phân tích vai trò và ảnh hưởng hai mặt của Vật lí đối với đời sống, khoa học và công nghệ}
\begin{ex}
% ID: L10C1B1-16-DS
% Mức: TH
Về bản chất của các lí thuyết vật lí và mối quan hệ giữa vật lí với công nghệ, nhận định nào sau đây là đúng?
\choiceTF
{\True Các lí thuyết vật lí là những mô hình được xây dựng để mô tả và giải thích các hiện tượng tự nhiên, chúng luôn được kiểm chứng và điều chỉnh dựa trên thực nghiệm.}
{\True Vật lí học cung cấp nền tảng khoa học cơ bản cho sự phát triển của công nghệ, đồng thời công nghệ cũng tạo ra các công cụ và phương tiện giúp vật lí khám phá sâu hơn về thế giới.}
{Các lí thuyết vật lí, một khi đã được chấp nhận, sẽ không bao giờ thay đổi hoặc được thay thế bởi các lí thuyết mới, vì chúng đã đạt đến sự hoàn hảo tuyệt đối.}
{Công nghệ chỉ đơn thuần là ứng dụng các nguyên lí vật lí đã có sẵn và không có bất kỳ vai trò nào trong việc thúc đẩy sự tiến bộ của vật lí học.}
\loigiai{
\begin{itemchoice}
\item (Đúng) Các lí thuyết vật lí là những mô hình được xây dựng để mô tả và giải thích các hiện tượng tự nhiên, chúng luôn được kiểm chứng và điều chỉnh dựa trên thực nghiệm. (Vì): Khoa học nói chung và vật lí nói riêng có tính chất tự sửa chữa. Các lí thuyết vật lí luôn mở để được kiểm chứng, điều chỉnh hoặc thay thế khi có bằng chứng mới hoặc lí thuyết toàn diện hơn
\item (Đúng) Vật lí học cung cấp nền tảng khoa học cơ bản cho sự phát triển của công nghệ, đồng thời công nghệ cũng tạo ra các công cụ và phương tiện giúp vật lí khám phá sâu hơn về thế giới. (Vì): Công nghệ là lĩnh vực ứng dụng kiến thức khoa học, trong đó vật lí đóng vai trò nền tảng để phát triển các công nghệ mới như điện tử, quang học, năng lượng, v.v
\item (Sai) Các lí thuyết vật lí, một khi đã được chấp nhận, sẽ không bao giờ thay đổi hoặc được thay thế bởi các lí thuyết mới, vì chúng đã đạt đến sự hoàn hảo tuyệt đối. (Vì): Khoa học là quá trình không ngừng phát triển. Các lí thuyết vật lí, dù đã được kiểm chứng, vẫn có thể được cải tiến hoặc thay thế bởi các lí thuyết mới có khả năng giải thích rộng hơn hoặc chính xác hơn
\item (Sai) Công nghệ chỉ đơn thuần là ứng dụng các nguyên lí vật lí đã có sẵn và không có bất kỳ vai trò nào trong việc thúc đẩy sự tiến bộ của vật lí học. (Vì): Vật lí học nghiên cứu mọi hiện tượng tự nhiên, từ các hạt cơ bản đến cấu trúc vũ trụ, không chỉ giới hạn trong phòng thí nghiệm mà còn thông qua quan sát thiên văn và các hiện tượng tự nhiên khác
\end{itemchoice}
}
\end{ex}

% ===== Câu 106 =====
% ID: L10C1B1-16-TL
% Mức: TH
\dangbt{Quá trình phát triển của Vật lí}
\begin{bt}
% ID: L10C1B1-16-TL
% Mức: TH
So sánh hai nhận định: “Một lí thuyết khoa học cần được đối chiếu với bằng chứng.” và “Mọi quan niệm cổ đều luôn đúng.”.
\loigiai{
Nhận định thứ nhất đúng: Một lí thuyết khoa học cần được đối chiếu với bằng chứng. Nhận định thứ hai sai: Mọi quan niệm cổ đều luôn đúng. Sự khác nhau nằm ở điều kiện và bản chất kiến thức được áp dụng.
}
\end{bt}

% ===== Câu 107 =====
% ID: L10C1B1-16-TLN
% Mức: TH
\begin{ex}
% ID: L10C1B1-16-TLN
% Mức: TH
Cho bốn nhận định: (1) Vật lí cổ điển phát triển mạnh từ thế kỉ XVII. (2) Mọi quan niệm cổ đều luôn đúng. (3) Newton có đóng góp quan trọng cho cơ học cổ điển. (4) Quan sát không cung cấp dữ liệu. Có bao nhiêu nhận định đúng? Chỉ ghi số nguyên.
\shortans{2}
\loigiai{
Nhận định (1), (3) đúng; (2), (4) sai. Vậy có 2 nhận định đúng.
}
\end{ex}

% ===== Câu 108 =====
% ID: L10C1B1-16-TN
% Mức: TH
\dangbt{Nhận diện và sắp xếp các bước của phương pháp thực nghiệm}
\begin{ex}
% ID: L10C1B1-16-TN
% Mức: TH
Phát biểu nào sau đây đúng về mối quan hệ giữa phương pháp lí thuyết và thực nghiệm?
\choice
{Phương pháp lí thuyết quan trọng hơn thực nghiệm}
{Hai phương pháp không liên quan nhau}
{\True Hai phương pháp hỗ trợ cho nhau, trong đó thực nghiệm có vai trò quyết định}
{Thực nghiệm là một phần của phương pháp lí thuyết}
\loigiai{
Hai phương pháp hỗ trợ lẫn nhau, nhưng thực nghiệm giữ vai trò \textbf{quyết định} trong việc xác minh các giả thuyết, mô hình lí thuyết.
}
\end{ex}

% ===== Câu 109 =====
% ID: L10C1B1-17-DS
% Mức: NB
\dangbt{Phương pháp mô hình trong nghiên cứu Vật lí}
\begin{ex}
% ID: L10C1B1-17-DS
% Mức: NB
Xét các phát biểu sau về “Phương pháp mô hình trong nghiên cứu Vật lí” (bộ 3).
\choiceTF
{\True Mỗi mô hình chỉ phù hợp trong phạm vi nhất định.}
{Sơ đồ không thể là mô hình.}
{\True Sơ đồ tia sáng là một dạng mô hình.}
{Mô hình không có khả năng dự đoán.}
\loigiai{
\begin{itemchoice}
\item (Đúng) Mỗi mô hình chỉ phù hợp trong phạm vi nhất định. (Vì): Mỗi mô hình chỉ phù hợp trong phạm vi nhất định. b) S: Sơ đồ không thể là mô hình. c) Đ: Sơ đồ tia sáng là một dạng mô hình. d) S: Mô hình không có khả năng dự đoán. Nhận diện mô hình vật chất, mô hình toán học và mô hình lí thuyết; xác định phạm vi áp dụng và kiểm chứng
\item (Sai) Sơ đồ không thể là mô hình.
\item (Đúng) Sơ đồ tia sáng là một dạng mô hình.
\item (Sai) Mô hình không có khả năng dự đoán.
\end{itemchoice}
}
\end{ex}

% ===== Câu 110 =====
% ID: L10C1B1-17-TL
% Mức: VD
\dangbt{Quá trình phát triển của Vật lí}
\begin{bt}
% ID: L10C1B1-17-TL
% Mức: VD
Phân tích phát biểu “Galileo góp phần đặt nền móng cho phương pháp thực nghiệm.” và giải thích vì sao phát biểu đó đúng.
\loigiai{
Phát biểu đúng. Galileo góp phần đặt nền móng cho phương pháp thực nghiệm. Nội dung này phù hợp với kiến thức trọng tâm của dạng bài.
}
\end{bt}

% ===== Câu 111 =====
% ID: L10C1B1-17-TLN
% Mức: VD
\begin{ex}
% ID: L10C1B1-17-TLN
% Mức: VD
Cho bốn nhận định: (1) Vật lí hiện đại nghiên cứu sâu thế giới vi mô. (2) Newton không liên quan cơ học. (3) Thuyết lượng tử thuộc Vật lí hiện đại. (4) Vật lí chỉ phát triển nhờ suy luận chủ quan. Có bao nhiêu nhận định đúng? Chỉ ghi số nguyên.
\shortans{2}
\loigiai{
Nhận định (1), (3) đúng; (2), (4) sai. Vậy có 2 nhận định đúng.
}
\end{ex}

% ===== Câu 112 =====
% ID: L10C1B1-17-TN
% Mức: TH
\dangbt{Nhận diện và sắp xếp các bước của phương pháp thực nghiệm}
\begin{ex}
% ID: L10C1B1-17-TN
% Mức: TH
Các hiện tượng vật lí nào sau đây liên quan đến phương pháp thực nghiệm ?
\choice
{Quả địa cầu là mô hình thu nhỏ của Trái đất}
{Để biểu diễn đường truyền của ánh sáng người ta dùng tia sáng}
{Ô tô khi chạy đường dài có thể xem ô tô như là một chất điểm}
{\True Thả rơi một vật từ trên cao xuống mặt đất}
\end{ex}

% ===== Câu 113 =====
% ID: L10C1B1-18-DS
% Mức: TH
\dangbt{Phương pháp mô hình trong nghiên cứu Vật lí}
\begin{ex}
% ID: L10C1B1-18-DS
% Mức: TH
Xét các phát biểu sau về “Phương pháp mô hình trong nghiên cứu Vật lí” (bộ 4).
\choiceTF
{\True Có thể điều chỉnh mô hình khi bằng chứng mới xuất hiện.}
{Không được sửa mô hình khi có dữ liệu mới.}
{\True Mô hình giúp tập trung vào đặc trưng quan trọng.}
{Mô hình phải giống hoàn toàn vật thật ở mọi chi tiết.}
\loigiai{
\begin{itemchoice}
\item (Đúng) Có thể điều chỉnh mô hình khi bằng chứng mới xuất hiện. (Vì): Có thể điều chỉnh mô hình khi bằng chứng mới xuất hiện. b) S: Không được sửa mô hình khi có dữ liệu mới. c) Đ: Mô hình giúp tập trung vào đặc trưng quan trọng. d) S: Mô hình phải giống hoàn toàn vật thật ở mọi chi tiết. Nhận diện mô hình vật chất, mô hình toán học và mô hình lí thuyết; xác định phạm vi áp dụng và kiểm chứng
\item (Sai) Không được sửa mô hình khi có dữ liệu mới.
\item (Đúng) Mô hình giúp tập trung vào đặc trưng quan trọng.
\item (Sai) Mô hình phải giống hoàn toàn vật thật ở mọi chi tiết.
\end{itemchoice}
}
\end{ex}

% ===== Câu 114 =====
% ID: L10C1B1-18-TL
% Mức: VD
\dangbt{Quá trình phát triển của Vật lí}
\begin{bt}
% ID: L10C1B1-18-TL
% Mức: VD
Xây dựng một quy trình ngắn để giải quyết nhiệm vụ thuộc dạng “Quá trình phát triển của Vật lí”.
\loigiai{
Quy trình: đọc yêu cầu; xác định khái niệm hoặc đại lượng; chọn nguyên tắc/công thức; xử lí dữ kiện; kiểm tra đơn vị và điều kiện; kết luận. Nhận biết các giai đoạn và đặc trưng phát triển của Vật lí, nhấn mạnh quan hệ giữa lí thuyết, thực nghiệm và thiết bị.
}
\end{bt}

% ===== Câu 115 =====
% ID: L10C1B1-18-TLN
% Mức: VDC
\begin{ex}
% ID: L10C1B1-18-TLN
% Mức: VDC
Cho bốn nhận định: (1) Thực nghiệm giúp kiểm tra các dự đoán khoa học. (2) Vật lí hiện đại không nghiên cứu vi mô. (3) Kết quả mới có thể làm lí thuyết được bổ sung hoặc điều chỉnh. (4) Thiết bị đo cản trở mọi phát hiện. Có bao nhiêu nhận định đúng? Chỉ ghi số nguyên.
\shortans{2}
\loigiai{
Nhận định (1), (3) đúng; (2), (4) sai. Vậy có 2 nhận định đúng.
}
\end{ex}

% ===== Câu 116 =====
% ID: L10C1B1-18-TN
% Mức: TH
\dangbt{Nhận diện và sắp xếp các bước của phương pháp thực nghiệm}
\begin{ex}
% ID: L10C1B1-18-TN
% Mức: TH
Phát biểu nào sau đây là đúng về phương pháp nghiên cứu khoa học?
\choice
{Không cần giả thuyết nếu có dữ liệu thực nghiệm}
{Chỉ cần quan sát và đưa ra kết luận}
{\True Giả thuyết cần được kiểm tra trước khi rút ra kết luận}
{Có thể kiểm tra giả thuyết trước rồi mới đề xuất vấn đề}
\loigiai{
Giả thuyết cần được kiểm chứng trước khi kết luận được đưa ra.
}
\end{ex}

% ===== Câu 117 =====
% ID: L10C1B1-19-DS
% Mức: VD
\dangbt{Phương pháp mô hình trong nghiên cứu Vật lí}
\begin{ex}
% ID: L10C1B1-19-DS
% Mức: VD
Xét các phát biểu sau về “Phương pháp mô hình trong nghiên cứu Vật lí” (bộ 1).
\choiceTF
{\True Mô hình là sự biểu diễn đơn giản hóa đối tượng hoặc quá trình.}
{Mọi mô hình đều đúng trong mọi điều kiện.}
{\True Mô hình lí thuyết giúp giải thích và dự đoán.}
{Mô hình không cần kiểm chứng.}
\loigiai{
\begin{itemchoice}
\item (Đúng) Mô hình là sự biểu diễn đơn giản hóa đối tượng hoặc quá trình. (Vì): Mô hình là sự biểu diễn đơn giản hóa đối tượng hoặc quá trình. b) S: Mọi mô hình đều đúng trong mọi điều kiện. c) Đ: Mô hình lí thuyết giúp giải thích và dự đoán. d) S: Mô hình không cần kiểm chứng. Nhận diện mô hình vật chất, mô hình toán học và mô hình lí thuyết; xác định phạm vi áp dụng và kiểm chứng
\item (Sai) Mọi mô hình đều đúng trong mọi điều kiện.
\item (Đúng) Mô hình lí thuyết giúp giải thích và dự đoán.
\item (Sai) Mô hình không cần kiểm chứng.
\end{itemchoice}
}
\end{ex}

% ===== Câu 118 =====
% ID: L10C1B1-19-TL
% Mức: VDC
\dangbt{Quá trình phát triển của Vật lí}
\begin{bt}
% ID: L10C1B1-19-TL
% Mức: VDC
Phân tích phát biểu “Vật lí hiện đại không nghiên cứu vi mô.”, chỉ ra sai sót và sửa lại.
\loigiai{
Phát biểu sai: Vật lí hiện đại không nghiên cứu vi mô. Cần sửa thành: Vật lí cổ điển phát triển mạnh từ thế kỉ XVII. Nhận biết các giai đoạn và đặc trưng phát triển của Vật lí, nhấn mạnh quan hệ giữa lí thuyết, thực nghiệm và thiết bị.
}
\end{bt}

% ===== Câu 119 =====
% ID: L10C1B1-19-TLN
% Mức: VDC
\begin{ex}
% ID: L10C1B1-19-TLN
% Mức: VDC
Cho bốn nhận định: (1) Một lí thuyết khoa học cần được đối chiếu với bằng chứng. (2) Thuyết lượng tử thuộc thời tiền Vật lí. (3) Sự phát triển thiết bị tạo điều kiện cho phát hiện mới. (4) Thí nghiệm không có vai trò kiểm chứng. Có bao nhiêu nhận định đúng? Chỉ ghi số nguyên.
\shortans{2}
\loigiai{
Nhận định (1), (3) đúng; (2), (4) sai. Vậy có 2 nhận định đúng.
}
\end{ex}

% ===== Câu 120 =====
% ID: L10C1B1-19-TN
% Mức: TH
\dangbt{Nhận diện và sắp xếp các bước của phương pháp thực nghiệm}
\begin{ex}
% ID: L10C1B1-19-TN
% Mức: TH
Phương pháp nào sau đây có vai trò quyết định trong nghiên cứu vật lí?
\choice
{Phương pháp lí thuyết}
{\True Phương pháp thực nghiệm}
{Phương pháp thống kê}
{Phương pháp mô phỏng}
\loigiai{
Trong nghiên cứu vật lí, \textbf{phương pháp thực nghiệm} có vai trò quyết định vì giúp kiểm chứng giả thuyết và lí thuyết.
}
\end{ex}

% ===== Câu 121 =====
% ID: L10C1B1-20-DS
% Mức: VD
\dangbt{Phương pháp mô hình trong nghiên cứu Vật lí}
\begin{ex}
% ID: L10C1B1-20-DS
% Mức: VD
Xét các phát biểu sau về “Phương pháp mô hình trong nghiên cứu Vật lí” (bộ 5).
\choiceTF
{\True Chất điểm là mô hình dùng khi kích thước vật có thể bỏ qua.}
{Một đối tượng chỉ có duy nhất một mô hình.}
{\True Mô hình vật chất có thể là mẫu thu nhỏ.}
{Mô hình toán học không dùng phương trình.}
\loigiai{
\begin{itemchoice}
\item (Đúng) Chất điểm là mô hình dùng khi kích thước vật có thể bỏ qua. (Vì): Chất điểm là mô hình dùng khi kích thước vật có thể bỏ qua. b) S: Một đối tượng chỉ có duy nhất một mô hình. c) Đ: Mô hình vật chất có thể là mẫu thu nhỏ. d) S: Mô hình toán học không dùng phương trình. Nhận diện mô hình vật chất, mô hình toán học và mô hình lí thuyết; xác định phạm vi áp dụng và kiểm chứng
\item (Sai) Một đối tượng chỉ có duy nhất một mô hình.
\item (Đúng) Mô hình vật chất có thể là mẫu thu nhỏ.
\item (Sai) Mô hình toán học không dùng phương trình.
\end{itemchoice}
}
\end{ex}

% ===== Câu 122 =====
% ID: L10C1B1-20-TL
% Mức: NB
\dangbt{Vai trò của Vật lí đối với khoa học, kĩ thuật, công nghệ và đời sống}
\begin{bt}
% ID: L10C1B1-20-TL
% Mức: NB
Đề xuất cách vận dụng nội dung “Vai trò của Vật lí đối với khoa học, kĩ thuật, công nghệ và đời sống” trong một tình huống học tập hoặc thực tiễn.
\loigiai{
Cần xác định tình huống, chọn kiến thức phù hợp, thực hiện theo quy trình, kiểm tra điều kiện và đánh giá kết quả. Phân tích được vai trò của Vật lí qua ví dụ cụ thể và đánh giá cả lợi ích lẫn hệ quả của ứng dụng công nghệ.
}
\end{bt}

% ===== Câu 123 =====
% ID: L10C1B1-20-TLN
% Mức: NB
\begin{ex}
% ID: L10C1B1-20-TLN
% Mức: NB
Cho bốn nhận định: (1) Máy chụp X-quang là một ứng dụng của Vật lí. (2) Năng lượng gió không liên quan kiến thức Vật lí. (3) Công nghệ bán dẫn dựa trên hiểu biết Vật lí chất rắn. (4) Kính thiên văn dùng để nghiên cứu vi khuẩn. Có bao nhiêu nhận định đúng? Chỉ ghi số nguyên.
\shortans{2}
\loigiai{
Nhận định (1), (3) đúng; (2), (4) sai. Vậy có 2 nhận định đúng.
}
\end{ex}

% ===== Câu 124 =====
% ID: L10C1B1-20-TN
% Mức: TH
\dangbt{Nhận diện và sắp xếp các bước của phương pháp thực nghiệm}
\begin{ex}
% ID: L10C1B1-20-TN
% Mức: TH
Chu trình phát triển điển hình của một lí thuyết vật lí thường bao gồm các bước nào sau đây?
\choice
{Chỉ tập trung vào việc thu thập dữ liệu thí nghiệm một cách ngẫu nhiên, không cần có giả thuyết ban đầu.}
{Xây dựng một lí thuyết hoàn chỉnh dựa trên suy luận logic, sau đó tiến hành một thí nghiệm duy nhất để xác nhận.}
{Dựa hoàn toàn vào các nguyên lí triết học và suy luận trừu tượng để đưa ra các định luật, không cần thực nghiệm.}
{\True Quan sát hiện tượng $\rightarrow$ Đề xuất giả thuyết $\rightarrow$ Xây dựng mô hình/lí thuyết $\rightarrow$ Kiểm chứng bằng thực nghiệm $\rightarrow$ Điều chỉnh/Hoàn thiện lí thuyết.}
\loigiai{
Quá trình phát triển của vật lí học, cũng như khoa học nói chung, thường tuân theo một chu trình lặp đi lặp lại. Nó bắt đầu từ việc quan sát các hiện tượng tự nhiên, sau đó các nhà khoa học đề xuất giả thuyết để giải thích. Từ giả thuyết, họ xây dựng các mô hình hoặc lí thuyết toán học. Bước tiếp theo là kiểm chứng các dự đoán của lí thuyết bằng thực nghiệm. Dựa trên kết quả thực nghiệm, lí thuyết có thể được điều chỉnh, hoàn thiện hoặc thậm chí bị bác bỏ để nhường chỗ cho một lí thuyết mới tốt hơn. Đây là một quá trình liên tục và có tính tương tác cao giữa lí thuyết và thực nghiệm.
}
\end{ex}

% ===== Câu 125 =====
% ID: L10C1B1-21-DS
% Mức: VDC
\dangbt{Phương pháp mô hình trong nghiên cứu Vật lí}
\begin{ex}
% ID: L10C1B1-21-DS
% Mức: VDC
Xét các phát biểu sau về “Phương pháp mô hình trong nghiên cứu Vật lí” (bộ 2).
\choiceTF
{\True Mô hình toán học dùng phương trình mô tả quan hệ giữa các đại lượng.}
{Chất điểm luôn là vật có kích thước bằng không trong thực tế.}
{\True Kết quả từ mô hình cần được đối chiếu thực nghiệm.}
{Mô hình càng nhiều chi tiết thừa càng tốt.}
\loigiai{
\begin{itemchoice}
\item (Đúng) Mô hình toán học dùng phương trình mô tả quan hệ giữa các đại lượng. (Vì): Mô hình toán học dùng phương trình mô tả quan hệ giữa các đại lượng. b) S: Chất điểm luôn là vật có kích thước bằng không trong thực tế. c) Đ: Kết quả từ mô hình cần được đối chiếu thực nghiệm. d) S: Mô hình càng nhiều chi tiết thừa càng tốt. Nhận diện mô hình vật chất, mô hình toán học và mô hình lí thuyết; xác định phạm vi áp dụng và kiểm chứng
\item (Sai) Chất điểm luôn là vật có kích thước bằng không trong thực tế.
\item (Đúng) Kết quả từ mô hình cần được đối chiếu thực nghiệm.
\item (Sai) Mô hình càng nhiều chi tiết thừa càng tốt.
\end{itemchoice}
}
\end{ex}

% ===== Câu 126 =====
% ID: L10C1B1-21-TL
% Mức: TH
\dangbt{Vai trò của Vật lí đối với khoa học, kĩ thuật, công nghệ và đời sống}
\begin{bt}
% ID: L10C1B1-21-TL
% Mức: TH
Trình bày kiến thức cốt lõi của dạng “Vai trò của Vật lí đối với khoa học, kĩ thuật, công nghệ và đời sống” và minh họa bằng một ví dụ.
\loigiai{
Phân tích được vai trò của Vật lí qua ví dụ cụ thể và đánh giá cả lợi ích lẫn hệ quả của ứng dụng công nghệ. Ví dụ phù hợp cần nêu rõ đối tượng, dữ kiện và kết luận theo kiến thức SGK.
}
\end{bt}

% ===== Câu 127 =====
% ID: L10C1B1-21-TLN
% Mức: NB
\begin{ex}
% ID: L10C1B1-21-TLN
% Mức: NB
Cho bốn nhận định: (1) Ứng dụng Vật lí có thể vừa mang lợi ích vừa gây tác động môi trường. (2) Siêu âm không sử dụng sóng. (3) Vật lí cung cấp cơ sở cho nhiều công nghệ hiện đại. (4) Mọi ứng dụng công nghệ đều không có rủi ro. Có bao nhiêu nhận định đúng? Chỉ ghi số nguyên.
\shortans{2}
\loigiai{
Nhận định (1), (3) đúng; (2), (4) sai. Vậy có 2 nhận định đúng.
}
\end{ex}

% ===== Câu 128 =====
% ID: L10C1B1-21-TN
% Mức: NB
\dangbt{Phân biệt các loại mô hình trong nghiên cứu Vật lí}
\begin{ex}
% ID: L10C1B1-21-TN
% Mức: NB
Các loại mô hình nào sau đây thường được dùng trong trường phổ thông ?
\choice
{Mô hình toán học}
{\True Cả ba mô hình trên}
{Mô hình lí thuyết}
{Mô hình vật chất}
\loigiai{
Trong nghiên cứu Vật lí ở trường phổ thông, các loại mô hình thường được sử dụng bao gồm: mô hình lí thuyết (dùng để phát biểu các định luật, nguyên lí), mô hình toán học (biểu diễn các định luật dưới dạng công thức toán học) và mô hình vật chất (biểu diễn các đối tượng vật lí bằng các vật thể đơn giản hơn). Do đó, cả ba loại mô hình này đều được dùng.
}
\end{ex}

% ===== Câu 129 =====
% ID: L10C1B1-22-DS
% Mức: NB
\dangbt{Phương pháp thực nghiệm trong nghiên cứu Vật lí}
\begin{ex}
% ID: L10C1B1-22-DS
% Mức: NB
Xét các phát biểu sau về “Phương pháp thực nghiệm trong nghiên cứu Vật lí” (bộ 3).
\choiceTF
{\True Kết quả đo cần được xử lí trước khi kết luận.}
{Không cần kiểm soát điều kiện.}
{\True Biến phụ thuộc là đại lượng được theo dõi theo biến độc lập.}
{Giả thuyết không cần kiểm tra.}
\loigiai{
\begin{itemchoice}
\item (Đúng) Kết quả đo cần được xử lí trước khi kết luận. (Vì): Kết quả đo cần được xử lí trước khi kết luận. b) S: Không cần kiểm soát điều kiện. c) Đ: Biến phụ thuộc là đại lượng được theo dõi theo biến độc lập. d) S: Giả thuyết không cần kiểm tra. Tiến trình cơ bản: xác định vấn đề, dự đoán, thiết kế và thực hiện thí nghiệm, xử lí dữ liệu, kết luận và điều chỉnh
\item (Sai) Không cần kiểm soát điều kiện.
\item (Đúng) Biến phụ thuộc là đại lượng được theo dõi theo biến độc lập.
\item (Sai) Giả thuyết không cần kiểm tra.
\end{itemchoice}
}
\end{ex}

% ===== Câu 130 =====
% ID: L10C1B1-22-TL
% Mức: TH
\dangbt{Vai trò của Vật lí đối với khoa học, kĩ thuật, công nghệ và đời sống}
\begin{bt}
% ID: L10C1B1-22-TL
% Mức: TH
So sánh hai nhận định: “Vật lí hỗ trợ chế tạo thiết bị đo chính xác.” và “Thiết bị đo không cần nguyên lí Vật lí.”.
\loigiai{
Nhận định thứ nhất đúng: Vật lí hỗ trợ chế tạo thiết bị đo chính xác. Nhận định thứ hai sai: Thiết bị đo không cần nguyên lí Vật lí. Sự khác nhau nằm ở điều kiện và bản chất kiến thức được áp dụng.
}
\end{bt}

% ===== Câu 131 =====
% ID: L10C1B1-22-TLN
% Mức: TH
\begin{ex}
% ID: L10C1B1-22-TLN
% Mức: TH
Cho bốn nhận định: (1) Thông tin liên lạc phát triển dựa trên điện từ học và quang học. (2) Thiết bị đo không cần nguyên lí Vật lí. (3) Siêu âm y học có cơ sở từ sự truyền sóng. (4) Tự động hóa không cần cảm biến. Có bao nhiêu nhận định đúng? Chỉ ghi số nguyên.
\shortans{2}
\loigiai{
Nhận định (1), (3) đúng; (2), (4) sai. Vậy có 2 nhận định đúng.
}
\end{ex}

% ===== Câu 132 =====
% ID: L10C1B1-22-TN
% Mức: NB
\dangbt{Phân biệt các loại mô hình trong nghiên cứu Vật lí}
\begin{ex}
% ID: L10C1B1-22-TN
% Mức: NB
Quả địa cầu được sử dụng trong phòng thí nghiệm để mô phỏng hình dạng và một số tính chất của Trái Đất thuộc loại mô hình nào sau đây?
\choice
{Mô hình lí thuyết}
{Mô hình toán học}
{\True Mô hình vật chất}
{Mô hình sinh học}
\loigiai{
Quả địa cầu là một vật thể thu nhỏ mô phỏng Trái Đất thật trong tự nhiên nên nó thuộc loại mô hình vật chất.
}
\end{ex}

% ===== Câu 133 =====
% ID: L10C1B1-23-DS
% Mức: TH
\dangbt{Phương pháp thực nghiệm trong nghiên cứu Vật lí}
\begin{ex}
% ID: L10C1B1-23-DS
% Mức: TH
Xét các phát biểu sau về “Phương pháp thực nghiệm trong nghiên cứu Vật lí” (bộ 4).
\choiceTF
{\True Biến độc lập là đại lượng người nghiên cứu chủ động thay đổi.}
{Không cần ghi chép số liệu.}
{\True Thí nghiệm cần có khả năng lặp lại.}
{Thực nghiệm chỉ cần kết luận, không cần dữ liệu.}
\loigiai{
\begin{itemchoice}
\item (Đúng) Biến độc lập là đại lượng người nghiên cứu chủ động thay đổi. (Vì): Biến độc lập là đại lượng người nghiên cứu chủ động thay đổi. b) S: Không cần ghi chép số liệu. c) Đ: Thí nghiệm cần có khả năng lặp lại. d) S: Thực nghiệm chỉ cần kết luận, không cần dữ liệu. Tiến trình cơ bản: xác định vấn đề, dự đoán, thiết kế và thực hiện thí nghiệm, xử lí dữ liệu, kết luận và điều chỉnh
\item (Sai) Không cần ghi chép số liệu.
\item (Đúng) Thí nghiệm cần có khả năng lặp lại.
\item (Sai) Thực nghiệm chỉ cần kết luận, không cần dữ liệu.
\end{itemchoice}
}
\end{ex}

% ===== Câu 134 =====
% ID: L10C1B1-23-TL
% Mức: VD
\dangbt{Vai trò của Vật lí đối với khoa học, kĩ thuật, công nghệ và đời sống}
\begin{bt}
% ID: L10C1B1-23-TL
% Mức: VD
Phân tích phát biểu “Máy chụp X-quang là một ứng dụng của Vật lí.” và giải thích vì sao phát biểu đó đúng.
\loigiai{
Phát biểu đúng. Máy chụp X-quang là một ứng dụng của Vật lí. Nội dung này phù hợp với kiến thức trọng tâm của dạng bài.
}
\end{bt}

% ===== Câu 135 =====
% ID: L10C1B1-23-TLN
% Mức: VD
\begin{ex}
% ID: L10C1B1-23-TLN
% Mức: VD
Cho bốn nhận định: (1) Kiến thức Vật lí góp phần phát triển năng lượng tái tạo. (2) Vật lí chỉ có vai trò trong phòng thí nghiệm. (3) Kính thiên văn giúp mở rộng nghiên cứu vũ trụ. (4) Vật lí không liên quan đến y học. Có bao nhiêu nhận định đúng? Chỉ ghi số nguyên.
\shortans{2}
\loigiai{
Nhận định (1), (3) đúng; (2), (4) sai. Vậy có 2 nhận định đúng.
}
\end{ex}

% ===== Câu 136 =====
% ID: L10C1B1-23-TN
% Mức: NB
\dangbt{Phân biệt các loại mô hình trong nghiên cứu Vật lí}
\begin{ex}
% ID: L10C1B1-23-TN
% Mức: NB
Khi nghiên cứu chuyển động của một ô tô chạy trên quãng đường dài từ Hà Nội đến Hải Phòng, người ta coi chiếc ô tô đó là một chất điểm. Đây là ví dụ về loại mô hình nào?
\choice
{Mô hình toán học}
{\True Mô hình lí thuyết}
{Mô hình vật chất}
{Mô hình cơ học}
\loigiai{
Chất điểm là một khái niệm trừu tượng, lí tưởng hóa nhằm đơn giản hóa đối tượng nghiên cứu nên thuộc loại mô hình lí thuyết.
}
\end{ex}

% ===== Câu 137 =====
% ID: L10C1B1-24-DS
% Mức: VD
\dangbt{Phương pháp thực nghiệm trong nghiên cứu Vật lí}
\begin{ex}
% ID: L10C1B1-24-DS
% Mức: VD
Xét các phát biểu sau về “Phương pháp thực nghiệm trong nghiên cứu Vật lí” (bộ 1).
\choiceTF
{\True Phương pháp thực nghiệm sử dụng thí nghiệm để kiểm tra dự đoán.}
{Có thể thay đổi đồng thời mọi yếu tố mà vẫn xác định chắc chắn nguyên nhân.}
{\True Phương án thí nghiệm phải xác định đại lượng cần đo.}
{Biến phụ thuộc luôn do người làm tùy ý đặt giá trị.}
\loigiai{
\begin{itemchoice}
\item (Đúng) Phương pháp thực nghiệm sử dụng thí nghiệm để kiểm tra dự đoán. (Vì): Phương pháp thực nghiệm sử dụng thí nghiệm để kiểm tra dự đoán. b) S: Có thể thay đổi đồng thời mọi yếu tố mà vẫn xác định chắc chắn nguyên nhân. c) Đ: Phương án thí nghiệm phải xác định đại lượng cần đo. d) S: Biến phụ thuộc luôn do người làm tùy ý đặt giá trị. Tiến trình cơ bản: xác định vấn đề, dự đoán, thiết kế và thực hiện thí nghiệm, xử lí dữ liệu, kết luận và điều chỉnh
\item (Sai) Có thể thay đổi đồng thời mọi yếu tố mà vẫn xác định chắc chắn nguyên nhân.
\item (Đúng) Phương án thí nghiệm phải xác định đại lượng cần đo.
\item (Sai) Biến phụ thuộc luôn do người làm tùy ý đặt giá trị.
\end{itemchoice}
}
\end{ex}

% ===== Câu 138 =====
% ID: L10C1B1-24-TL
% Mức: VD
\dangbt{Vai trò của Vật lí đối với khoa học, kĩ thuật, công nghệ và đời sống}
\begin{bt}
% ID: L10C1B1-24-TL
% Mức: VD
Xây dựng một quy trình ngắn để giải quyết nhiệm vụ thuộc dạng “Vai trò của Vật lí đối với khoa học, kĩ thuật, công nghệ và đời sống”.
\loigiai{
Quy trình: đọc yêu cầu; xác định khái niệm hoặc đại lượng; chọn nguyên tắc/công thức; xử lí dữ kiện; kiểm tra đơn vị và điều kiện; kết luận. Phân tích được vai trò của Vật lí qua ví dụ cụ thể và đánh giá cả lợi ích lẫn hệ quả của ứng dụng công nghệ.
}
\end{bt}

% ===== Câu 139 =====
% ID: L10C1B1-24-TLN
% Mức: VDC
\begin{ex}
% ID: L10C1B1-24-TLN
% Mức: VDC
Cho bốn nhận định: (1) Vật lí cung cấp cơ sở cho nhiều công nghệ hiện đại. (2) Mọi ứng dụng công nghệ đều không có rủi ro. (3) Ứng dụng Vật lí có thể vừa mang lợi ích vừa gây tác động môi trường. (4) Siêu âm không sử dụng sóng. Có bao nhiêu nhận định đúng? Chỉ ghi số nguyên.
\shortans{2}
\loigiai{
Nhận định (1), (3) đúng; (2), (4) sai. Vậy có 2 nhận định đúng.
}
\end{ex}

% ===== Câu 140 =====
% ID: L10C1B1-24-TN
% Mức: NB
\dangbt{Phân tích vai trò và ảnh hưởng hai mặt của Vật lí đối với đời sống, khoa học và công nghệ}
\begin{ex}
% ID: L10C1B1-24-TN
% Mức: NB
Việc lắp ráp pin cho nhà máy điện mặt trời thể hiện vai trò nào sau đây?
\choice
{Chăm sóc đời sống con người}
{\True Ứng dụng công nghệ vào đời sống, sản xuất}
{Nâng cao hiểu biết của con người về tự nhiên}
{Nghiên cứu khoa học}
\loigiai{
Lắp đặt pin năng lượng mặt trời là ứng dụng thực tế của công nghệ vật lí phục vụ sản xuất năng lượng.
}
\end{ex}

% ===== Câu 141 =====
% ID: L10C1B1-25-DS
% Mức: VD
\dangbt{Phương pháp thực nghiệm trong nghiên cứu Vật lí}
\begin{ex}
% ID: L10C1B1-25-DS
% Mức: VD
Xét các phát biểu sau về “Phương pháp thực nghiệm trong nghiên cứu Vật lí” (bộ 5).
\choiceTF
{\True Các yếu tố khác nên được kiểm soát.}
{Thí nghiệm không cần lặp lại.}
{\True Bước đầu cần xác định vấn đề nghiên cứu.}
{Kết quả trái dự đoán phải bị bỏ đi.}
\loigiai{
\begin{itemchoice}
\item (Đúng) Các yếu tố khác nên được kiểm soát. (Vì): Các yếu tố khác nên được kiểm soát. b) S: Thí nghiệm không cần lặp lại. c) Đ: Bước đầu cần xác định vấn đề nghiên cứu. d) S: Kết quả trái dự đoán phải bị bỏ đi. Tiến trình cơ bản: xác định vấn đề, dự đoán, thiết kế và thực hiện thí nghiệm, xử lí dữ liệu, kết luận và điều chỉnh
\item (Sai) Thí nghiệm không cần lặp lại.
\item (Đúng) Bước đầu cần xác định vấn đề nghiên cứu.
\item (Sai) Kết quả trái dự đoán phải bị bỏ đi.
\end{itemchoice}
}
\end{ex}

% ===== Câu 142 =====
% ID: L10C1B1-25-TL
% Mức: VDC
\dangbt{Vai trò của Vật lí đối với khoa học, kĩ thuật, công nghệ và đời sống}
\begin{bt}
% ID: L10C1B1-25-TL
% Mức: VDC
Phân tích phát biểu “Mọi ứng dụng công nghệ đều không có rủi ro.”, chỉ ra sai sót và sửa lại.
\loigiai{
Phát biểu sai: Mọi ứng dụng công nghệ đều không có rủi ro. Cần sửa thành: Thông tin liên lạc phát triển dựa trên điện từ học và quang học. Phân tích được vai trò của Vật lí qua ví dụ cụ thể và đánh giá cả lợi ích lẫn hệ quả của ứng dụng công nghệ.
}
\end{bt}

% ===== Câu 143 =====
% ID: L10C1B1-25-TLN
% Mức: VDC
\begin{ex}
% ID: L10C1B1-25-TLN
% Mức: VDC
Cho bốn nhận định: (1) Vật lí hỗ trợ chế tạo thiết bị đo chính xác. (2) Bán dẫn là thành tựu không liên quan Vật lí. (3) Vật lí góp phần tự động hóa sản xuất. (4) Internet không sử dụng thành tựu Vật lí. Có bao nhiêu nhận định đúng? Chỉ ghi số nguyên.
\shortans{2}
\loigiai{
Nhận định (1), (3) đúng; (2), (4) sai. Vậy có 2 nhận định đúng.
}
\end{ex}

% ===== Câu 144 =====
% ID: L10C1B1-25-TN
% Mức: NB
\dangbt{Phân tích vai trò và ảnh hưởng hai mặt của Vật lí đối với đời sống, khoa học và công nghệ}
\begin{ex}
% ID: L10C1B1-25-TN
% Mức: NB
Cuộc cách mạng công nghiệp lần thứ nhất khởi đầu vào thế kỉ XVIII gắn liền với sự phát minh ra thiết bị nào sau đây dựa trên những nghiên cứu về Nhiệt của Vật lí?
\choice
{Động cơ điện}
{\True Máy hơi nước}
{Chất bán dẫn}
{Trí tuệ nhân tạo}
\loigiai{
Sự phát minh ra máy hơi nước của James Watt năm $1765$ dựa trên những nghiên cứu về Nhiệt của Vật lí đã tạo nên bước khởi đầu cho cuộc cách mạng công nghiệp lần thứ nhất, thay thế sức lực cơ bắp bằng sức lực máy móc.
}
\end{ex}

% ===== Câu 145 =====
% ID: L10C1B1-26-DS
% Mức: VDC
\dangbt{Phương pháp thực nghiệm trong nghiên cứu Vật lí}
\begin{ex}
% ID: L10C1B1-26-DS
% Mức: VDC
Xét các phát biểu sau về “Phương pháp thực nghiệm trong nghiên cứu Vật lí” (bộ 2).
\choiceTF
{\True Sau khi nêu vấn đề có thể đề xuất giả thuyết hoặc dự đoán.}
{Không cần xác định vấn đề trước khi đo.}
{\True Nếu kết quả không phù hợp, cần xem lại giả thuyết hoặc phương án.}
{Một phép đo duy nhất luôn tuyệt đối chính xác.}
\loigiai{
\begin{itemchoice}
\item (Đúng) Sau khi nêu vấn đề có thể đề xuất giả thuyết hoặc dự đoán. (Vì): Sau khi nêu vấn đề có thể đề xuất giả thuyết hoặc dự đoán. b) S: Không cần xác định vấn đề trước khi đo. c) Đ: Nếu kết quả không phù hợp, cần xem lại giả thuyết hoặc phương án. d) S: Một phép đo duy nhất luôn tuyệt đối chính xác. Tiến trình cơ bản: xác định vấn đề, dự đoán, thiết kế và thực hiện thí nghiệm, xử lí dữ liệu, kết luận và điều chỉnh
\item (Sai) Không cần xác định vấn đề trước khi đo.
\item (Đúng) Nếu kết quả không phù hợp, cần xem lại giả thuyết hoặc phương án.
\item (Sai) Một phép đo duy nhất luôn tuyệt đối chính xác.
\end{itemchoice}
}
\end{ex}

% ===== Câu 146 =====
% ID: L10C1B1-26-TL
% Mức: NB
\dangbt{Vận dụng tiến trình nghiên cứu Vật lí vào tình huống thực tiễn}
\begin{bt}
% ID: L10C1B1-26-TL
% Mức: NB
Đề xuất cách vận dụng nội dung “Vận dụng tiến trình nghiên cứu Vật lí vào tình huống thực tiễn” trong một tình huống học tập hoặc thực tiễn.
\loigiai{
Cần xác định tình huống, chọn kiến thức phù hợp, thực hiện theo quy trình, kiểm tra điều kiện và đánh giá kết quả. Vận dụng đầy đủ chu trình nghiên cứu vào vấn đề gần gũi, từ câu hỏi đến đánh giá và cải tiến.
}
\end{bt}

% ===== Câu 147 =====
% ID: L10C1B1-26-TLN
% Mức: NB
\dangbt{Vận dụng tiến trình của phương pháp mô hình}
\begin{ex}
% ID: L10C1B1-26-TLN
% Mức: NB
Trong tiến trình của phương pháp mô hình, bước thứ ba ngay sau bước "Xây dựng mô hình (giả thuyết)" là bước nào?
\shortans{Kiểm tra sự phù hợp của mô hình}
\loigiai{
Theo sơ đồ 4 bước của tiến trình phương pháp mô hình: Xác định đối tượng $\rightarrow$ Xây dựng mô hình $\rightarrow$ Kiểm tra sự phù hợp $\rightarrow$ Kết luận.
}
\end{ex}

% ===== Câu 148 =====
% ID: L10C1B1-26-TN
% Mức: NB
\dangbt{Phân tích vai trò và ảnh hưởng hai mặt của Vật lí đối với đời sống, khoa học và công nghệ}
\begin{ex}
% ID: L10C1B1-26-TN
% Mức: NB
Ứng dụng nào sau đây của Vật lí có liên quan đến lĩnh vực y tế?
\choice
{Internet vạn vật (IoT)}
{Cảm biến không dây trong nông nghiệp}
{\True Chụp cộng hưởng từ (MRI)}
{Điện toán đám mây}
\loigiai{
MRI (Magnetic Resonance Imaging – chụp cộng hưởng từ) là một ứng dụng nổi bật của vật lí trong chẩn đoán và điều trị y học.
}
\end{ex}

% ===== Câu 149 =====
% ID: L10C1B1-27-DS
% Mức: NB
\dangbt{Quá trình phát triển của Vật lí}
\begin{ex}
% ID: L10C1B1-27-DS
% Mức: NB
Xét các phát biểu sau về “Quá trình phát triển của Vật lí” (bộ 3).
\choiceTF
{\True Một lí thuyết khoa học cần được đối chiếu với bằng chứng.}
{Newton không liên quan cơ học.}
{\True Newton có đóng góp quan trọng cho cơ học cổ điển.}
{Khoa học không bao giờ điều chỉnh lí thuyết.}
\loigiai{
\begin{itemchoice}
\item (Đúng) Một lí thuyết khoa học cần được đối chiếu với bằng chứng. (Vì): Một lí thuyết khoa học cần được đối chiếu với bằng chứng. b) S: Newton không liên quan cơ học. c) Đ: Newton có đóng góp quan trọng cho cơ học cổ điển. d) S: Khoa học không bao giờ điều chỉnh lí thuyết. Nhận biết các giai đoạn và đặc trưng phát triển của Vật lí, nhấn mạnh quan hệ giữa lí thuyết, thực nghiệm và thiết bị
\item (Sai) Newton không liên quan cơ học.
\item (Đúng) Newton có đóng góp quan trọng cho cơ học cổ điển.
\item (Sai) Khoa học không bao giờ điều chỉnh lí thuyết.
\end{itemchoice}
}
\end{ex}

% ===== Câu 150 =====
% ID: L10C1B1-27-TL
% Mức: TH
\dangbt{Vận dụng tiến trình nghiên cứu Vật lí vào tình huống thực tiễn}
\begin{bt}
% ID: L10C1B1-27-TL
% Mức: TH
Trình bày kiến thức cốt lõi của dạng “Vận dụng tiến trình nghiên cứu Vật lí vào tình huống thực tiễn” và minh họa bằng một ví dụ.
\loigiai{
Vận dụng đầy đủ chu trình nghiên cứu vào vấn đề gần gũi, từ câu hỏi đến đánh giá và cải tiến. Ví dụ phù hợp cần nêu rõ đối tượng, dữ kiện và kết luận theo kiến thức SGK.
}
\end{bt}

% ===== Câu 151 =====
% ID: L10C1B1-27-TLN
% Mức: NB
\begin{ex}
% ID: L10C1B1-27-TLN
% Mức: NB
Cho bốn nhận định: (1) Dự đoán nên dựa trên kiến thức hoặc quan sát ban đầu. (2) Không cần ghi đơn vị trong bảng số liệu. (3) Cần nêu nguồn sai số và giới hạn của kết quả. (4) Kế hoạch không cần dụng cụ. Có bao nhiêu nhận định đúng? Chỉ ghi số nguyên.
\shortans{2}
\loigiai{
Nhận định (1), (3) đúng; (2), (4) sai. Vậy có 2 nhận định đúng.
}
\end{ex}

% ===== Câu 152 =====
% ID: L10C1B1-27-TN
% Mức: NB
\dangbt{Phân tích vai trò và ảnh hưởng hai mặt của Vật lí đối với đời sống, khoa học và công nghệ}
\begin{ex}
% ID: L10C1B1-27-TN
% Mức: NB
Phát biểu nào sau đây thể hiện ảnh hưởng tiêu cực của việc ứng dụng các thành tựu Vật lí vào công nghệ?
\choice
{Giải phóng sức lao động cơ bắp của con người nhờ máy móc}
{Mở rộng phương thức kết nối thông tin liên lạc toàn cầu}
{\True Khí thải từ các nhà máy công nghiệp gây ra hiệu ứng nhà kính và biến đổi khí hậu}
{Sử dụng năng lượng sạch như điện gió và năng lượng mặt trời}
\loigiai{
Khí thải từ các nhà máy công nghiệp đốt nhiên liệu hóa thạch gây ra ô nhiễm không khí và hiệu ứng nhà kính là một ví dụ điển hình về ảnh hưởng tiêu cực của việc ứng dụng công nghệ nếu không được quản lí và bảo vệ môi trường phù hợp.
}
\end{ex}

% ===== Câu 153 =====
% ID: L10C1B1-28-DS
% Mức: TH
\dangbt{Quá trình phát triển của Vật lí}
\begin{ex}
% ID: L10C1B1-28-DS
% Mức: TH
Xét các phát biểu sau về “Quá trình phát triển của Vật lí” (bộ 4).
\choiceTF
{\True Quan sát và đo lường thúc đẩy sự phát triển của Vật lí.}
{Thiết bị đo cản trở mọi phát hiện.}
{\True Sự phát triển thiết bị tạo điều kiện cho phát hiện mới.}
{Vật lí chỉ phát triển nhờ suy luận chủ quan.}
\loigiai{
\begin{itemchoice}
\item (Đúng) Quan sát và đo lường thúc đẩy sự phát triển của Vật lí. (Vì): Quan sát và đo lường thúc đẩy sự phát triển của Vật lí. b) S: Thiết bị đo cản trở mọi phát hiện. c) Đ: Sự phát triển thiết bị tạo điều kiện cho phát hiện mới. d) S: Vật lí chỉ phát triển nhờ suy luận chủ quan. Nhận biết các giai đoạn và đặc trưng phát triển của Vật lí, nhấn mạnh quan hệ giữa lí thuyết, thực nghiệm và thiết bị
\item (Sai) Thiết bị đo cản trở mọi phát hiện.
\item (Đúng) Sự phát triển thiết bị tạo điều kiện cho phát hiện mới.
\item (Sai) Vật lí chỉ phát triển nhờ suy luận chủ quan.
\end{itemchoice}
}
\end{ex}

% ===== Câu 154 =====
% ID: L10C1B1-28-TL
% Mức: TH
\dangbt{Vận dụng tiến trình nghiên cứu Vật lí vào tình huống thực tiễn}
\begin{bt}
% ID: L10C1B1-28-TL
% Mức: TH
So sánh hai nhận định: “Đồ thị có thể làm rõ quan hệ giữa hai đại lượng.” và “Đồ thị không thể hiện quan hệ đại lượng.”.
\loigiai{
Nhận định thứ nhất đúng: Đồ thị có thể làm rõ quan hệ giữa hai đại lượng. Nhận định thứ hai sai: Đồ thị không thể hiện quan hệ đại lượng. Sự khác nhau nằm ở điều kiện và bản chất kiến thức được áp dụng.
}
\end{bt}

% ===== Câu 155 =====
% ID: L10C1B1-28-TLN
% Mức: NB
\begin{ex}
% ID: L10C1B1-28-TLN
% Mức: NB
Cho bốn nhận định: (1) Kết luận phải dựa trên dữ liệu. (2) Có thể chọn bỏ dữ liệu tùy ý. (3) Câu hỏi nghiên cứu cần rõ ràng và có thể kiểm tra. (4) Không cần quan tâm an toàn. Có bao nhiêu nhận định đúng? Chỉ ghi số nguyên.
\shortans{2}
\loigiai{
Nhận định (1), (3) đúng; (2), (4) sai. Vậy có 2 nhận định đúng.
}
\end{ex}

% ===== Câu 156 =====
% ID: L10C1B1-28-TN
% Mức: NB
\dangbt{Phân tích vai trò và ảnh hưởng hai mặt của Vật lí đối với đời sống, khoa học và công nghệ}
\begin{ex}
% ID: L10C1B1-28-TN
% Mức: NB
Kết luận sai về ảnh hưởng của vật lí đến một số lĩnh vực trong đời sống và kĩ thuật
\choice
{\True Vật lí đem lại cho con người những lợi ích tuyệt vời và không gây ra một ảnh hưởng xấu nào}
{Vật lí ảnh hưởng mạnh mẽ và có tác động làm thay đổi mọi lĩnh vực hoạt động của con người}
{Kiến thức vật lí trong các phân ngành được áp dụng kết hợp để tạo ra kết quả tối ưu}
{Vật lí là cơ sở của khoa học tự nhiên và công nghệ}
\loigiai{
Vật lí không chỉ đem lại lợi ích mà còn có thể gây ra hậu quả xấu nếu bị lạm dụng (như vũ khí hạt nhân).
}
\end{ex}

% ===== Câu 157 =====
% ID: L10C1B1-29-DS
% Mức: VD
\dangbt{Quá trình phát triển của Vật lí}
\begin{ex}
% ID: L10C1B1-29-DS
% Mức: VD
Xét các phát biểu sau về “Quá trình phát triển của Vật lí” (bộ 1).
\choiceTF
{\True Thực nghiệm giúp kiểm tra các dự đoán khoa học.}
{Thí nghiệm không có vai trò kiểm chứng.}
{\True Vật lí hiện đại nghiên cứu sâu thế giới vi mô.}
{Mọi quan niệm cổ đều luôn đúng.}
\loigiai{
\begin{itemchoice}
\item (Đúng) Thực nghiệm giúp kiểm tra các dự đoán khoa học. (Vì): Thực nghiệm giúp kiểm tra các dự đoán khoa học. b) S: Thí nghiệm không có vai trò kiểm chứng. c) Đ: Vật lí hiện đại nghiên cứu sâu thế giới vi mô. d) S: Mọi quan niệm cổ đều luôn đúng. Nhận biết các giai đoạn và đặc trưng phát triển của Vật lí, nhấn mạnh quan hệ giữa lí thuyết, thực nghiệm và thiết bị
\item (Sai) Thí nghiệm không có vai trò kiểm chứng.
\item (Đúng) Vật lí hiện đại nghiên cứu sâu thế giới vi mô.
\item (Sai) Mọi quan niệm cổ đều luôn đúng.
\end{itemchoice}
}
\end{ex}

% ===== Câu 158 =====
% ID: L10C1B1-29-TL
% Mức: VD
\dangbt{Vận dụng tiến trình nghiên cứu Vật lí vào tình huống thực tiễn}
\begin{bt}
% ID: L10C1B1-29-TL
% Mức: VD
Phân tích phát biểu “Dự đoán nên dựa trên kiến thức hoặc quan sát ban đầu.” và giải thích vì sao phát biểu đó đúng.
\loigiai{
Phát biểu đúng. Dự đoán nên dựa trên kiến thức hoặc quan sát ban đầu. Nội dung này phù hợp với kiến thức trọng tâm của dạng bài.
}
\end{bt}

% ===== Câu 159 =====
% ID: L10C1B1-29-TLN
% Mức: TH
\begin{ex}
% ID: L10C1B1-29-TLN
% Mức: TH
Cho bốn nhận định: (1) Kế hoạch cần nêu dụng cụ và cách thu thập dữ liệu. (2) Đồ thị không thể hiện quan hệ đại lượng. (3) Lặp phép đo giúp đánh giá độ ổn định. (4) Không được cải tiến phương án. Có bao nhiêu nhận định đúng? Chỉ ghi số nguyên.
\shortans{2}
\loigiai{
Nhận định (1), (3) đúng; (2), (4) sai. Vậy có 2 nhận định đúng.
}
\end{ex}

% ===== Câu 160 =====
% ID: L10C1B1-29-TN
% Mức: TH
\dangbt{Phân tích vai trò và ảnh hưởng hai mặt của Vật lí đối với đời sống, khoa học và công nghệ}
\begin{ex}
% ID: L10C1B1-29-TN
% Mức: TH
Phát biểu nào sau đây mô tả đúng nhất vai trò của Vật lí trong đời sống và khoa học?
\choice
{Vật lí chỉ nghiên cứu các hiện tượng tự nhiên, không liên quan đến đời sống con người.}
{\True Vật lí là ngành khoa học cơ bản nghiên cứu về các quy luật phổ biến của tự nhiên, từ đó ứng dụng vào nhiều lĩnh vực của đời sống và khoa học kĩ thuật.}
{Vật lí chỉ tập trung vào các hiện tượng vĩ mô, bỏ qua các hiện tượng vi mô.}
{Vật lí là một ngành khoa học tĩnh, không có sự phát triển và thay đổi theo thời gian.}
\loigiai{
Vật lí cung cấp những nguyên lý cơ bản giải thích nhiều hiện tượng trong tự nhiên và là cơ sở lý thuyết cho sự phát triển của các ngành khoa học khác như Hóa học, Sinh học, Khoa học Trái Đất, và đặc biệt là Kỹ thuật, Công nghệ. Nhiều ứng dụng công nghệ hiện đại đều dựa trên các khám phá và định luật vật lí.
}
\end{ex}

% ===== Câu 161 =====
% ID: L10C1B1-30-DS
% Mức: VD
\dangbt{Quá trình phát triển của Vật lí}
\begin{ex}
% ID: L10C1B1-30-DS
% Mức: VD
Xét các phát biểu sau về “Quá trình phát triển của Vật lí” (bộ 5).
\choiceTF
{\True Thuyết lượng tử thuộc Vật lí hiện đại.}
{Quan sát không cung cấp dữ liệu.}
{\True Galileo góp phần đặt nền móng cho phương pháp thực nghiệm.}
{Vật lí hiện đại không nghiên cứu vi mô.}
\loigiai{
\begin{itemchoice}
\item (Đúng) Thuyết lượng tử thuộc Vật lí hiện đại. (Vì): Thuyết lượng tử thuộc Vật lí hiện đại. b) S: Quan sát không cung cấp dữ liệu. c) Đ: Galileo góp phần đặt nền móng cho phương pháp thực nghiệm. d) S: Vật lí hiện đại không nghiên cứu vi mô. Nhận biết các giai đoạn và đặc trưng phát triển của Vật lí, nhấn mạnh quan hệ giữa lí thuyết, thực nghiệm và thiết bị
\item (Sai) Quan sát không cung cấp dữ liệu.
\item (Đúng) Galileo góp phần đặt nền móng cho phương pháp thực nghiệm.
\item (Sai) Vật lí hiện đại không nghiên cứu vi mô.
\end{itemchoice}
}
\end{ex}

% ===== Câu 162 =====
% ID: L10C1B1-30-TL
% Mức: VD
\dangbt{Vận dụng tiến trình nghiên cứu Vật lí vào tình huống thực tiễn}
\begin{bt}
% ID: L10C1B1-30-TL
% Mức: VD
Xây dựng một quy trình ngắn để giải quyết nhiệm vụ thuộc dạng “Vận dụng tiến trình nghiên cứu Vật lí vào tình huống thực tiễn”.
\loigiai{
Quy trình: đọc yêu cầu; xác định khái niệm hoặc đại lượng; chọn nguyên tắc/công thức; xử lí dữ kiện; kiểm tra đơn vị và điều kiện; kết luận. Vận dụng đầy đủ chu trình nghiên cứu vào vấn đề gần gũi, từ câu hỏi đến đánh giá và cải tiến.
}
\end{bt}

% ===== Câu 163 =====
% ID: L10C1B1-30-TLN
% Mức: VD
\begin{ex}
% ID: L10C1B1-30-TLN
% Mức: VD
Cho bốn nhận định: (1) Bảng số liệu giúp tổ chức kết quả đo. (2) Lặp đo luôn vô ích. (3) Điều kiện an toàn phải được đưa vào kế hoạch. (4) Câu hỏi càng mơ hồ càng dễ nghiên cứu. Có bao nhiêu nhận định đúng? Chỉ ghi số nguyên.
\shortans{2}
\loigiai{
Nhận định (1), (3) đúng; (2), (4) sai. Vậy có 2 nhận định đúng.
}
\end{ex}

% ===== Câu 164 =====
% ID: L10C1B1-30-TN
% Mức: TH
\dangbt{Phân tích vai trò và ảnh hưởng hai mặt của Vật lí đối với đời sống, khoa học và công nghệ}
\begin{ex}
% ID: L10C1B1-30-TN
% Mức: TH
Câu nào sau đây đúng về mối quan hệ giữa Vật lý và các khoa học khác:
\choice
{Vật lý độc lập với các khoa học khác}
{\True Vật lý là nền tảng cho nhiều khoa học khác}
{Các khoa học khác không ảnh hưởng đến Vật lý}
{Vật lý chỉ liên quan đến toán học}
\loigiai{
Vật lý là nền tảng cho nhiều khoa học khác như hóa học, sinh học, và kỹ thuật, vì nó cung cấp các quy luật cơ bản của tự nhiên.
}
\end{ex}

% ===== Câu 165 =====
% ID: L10C1B1-31-DS
% Mức: VDC
\dangbt{Quá trình phát triển của Vật lí}
\begin{ex}
% ID: L10C1B1-31-DS
% Mức: VDC
Xét các phát biểu sau về “Quá trình phát triển của Vật lí” (bộ 2).
\choiceTF
{\True Vật lí cổ điển phát triển mạnh từ thế kỉ XVII.}
{Lí thuyết khoa học không cần bằng chứng.}
{\True Kết quả mới có thể làm lí thuyết được bổ sung hoặc điều chỉnh.}
{Thuyết lượng tử thuộc thời tiền Vật lí.}
\loigiai{
\begin{itemchoice}
\item (Đúng) Vật lí cổ điển phát triển mạnh từ thế kỉ XVII. (Vì): Vật lí cổ điển phát triển mạnh từ thế kỉ XVII. b) S: Lí thuyết khoa học không cần bằng chứng. c) Đ: Kết quả mới có thể làm lí thuyết được bổ sung hoặc điều chỉnh. d) S: Thuyết lượng tử thuộc thời tiền Vật lí. Nhận biết các giai đoạn và đặc trưng phát triển của Vật lí, nhấn mạnh quan hệ giữa lí thuyết, thực nghiệm và thiết bị
\item (Sai) Lí thuyết khoa học không cần bằng chứng.
\item (Đúng) Kết quả mới có thể làm lí thuyết được bổ sung hoặc điều chỉnh.
\item (Sai) Thuyết lượng tử thuộc thời tiền Vật lí.
\end{itemchoice}
}
\end{ex}

% ===== Câu 166 =====
% ID: L10C1B1-31-TL
% Mức: VDC
\dangbt{Vận dụng tiến trình nghiên cứu Vật lí vào tình huống thực tiễn}
\begin{bt}
% ID: L10C1B1-31-TL
% Mức: VDC
Phân tích phát biểu “Không cần quan tâm an toàn.”, chỉ ra sai sót và sửa lại.
\loigiai{
Phát biểu sai: Không cần quan tâm an toàn. Cần sửa thành: Kế hoạch cần nêu dụng cụ và cách thu thập dữ liệu. Vận dụng đầy đủ chu trình nghiên cứu vào vấn đề gần gũi, từ câu hỏi đến đánh giá và cải tiến.
}
\end{bt}

% ===== Câu 167 =====
% ID: L10C1B1-31-TLN
% Mức: VDC
\begin{ex}
% ID: L10C1B1-31-TLN
% Mức: VDC
Cho bốn nhận định: (1) Câu hỏi nghiên cứu cần rõ ràng và có thể kiểm tra. (2) Không cần quan tâm an toàn. (3) Kết luận phải dựa trên dữ liệu. (4) Có thể chọn bỏ dữ liệu tùy ý. Có bao nhiêu nhận định đúng? Chỉ ghi số nguyên.
\shortans{2}
\loigiai{
Nhận định (1), (3) đúng; (2), (4) sai. Vậy có 2 nhận định đúng.
}
\end{ex}

% ===== Câu 168 =====
% ID: L10C1B1-31-TN
% Mức: TH
\dangbt{Phân tích vai trò và ảnh hưởng hai mặt của Vật lí đối với đời sống, khoa học và công nghệ}
\begin{ex}
% ID: L10C1B1-31-TN
% Mức: TH
Vật lý đóng vai trò quan trọng trong sự phát triển của:
\choice
{Công nghệ thông tin}
{$Y$ học}
{Năng lượng}
{\True Tất cả các đáp án trên}
\loigiai{
Vật lý đóng vai trò quan trọng trong nhiều lĩnh vực như công nghệ thông tin, y học, và năng lượng, giúp cải thiện chất lượng cuộc sống và phát triển kinh tế.
}
\end{ex}

% ===== Câu 169 =====
% ID: L10C1B1-32-DS
% Mức: NB
\dangbt{Vai trò của Vật lí đối với khoa học, kĩ thuật, công nghệ và đời sống}
\begin{ex}
% ID: L10C1B1-32-DS
% Mức: NB
Xét các phát biểu sau về “Vai trò của Vật lí đối với khoa học, kĩ thuật, công nghệ và đời sống” (bộ 3).
\choiceTF
{\True Vật lí hỗ trợ chế tạo thiết bị đo chính xác.}
{Vật lí chỉ có vai trò trong phòng thí nghiệm.}
{\True Siêu âm y học có cơ sở từ sự truyền sóng.}
{Kính thiên văn dùng để nghiên cứu vi khuẩn.}
\loigiai{
\begin{itemchoice}
\item (Đúng) Vật lí hỗ trợ chế tạo thiết bị đo chính xác. (Vì): Vật lí hỗ trợ chế tạo thiết bị đo chính xác. b) S: Vật lí chỉ có vai trò trong phòng thí nghiệm. c) Đ: Siêu âm y học có cơ sở từ sự truyền sóng. d) S: Kính thiên văn dùng để nghiên cứu vi khuẩn. Phân tích được vai trò của Vật lí qua ví dụ cụ thể và đánh giá cả lợi ích lẫn hệ quả của ứng dụng công nghệ
\item (Sai) Vật lí chỉ có vai trò trong phòng thí nghiệm.
\item (Đúng) Siêu âm y học có cơ sở từ sự truyền sóng.
\item (Sai) Kính thiên văn dùng để nghiên cứu vi khuẩn.
\end{itemchoice}
}
\end{ex}

% ===== Câu 170 =====
% ID: L10C1B1-32-TL
% Mức: NB
\dangbt{Đối tượng nghiên cứu và mục tiêu của môn Vật lí}
\begin{bt}
% ID: L10C1B1-32-TL
% Mức: NB
Đề xuất cách vận dụng nội dung “Đối tượng nghiên cứu và mục tiêu của môn Vật lí” trong một tình huống học tập hoặc thực tiễn.
\loigiai{
Cần xác định tình huống, chọn kiến thức phù hợp, thực hiện theo quy trình, kiểm tra điều kiện và đánh giá kết quả. Xác định đúng đối tượng, phạm vi và mục tiêu của Vật lí; phân biệt với đối tượng trung tâm của các môn khoa học khác.
}
\end{bt}

% ===== Câu 171 =====
% ID: L10C1B1-32-TLN
% Mức: VDC
\dangbt{Vận dụng tiến trình nghiên cứu Vật lí vào tình huống thực tiễn}
\begin{ex}
% ID: L10C1B1-32-TLN
% Mức: VDC
Cho bốn nhận định: (1) Đồ thị có thể làm rõ quan hệ giữa hai đại lượng. (2) Sai số không cần được xem xét. (3) Có thể đề xuất cải tiến sau khi đánh giá kết quả. (4) Kết luận có thể trái dữ liệu mà không cần giải thích. Có bao nhiêu nhận định đúng? Chỉ ghi số nguyên.
\shortans{2}
\loigiai{
Nhận định (1), (3) đúng; (2), (4) sai. Vậy có 2 nhận định đúng.
}
\end{ex}

% ===== Câu 172 =====
% ID: L10C1B1-32-TN
% Mức: TH
\dangbt{Phân tích vai trò và ảnh hưởng hai mặt của Vật lí đối với đời sống, khoa học và công nghệ}
\begin{ex}
% ID: L10C1B1-32-TN
% Mức: TH
Vật lý có vai trò quan trọng trong việc giải quyết các vấn đề toàn cầu như:
\choice
{Biến đổi khí hậu}
{Năng lượng sạch}
{$Y$ tế}
{\True Tất cả các đáp án trên}
\loigiai{
Vật lý có vai trò quan trọng trong việc giải quyết các vấn đề toàn cầu, như phát triển năng lượng sạch, hiểu rõ biến đổi khí hậu, và cải thiện công nghệ y tế.
}
\end{ex}

% ===== Câu 173 =====
% ID: L10C1B1-33-DS
% Mức: TH
\dangbt{Vai trò của Vật lí đối với khoa học, kĩ thuật, công nghệ và đời sống}
\begin{ex}
% ID: L10C1B1-33-DS
% Mức: TH
Xét các phát biểu sau về “Vai trò của Vật lí đối với khoa học, kĩ thuật, công nghệ và đời sống” (bộ 4).
\choiceTF
{\True Công nghệ bán dẫn dựa trên hiểu biết Vật lí chất rắn.}
{Siêu âm không sử dụng sóng.}
{\True Vật lí góp phần tự động hóa sản xuất.}
{Vật lí không liên quan đến y học.}
\loigiai{
\begin{itemchoice}
\item (Đúng) Công nghệ bán dẫn dựa trên hiểu biết Vật lí chất rắn. (Vì): Công nghệ bán dẫn dựa trên hiểu biết Vật lí chất rắn. b) S: Siêu âm không sử dụng sóng. c) Đ: Vật lí góp phần tự động hóa sản xuất. d) S: Vật lí không liên quan đến y học. Phân tích được vai trò của Vật lí qua ví dụ cụ thể và đánh giá cả lợi ích lẫn hệ quả của ứng dụng công nghệ
\item (Sai) Siêu âm không sử dụng sóng.
\item (Đúng) Vật lí góp phần tự động hóa sản xuất.
\item (Sai) Vật lí không liên quan đến y học.
\end{itemchoice}
}
\end{ex}

% ===== Câu 174 =====
% ID: L10C1B1-33-TL
% Mức: TH
\dangbt{Đối tượng nghiên cứu và mục tiêu của môn Vật lí}
\begin{bt}
% ID: L10C1B1-33-TL
% Mức: TH
Trình bày kiến thức cốt lõi của dạng “Đối tượng nghiên cứu và mục tiêu của môn Vật lí” và minh họa bằng một ví dụ.
\loigiai{
Xác định đúng đối tượng, phạm vi và mục tiêu của Vật lí; phân biệt với đối tượng trung tâm của các môn khoa học khác. Ví dụ phù hợp cần nêu rõ đối tượng, dữ kiện và kết luận theo kiến thức SGK.
}
\end{bt}

% ===== Câu 175 =====
% ID: L10C1B1-33-TLN
% Mức: NB
\dangbt{Xác định đối tượng nghiên cứu và mục tiêu của môn Vật lí}
\begin{ex}
% ID: L10C1B1-33-TLN
% Mức: NB
Ngành khoa học tự nhiên nào nghiên cứu về các dạng vận động của vật chất và năng lượng được gọi là gì?
\shortans{Vật lí học}
\loigiai{
Theo định nghĩa trong sách giáo khoa, ngành khoa học nghiên cứu về các dạng vận động của vật chất và năng lượng được gọi là Vật lí học.
}
\end{ex}

% ===== Câu 176 =====
% ID: L10C1B1-33-TN
% Mức: NB
\dangbt{Phương pháp mô hình trong nghiên cứu Vật lí}
\begin{ex}
% ID: L10C1B1-33-TN
% Mức: NB
Trong dạng “Phương pháp mô hình trong nghiên cứu Vật lí”, phát biểu nào sau đây đúng?
\choice
{\True Mô hình là sự biểu diễn đơn giản hóa đối tượng hoặc quá trình.}
{Mô hình phải giống hoàn toàn vật thật ở mọi chi tiết.}
{Chất điểm luôn là vật có kích thước bằng không trong thực tế.}
{Mô hình càng nhiều chi tiết thừa càng tốt.}
\loigiai{
Phát biểu đúng là: Mô hình là sự biểu diễn đơn giản hóa đối tượng hoặc quá trình. Nhận diện mô hình vật chất, mô hình toán học và mô hình lí thuyết; xác định phạm vi áp dụng và kiểm chứng.
}
\end{ex}

% ===== Câu 177 =====
% ID: L10C1B1-34-DS
% Mức: VD
\dangbt{Vai trò của Vật lí đối với khoa học, kĩ thuật, công nghệ và đời sống}
\begin{ex}
% ID: L10C1B1-34-DS
% Mức: VD
Xét các phát biểu sau về “Vai trò của Vật lí đối với khoa học, kĩ thuật, công nghệ và đời sống” (bộ 1).
\choiceTF
{\True Vật lí cung cấp cơ sở cho nhiều công nghệ hiện đại.}
{Internet không sử dụng thành tựu Vật lí.}
{\True Kiến thức Vật lí góp phần phát triển năng lượng tái tạo.}
{Thiết bị đo không cần nguyên lí Vật lí.}
\loigiai{
\begin{itemchoice}
\item (Đúng) Vật lí cung cấp cơ sở cho nhiều công nghệ hiện đại. (Vì): Vật lí cung cấp cơ sở cho nhiều công nghệ hiện đại. b) S: Internet không sử dụng thành tựu Vật lí. c) Đ: Kiến thức Vật lí góp phần phát triển năng lượng tái tạo. d) S: Thiết bị đo không cần nguyên lí Vật lí. Phân tích được vai trò của Vật lí qua ví dụ cụ thể và đánh giá cả lợi ích lẫn hệ quả của ứng dụng công nghệ
\item (Sai) Internet không sử dụng thành tựu Vật lí.
\item (Đúng) Kiến thức Vật lí góp phần phát triển năng lượng tái tạo.
\item (Sai) Thiết bị đo không cần nguyên lí Vật lí.
\end{itemchoice}
}
\end{ex}

% ===== Câu 178 =====
% ID: L10C1B1-34-TL
% Mức: TH
\dangbt{Đối tượng nghiên cứu và mục tiêu của môn Vật lí}
\begin{bt}
% ID: L10C1B1-34-TL
% Mức: TH
So sánh hai nhận định: “Chuyển động của hành tinh thuộc cấp độ vĩ mô.” và “Sự chuyển động của electron không thuộc Vật lí.”.
\loigiai{
Nhận định thứ nhất đúng: Chuyển động của hành tinh thuộc cấp độ vĩ mô. Nhận định thứ hai sai: Sự chuyển động của electron không thuộc Vật lí. Sự khác nhau nằm ở điều kiện và bản chất kiến thức được áp dụng.
}
\end{bt}

% ===== Câu 179 =====
% ID: L10C1B1-34-TLN
% Mức: NB
\begin{ex}
% ID: L10C1B1-34-TLN
% Mức: NB
Cho bốn nhận định: (1) Mục tiêu của Vật lí là tìm các quy luật tổng quát chi phối vật chất và năng lượng. (2) Vật lí chỉ nghiên cứu cấp độ vĩ mô. (3) Nghiên cứu dòng điện thuộc lĩnh vực Vật lí. (4) Mọi vấn đề sinh học đều chỉ thuộc Vật lí. Có bao nhiêu nhận định đúng? Chỉ ghi số nguyên.
\shortans{2}
\loigiai{
Nhận định (1), (3) đúng; (2), (4) sai. Vậy có 2 nhận định đúng.
}
\end{ex}

% ===== Câu 180 =====
% ID: L10C1B1-34-TN
% Mức: NB
\dangbt{Phương pháp mô hình trong nghiên cứu Vật lí}
\begin{ex}
% ID: L10C1B1-34-TN
% Mức: NB
Trong dạng “Phương pháp mô hình trong nghiên cứu Vật lí”, phát biểu nào sau đây đúng?
\choice
{\True Mỗi mô hình chỉ phù hợp trong phạm vi nhất định.}
{Mô hình không cần kiểm chứng.}
{Không được sửa mô hình khi có dữ liệu mới.}
{Mô hình phải giống hoàn toàn vật thật ở mọi chi tiết.}
\loigiai{
Phát biểu đúng là: Mỗi mô hình chỉ phù hợp trong phạm vi nhất định. Nhận diện mô hình vật chất, mô hình toán học và mô hình lí thuyết; xác định phạm vi áp dụng và kiểm chứng.
}
\end{ex}

% ===== Câu 181 =====
% ID: L10C1B1-35-DS
% Mức: VD
\dangbt{Vai trò của Vật lí đối với khoa học, kĩ thuật, công nghệ và đời sống}
\begin{ex}
% ID: L10C1B1-35-DS
% Mức: VD
Xét các phát biểu sau về “Vai trò của Vật lí đối với khoa học, kĩ thuật, công nghệ và đời sống” (bộ 5).
\choiceTF
{\True Kính thiên văn giúp mở rộng nghiên cứu vũ trụ.}
{Tự động hóa không cần cảm biến.}
{\True Máy chụp X-quang là một ứng dụng của Vật lí.}
{Mọi ứng dụng công nghệ đều không có rủi ro.}
\loigiai{
\begin{itemchoice}
\item (Đúng) Kính thiên văn giúp mở rộng nghiên cứu vũ trụ. (Vì): Kính thiên văn giúp mở rộng nghiên cứu vũ trụ. b) S: Tự động hóa không cần cảm biến. c) Đ: Máy chụp X-quang là một ứng dụng của Vật lí. d) S: Mọi ứng dụng công nghệ đều không có rủi ro. Phân tích được vai trò của Vật lí qua ví dụ cụ thể và đánh giá cả lợi ích lẫn hệ quả của ứng dụng công nghệ
\item (Sai) Tự động hóa không cần cảm biến.
\item (Đúng) Máy chụp X-quang là một ứng dụng của Vật lí.
\item (Sai) Mọi ứng dụng công nghệ đều không có rủi ro.
\end{itemchoice}
}
\end{ex}

% ===== Câu 182 =====
% ID: L10C1B1-35-TL
% Mức: VD
\dangbt{Đối tượng nghiên cứu và mục tiêu của môn Vật lí}
\begin{bt}
% ID: L10C1B1-35-TL
% Mức: VD
Phân tích phát biểu “Mục tiêu của Vật lí là tìm các quy luật tổng quát chi phối vật chất và năng lượng.” và giải thích vì sao phát biểu đó đúng.
\loigiai{
Phát biểu đúng. Mục tiêu của Vật lí là tìm các quy luật tổng quát chi phối vật chất và năng lượng. Nội dung này phù hợp với kiến thức trọng tâm của dạng bài.
}
\end{bt}

% ===== Câu 183 =====
% ID: L10C1B1-35-TLN
% Mức: NB
\begin{ex}
% ID: L10C1B1-35-TLN
% Mức: NB
Cho bốn nhận định: (1) Vật lí là cơ sở của nhiều ngành khoa học tự nhiên. (2) Vật lí không thể dự đoán hiện tượng. (3) Vật lí nghiên cứu các dạng vận động của vật chất và năng lượng. (4) Mục tiêu của Vật lí chỉ là ghi chép hiện tượng. Có bao nhiêu nhận định đúng? Chỉ ghi số nguyên.
\shortans{2}
\loigiai{
Nhận định (1), (3) đúng; (2), (4) sai. Vậy có 2 nhận định đúng.
}
\end{ex}

% ===== Câu 184 =====
% ID: L10C1B1-35-TN
% Mức: NB
\dangbt{Phương pháp mô hình trong nghiên cứu Vật lí}
\begin{ex}
% ID: L10C1B1-35-TN
% Mức: NB
Trong dạng “Phương pháp mô hình trong nghiên cứu Vật lí”, phát biểu nào sau đây đúng?
\choice
{\True Chất điểm là mô hình dùng khi kích thước vật có thể bỏ qua.}
{Mô hình không có khả năng dự đoán.}
{Mọi mô hình đều đúng trong mọi điều kiện.}
{Mô hình không cần kiểm chứng.}
\loigiai{
Phát biểu đúng là: Chất điểm là mô hình dùng khi kích thước vật có thể bỏ qua. Nhận diện mô hình vật chất, mô hình toán học và mô hình lí thuyết; xác định phạm vi áp dụng và kiểm chứng.
}
\end{ex}

% ===== Câu 185 =====
% ID: L10C1B1-36-DS
% Mức: VDC
\dangbt{Vai trò của Vật lí đối với khoa học, kĩ thuật, công nghệ và đời sống}
\begin{ex}
% ID: L10C1B1-36-DS
% Mức: VDC
Xét các phát biểu sau về “Vai trò của Vật lí đối với khoa học, kĩ thuật, công nghệ và đời sống” (bộ 2).
\choiceTF
{\True Thông tin liên lạc phát triển dựa trên điện từ học và quang học.}
{Năng lượng gió không liên quan kiến thức Vật lí.}
{\True Ứng dụng Vật lí có thể vừa mang lợi ích vừa gây tác động môi trường.}
{Bán dẫn là thành tựu không liên quan Vật lí.}
\loigiai{
\begin{itemchoice}
\item (Đúng) Thông tin liên lạc phát triển dựa trên điện từ học và quang học. (Vì): Thông tin liên lạc phát triển dựa trên điện từ học và quang học. b) S: Năng lượng gió không liên quan kiến thức Vật lí. c) Đ: Ứng dụng Vật lí có thể vừa mang lợi ích vừa gây tác động môi trường. d) S: Bán dẫn là thành tựu không liên quan Vật lí. Phân tích được vai trò của Vật lí qua ví dụ cụ thể và đánh giá cả lợi ích lẫn hệ quả của ứng dụng công nghệ
\item (Sai) Năng lượng gió không liên quan kiến thức Vật lí.
\item (Đúng) Ứng dụng Vật lí có thể vừa mang lợi ích vừa gây tác động môi trường.
\item (Sai) Bán dẫn là thành tựu không liên quan Vật lí.
\end{itemchoice}
}
\end{ex}

% ===== Câu 186 =====
% ID: L10C1B1-36-TL
% Mức: VD
\dangbt{Đối tượng nghiên cứu và mục tiêu của môn Vật lí}
\begin{bt}
% ID: L10C1B1-36-TL
% Mức: VD
Xây dựng một quy trình ngắn để giải quyết nhiệm vụ thuộc dạng “Đối tượng nghiên cứu và mục tiêu của môn Vật lí”.
\loigiai{
Quy trình: đọc yêu cầu; xác định khái niệm hoặc đại lượng; chọn nguyên tắc/công thức; xử lí dữ kiện; kiểm tra đơn vị và điều kiện; kết luận. Xác định đúng đối tượng, phạm vi và mục tiêu của Vật lí; phân biệt với đối tượng trung tâm của các môn khoa học khác.
}
\end{bt}

% ===== Câu 187 =====
% ID: L10C1B1-36-TLN
% Mức: TH
\begin{ex}
% ID: L10C1B1-36-TLN
% Mức: TH
Cho bốn nhận định: (1) Vật lí nghiên cứu cả cấp độ vi mô và vĩ mô. (2) Sự chuyển động của electron không thuộc Vật lí. (3) Nghiên cứu ánh sáng thuộc lĩnh vực Vật lí. (4) Vật lí không có ứng dụng trong đời sống. Có bao nhiêu nhận định đúng? Chỉ ghi số nguyên.
\shortans{2}
\loigiai{
Nhận định (1), (3) đúng; (2), (4) sai. Vậy có 2 nhận định đúng.
}
\end{ex}

% ===== Câu 188 =====
% ID: L10C1B1-36-TN
% Mức: TH
\dangbt{Phương pháp mô hình trong nghiên cứu Vật lí}
\begin{ex}
% ID: L10C1B1-36-TN
% Mức: TH
Trong dạng “Phương pháp mô hình trong nghiên cứu Vật lí”, phát biểu nào sau đây đúng?
\choice
{Mọi mô hình đều đúng trong mọi điều kiện.}
{\True Mô hình vật chất có thể là mẫu thu nhỏ.}
{Mô hình không cần kiểm chứng.}
{Không được sửa mô hình khi có dữ liệu mới.}
\loigiai{
Phát biểu đúng là: Mô hình vật chất có thể là mẫu thu nhỏ. Nhận diện mô hình vật chất, mô hình toán học và mô hình lí thuyết; xác định phạm vi áp dụng và kiểm chứng.
}
\end{ex}

% ===== Câu 189 =====
% ID: L10C1B1-37-DS
% Mức: NB
\dangbt{Vận dụng tiến trình của phương pháp mô hình}
\begin{ex}
% ID: L10C1B1-37-DS
% Mức: NB
Khi tìm hiểu về sơ đồ các bước thực hiện của phương pháp mô hình, hãy chọn tính đúng sai của các nhận định dưới đây:
\choiceTF
{\True Bước đầu tiên của tiến trình phương pháp mô hình là xác định đối tượng cần được mô hình hóa.}
{Quy trình xây dựng mô hình luôn diễn ra một chiều thẳng từ bước một đến bước bốn mà không bao giờ có sự quay đầu điều chỉnh.}
{\True Mô hình toán học của chuyển động thẳng đều $s = vt$ cần được kiểm tra sự phù hợp với các kết quả thực tế đo đạc.}
{Khi kiểm tra sự phù hợp của mô hình mà thấy chưa phù hợp, ta có thể bỏ qua sai số thực tế để đưa ra kết luận ngay.}
\loigiai{
\begin{itemchoice}
\item (Đúng) Bước đầu tiên của tiến trình phương pháp mô hình là xác định đối tượng cần được mô hình hóa. (Vì): Mệnh đề một đúng vì đây là bước khởi đầu để định hình đối tượng cần nghiên cứu
\item (Sai) Quy trình xây dựng mô hình luôn diễn ra một chiều thẳng từ bước một đến bước bốn mà không bao giờ có sự quay đầu điều chỉnh. (Vì): Mệnh đề hai sai vì khi kiểm tra thấy mô hình chưa phù hợp, ta cần phải thực hiện bước phụ quay lại để điều chỉnh mô hình
\item (Đúng) Mô hình toán học của chuyển động thẳng đều $s = vt$ cần được kiểm tra sự phù hợp với các kết quả thực tế đo đạc. (Vì): Mệnh đề ba đúng vì mọi mô hình toán học trong Vật lí đều cần vượt qua phép kiểm tra thực tế để khẳng định tính chính xác
\item (Sai) Khi kiểm tra sự phù hợp của mô hình mà thấy chưa phù hợp, ta có thể bỏ qua sai số thực tế để đưa ra kết luận ngay. (Vì): Mệnh đề bốn sai vì nhà nghiên cứu không được bỏ qua sai số thực tế mà bắt buộc phải điều chỉnh lại mô hình để đảm bảo tính khoa học
\end{itemchoice}
}
\end{ex}

% ===== Câu 190 =====
% ID: L10C1B1-37-TL
% Mức: VDC
\dangbt{Đối tượng nghiên cứu và mục tiêu của môn Vật lí}
\begin{bt}
% ID: L10C1B1-37-TL
% Mức: VDC
Phân tích phát biểu “Mục tiêu của Vật lí chỉ là ghi chép hiện tượng.”, chỉ ra sai sót và sửa lại.
\loigiai{
Phát biểu sai: Mục tiêu của Vật lí chỉ là ghi chép hiện tượng. Cần sửa thành: Vật lí nghiên cứu cả cấp độ vi mô và vĩ mô. Xác định đúng đối tượng, phạm vi và mục tiêu của Vật lí; phân biệt với đối tượng trung tâm của các môn khoa học khác.
}
\end{bt}

% ===== Câu 191 =====
% ID: L10C1B1-37-TLN
% Mức: VD
\begin{ex}
% ID: L10C1B1-37-TLN
% Mức: VD
Cho bốn nhận định: (1) Electron thuộc đối tượng ở cấp độ vi mô. (2) Vật lí hoàn toàn tách rời công nghệ. (3) Nghiên cứu sự truyền nhiệt thuộc lĩnh vực Vật lí. (4) Vật lí chỉ nghiên cứu vật thể nhìn thấy bằng mắt. Có bao nhiêu nhận định đúng? Chỉ ghi số nguyên.
\shortans{2}
\loigiai{
Nhận định (1), (3) đúng; (2), (4) sai. Vậy có 2 nhận định đúng.
}
\end{ex}

% ===== Câu 192 =====
% ID: L10C1B1-37-TN
% Mức: TH
\dangbt{Phương pháp mô hình trong nghiên cứu Vật lí}
\begin{ex}
% ID: L10C1B1-37-TN
% Mức: TH
Trong dạng “Phương pháp mô hình trong nghiên cứu Vật lí”, phát biểu nào sau đây đúng?
\choice
{Sơ đồ không thể là mô hình.}
{\True Kết quả từ mô hình cần được đối chiếu thực nghiệm.}
{Mô hình không có khả năng dự đoán.}
{Mọi mô hình đều đúng trong mọi điều kiện.}
\loigiai{
Phát biểu đúng là: Kết quả từ mô hình cần được đối chiếu thực nghiệm. Nhận diện mô hình vật chất, mô hình toán học và mô hình lí thuyết; xác định phạm vi áp dụng và kiểm chứng.
}
\end{ex}

% ===== Câu 193 =====
% ID: L10C1B1-38-DS
% Mức: NB
\dangbt{Vận dụng tiến trình của phương pháp mô hình}
\begin{ex}
% ID: L10C1B1-38-DS
% Mức: NB
Hoạt động nghiên cứu khoa học:
\choiceTF
{Trồng hoa trong nhà kính}
{\True Tìm vaccine trong phòng thí nghiệm}
{Sản xuất muối ăn}
{\True Thực hiện thí nghiệm vật liệu mới}
\loigiai{
\begin{itemchoice}
\item (Sai) Trồng hoa trong nhà kính. (Vì): Phát biểu này sai. Là nông nghiệp ứng dụng
\item (Đúng) Tìm vaccine trong phòng thí nghiệm. (Vì): Phát biểu này đúng. Là nghiên cứu bản chất virus
\item (Sai) Sản xuất muối ăn. (Vì): Phát biểu này sai. Là sản xuất truyền thống
\item (Đúng) Thực hiện thí nghiệm vật liệu mới. (Vì): Phát biểu này đúng. Là thực nghiệm khoa học
\end{itemchoice}
}
\end{ex}

% ===== Câu 194 =====
% ID: L10C1B1-38-TLN
% Mức: VDC
\dangbt{Đối tượng nghiên cứu và mục tiêu của môn Vật lí}
\begin{ex}
% ID: L10C1B1-38-TLN
% Mức: VDC
Cho bốn nhận định: (1) Vật lí nghiên cứu các dạng vận động của vật chất và năng lượng. (2) Mục tiêu của Vật lí chỉ là ghi chép hiện tượng. (3) Vật lí là cơ sở của nhiều ngành khoa học tự nhiên. (4) Vật lí không thể dự đoán hiện tượng. Có bao nhiêu nhận định đúng? Chỉ ghi số nguyên.
\shortans{2}
\loigiai{
Nhận định (1), (3) đúng; (2), (4) sai. Vậy có 2 nhận định đúng.
}
\end{ex}

% ===== Câu 195 =====
% ID: L10C1B1-38-TN
% Mức: TH
\dangbt{Phương pháp mô hình trong nghiên cứu Vật lí}
\begin{ex}
% ID: L10C1B1-38-TN
% Mức: TH
Trong dạng “Phương pháp mô hình trong nghiên cứu Vật lí”, phát biểu nào sau đây đúng?
\choice
{Một đối tượng chỉ có duy nhất một mô hình.}
{\True Mô hình giúp tập trung vào đặc trưng quan trọng.}
{Mô hình toán học không dùng phương trình.}
{Sơ đồ không thể là mô hình.}
\loigiai{
Phát biểu đúng là: Mô hình giúp tập trung vào đặc trưng quan trọng. Nhận diện mô hình vật chất, mô hình toán học và mô hình lí thuyết; xác định phạm vi áp dụng và kiểm chứng.
}
\end{ex}

% ===== Câu 196 =====
% ID: L10C1B1-39-DS
% Mức: NB
\dangbt{Vận dụng tiến trình của phương pháp mô hình}
\begin{ex}
% ID: L10C1B1-39-DS
% Mức: NB
Về cấp độ nghiên cứu trong Vật lí:
\choiceTF
{\True Vi mô mô phỏng hệ như nguyên tử, phân tử}
{Vĩ mô mô tả các hạt cơ bản}
{\True Vĩ mô là khảo sát hệ như xe, tàu, hành tinh}
{Vi mô mô tả hiện tượng lớn như bão từ}
\loigiai{
\begin{itemchoice}
\item (Đúng) Vi mô mô phỏng hệ như nguyên tử, phân tử. (Vì): Phát biểu này đúng. Là cấp độ vi mô
\item (Sai) Vĩ mô mô tả các hạt cơ bản. (Vì): Phát biểu này sai. Vĩ mô không mô tả hạt cơ bản
\item (Đúng) Vĩ mô là khảo sát hệ như xe, tàu, hành tinh. (Vì): Phát biểu này đúng. Đúng định nghĩa vĩ mô
\item (Sai) Vi mô mô tả hiện tượng lớn như bão từ. (Vì): Phát biểu này sai. Vi mô không mô tả hiện tượng lớn
\end{itemchoice}
}
\end{ex}

% ===== Câu 197 =====
% ID: L10C1B1-39-TLN
% Mức: VDC
\dangbt{Đối tượng nghiên cứu và mục tiêu của môn Vật lí}
\begin{ex}
% ID: L10C1B1-39-TLN
% Mức: VDC
Cho bốn nhận định: (1) Chuyển động của hành tinh thuộc cấp độ vĩ mô. (2) Phản ứng tạo chất mới là đối tượng trung tâm duy nhất của Vật lí. (3) Kiến thức Vật lí có thể dùng để giải quyết vấn đề thực tiễn. (4) Vật lí không nghiên cứu năng lượng. Có bao nhiêu nhận định đúng? Chỉ ghi số nguyên.
\shortans{2}
\loigiai{
Nhận định (1), (3) đúng; (2), (4) sai. Vậy có 2 nhận định đúng.
}
\end{ex}

% ===== Câu 198 =====
% ID: L10C1B1-39-TN
% Mức: VD
\dangbt{Phương pháp mô hình trong nghiên cứu Vật lí}
\begin{ex}
% ID: L10C1B1-39-TN
% Mức: VD
Trong dạng “Phương pháp mô hình trong nghiên cứu Vật lí”, phát biểu nào sau đây đúng?
\choice
{Mô hình toán học không dùng phương trình.}
{Sơ đồ không thể là mô hình.}
{\True Mô hình toán học dùng phương trình mô tả quan hệ giữa các đại lượng.}
{Mô hình không có khả năng dự đoán.}
\loigiai{
Phát biểu đúng là: Mô hình toán học dùng phương trình mô tả quan hệ giữa các đại lượng. Nhận diện mô hình vật chất, mô hình toán học và mô hình lí thuyết; xác định phạm vi áp dụng và kiểm chứng.
}
\end{ex}

% ===== Câu 199 =====
% ID: L10C1B1-40-DS
% Mức: NB
\dangbt{Vận dụng tiến trình của phương pháp mô hình}
\begin{ex}
% ID: L10C1B1-40-DS
% Mức: NB
Nhận xét các nhận định sau về đối tượng nghiên cứu và mục tiêu của môn Vật lí.
\choiceTF
{\True Đối tượng nghiên cứu của Vật lí là các dạng vật chất và năng lượng}
{\True Việc học tập Vật lí giúp người học hình thành kiến thức cơ bản về khoa học và vận dụng để khám phá, giải quyết các vấn đề có liên quan trong đời sống}
{Quá trình hình thành chất mới là đối tượng nghiên cứu của môn Vật lí}
{Mục tiêu nghiên cứu và học tập môn Vật lí giúp ta giải thích được sự sinh trưởng và phát triển của sinh vật (con người, động vật, thực vật, …)}
\loigiai{
\begin{itemchoice}
\item (Đúng) Đối tượng nghiên cứu của Vật lí là các dạng vật chất và năng lượng. (Vì): đối tượng nghiên cứu của Vật lí là các dạng vật chất và năng lượng
\item (Đúng) Việc học tập Vật lí giúp người học hình thành kiến thức cơ bản về khoa học và vận dụng để khám phá, giải quyết các vấn đề có liên quan trong đời sống. (Vì): mục tiêu học tập môn Vật lí giúp học sinh có được những kiến thức, kĩ năng cơ bản về Vật lí; Vận dụng được những kiến thức, kĩ năng đã học để khám phá, giải quyết các vấn đề có liên quan trong học tập cũng như trong đời sống; Nhận biết được năng lực, sở trường của bản thân, định hướng nghề nghiệp
\item (Sai) Quá trình hình thành chất mới là đối tượng nghiên cứu của môn Vật lí. (Vì): đối tượng nghiên cứu của Vật lí là các dạng vật chất và năng lượng và không bao gồm sự hình thành chất mới sau quá trình phản ứng hay sự phát triển và sinh trưởng của sinh vật
\item (Sai) Mục tiêu nghiên cứu và học tập môn Vật lí giúp ta giải thích được sự sinh trưởng và phát triển của sinh vật (con người, động vật, thực vật, …). (Vì): Phát biểu này sai. Mục tiêu nghiên cứu và học tập môn Vật lí giúp ta giải thích được sự sinh trưởng và phát triển của sinh vật
\end{itemchoice}
}
\end{ex}

% ===== Câu 200 =====
% ID: L10C1B1-40-TN
% Mức: VD
\dangbt{Phương pháp mô hình trong nghiên cứu Vật lí}
\begin{ex}
% ID: L10C1B1-40-TN
% Mức: VD
Trong dạng “Phương pháp mô hình trong nghiên cứu Vật lí”, phát biểu nào sau đây đúng?
\choice
{Mô hình càng nhiều chi tiết thừa càng tốt.}
{Một đối tượng chỉ có duy nhất một mô hình.}
{\True Có thể điều chỉnh mô hình khi bằng chứng mới xuất hiện.}
{Mô hình toán học không dùng phương trình.}
\loigiai{
Phát biểu đúng là: Có thể điều chỉnh mô hình khi bằng chứng mới xuất hiện. Nhận diện mô hình vật chất, mô hình toán học và mô hình lí thuyết; xác định phạm vi áp dụng và kiểm chứng.
}
\end{ex}

% ===== Câu 201 =====
% ID: L10C1B1-41-DS
% Mức: NB
\dangbt{Vận dụng tiến trình của phương pháp mô hình}
\begin{ex}
% ID: L10C1B1-41-DS
% Mức: NB
Trong quá trình nghiên cứu và phát triển kiến thức vật lí, các yếu tố sau đây đóng vai trò như thế nào?
\choiceTF
{Mô hình vật lí luôn là bản sao chính xác và đầy đủ của hiện thực, không bỏ qua bất kỳ chi tiết nào}
{\True Toán học là công cụ ngôn ngữ chính yếu để mô tả các định luật vật lí và biểu diễn các mối quan hệ giữa các đại lượng}
{Các định luật vật lí chủ yếu được xây dựng dựa trên suy luận logic từ các lí thuyết đã có mà không cần đến thực nghiệm}
{\True Mục tiêu của vật lí là tìm kiếm các quy luật chung nhất chi phối mọi hiện tượng tự nhiên, từ vi mô đến vĩ mô}
\loigiai{
\begin{itemchoice}
\item (Sai) Mô hình vật lí luôn là bản sao chính xác và đầy đủ của hiện thực, không bỏ qua bất kỳ chi tiết nào. (Vì): Phát biểu này sai. Mô hình vật lí là sự đơn giản hóa của hiện thực, chỉ giữ lại những yếu tố quan trọng nhất để dễ dàng nghiên cứu và dự đoán
\item (Đúng) Toán học là công cụ ngôn ngữ chính yếu để mô tả các định luật vật lí và biểu diễn các mối quan hệ giữa các đại lượng. (Vì): Phát biểu này đúng. Toán học cung cấp khung ngôn ngữ chặt chẽ và chính xác để biểu diễn các khái niệm, định luật và mối quan hệ trong vật lí
\item (Sai) Các định luật vật lí chủ yếu được xây dựng dựa trên suy luận logic từ các lí thuyết đã có mà không cần đến thực nghiệm. (Vì): Phát biểu này sai. Các định luật vật lí được xây dựng dựa trên sự tổng hợp từ nhiều quan sát và thực nghiệm, sau đó mới được diễn đạt bằng ngôn ngữ toán học và có thể được giải thích bởi các lí thuyết
\item (Đúng) Mục tiêu của vật lí là tìm kiếm các quy luật chung nhất chi phối mọi hiện tượng tự nhiên, từ vi mô đến vĩ mô. (Vì): Phát biểu này đúng. Vật lí học là khoa học cơ bản nghiên cứu các quy luật vận động của vật chất và năng lượng, tìm kiếm các nguyên lí và định luật phổ quát áp dụng cho toàn bộ vũ trụ
\end{itemchoice}
}
\end{ex}

% ===== Câu 202 =====
% ID: L10C1B1-41-TN
% Mức: VDC
\dangbt{Phương pháp mô hình trong nghiên cứu Vật lí}
\begin{ex}
% ID: L10C1B1-41-TN
% Mức: VDC
Trong dạng “Phương pháp mô hình trong nghiên cứu Vật lí”, phát biểu nào sau đây đúng?
\choice
{Chất điểm luôn là vật có kích thước bằng không trong thực tế.}
{Mô hình càng nhiều chi tiết thừa càng tốt.}
{Một đối tượng chỉ có duy nhất một mô hình.}
{\True Mô hình lí thuyết giúp giải thích và dự đoán.}
\loigiai{
Phát biểu đúng là: Mô hình lí thuyết giúp giải thích và dự đoán. Nhận diện mô hình vật chất, mô hình toán học và mô hình lí thuyết; xác định phạm vi áp dụng và kiểm chứng.
}
\end{ex}

% ===== Câu 203 =====
% ID: L10C1B1-42-DS
% Mức: NB
\dangbt{Vận dụng tiến trình nghiên cứu Vật lí vào tình huống thực tiễn}
\begin{ex}
% ID: L10C1B1-42-DS
% Mức: NB
Xét các phát biểu sau về “Vận dụng tiến trình nghiên cứu Vật lí vào tình huống thực tiễn” (bộ 3).
\choiceTF
{\True Đồ thị có thể làm rõ quan hệ giữa hai đại lượng.}
{Lặp đo luôn vô ích.}
{\True Lặp phép đo giúp đánh giá độ ổn định.}
{Kế hoạch không cần dụng cụ.}
\loigiai{
\begin{itemchoice}
\item (Đúng) Đồ thị có thể làm rõ quan hệ giữa hai đại lượng. (Vì): ồ thị có thể làm rõ quan hệ giữa hai đại lượng. b) S: Lặp đo luôn vô ích. c) Đ: Lặp phép đo giúp đánh giá độ ổn định. d) S: Kế hoạch không cần dụng cụ. Vận dụng đầy đủ chu trình nghiên cứu vào vấn đề gần gũi, từ câu hỏi đến đánh giá và cải tiến
\item (Sai) Lặp đo luôn vô ích.
\item (Đúng) Lặp phép đo giúp đánh giá độ ổn định.
\item (Sai) Kế hoạch không cần dụng cụ.
\end{itemchoice}
}
\end{ex}

% ===== Câu 204 =====
% ID: L10C1B1-42-TN
% Mức: VDC
\dangbt{Phương pháp mô hình trong nghiên cứu Vật lí}
\begin{ex}
% ID: L10C1B1-42-TN
% Mức: VDC
Trong dạng “Phương pháp mô hình trong nghiên cứu Vật lí”, phát biểu nào sau đây đúng?
\choice
{Không được sửa mô hình khi có dữ liệu mới.}
{Mô hình phải giống hoàn toàn vật thật ở mọi chi tiết.}
{Chất điểm luôn là vật có kích thước bằng không trong thực tế.}
{\True Sơ đồ tia sáng là một dạng mô hình.}
\loigiai{
Phát biểu đúng là: Sơ đồ tia sáng là một dạng mô hình. Nhận diện mô hình vật chất, mô hình toán học và mô hình lí thuyết; xác định phạm vi áp dụng và kiểm chứng.
}
\end{ex}

% ===== Câu 205 =====
% ID: L10C1B1-43-DS
% Mức: TH
\dangbt{Vận dụng tiến trình nghiên cứu Vật lí vào tình huống thực tiễn}
\begin{ex}
% ID: L10C1B1-43-DS
% Mức: TH
Xét các phát biểu sau về “Vận dụng tiến trình nghiên cứu Vật lí vào tình huống thực tiễn” (bộ 4).
\choiceTF
{\True Cần nêu nguồn sai số và giới hạn của kết quả.}
{Có thể chọn bỏ dữ liệu tùy ý.}
{\True Có thể đề xuất cải tiến sau khi đánh giá kết quả.}
{Câu hỏi càng mơ hồ càng dễ nghiên cứu.}
\loigiai{
\begin{itemchoice}
\item (Đúng) Cần nêu nguồn sai số và giới hạn của kết quả. (Vì): Cần nêu nguồn sai số và giới hạn của kết quả. b) S: Có thể chọn bỏ dữ liệu tùy ý. c) Đ: Có thể đề xuất cải tiến sau khi đánh giá kết quả. d) S: Câu hỏi càng mơ hồ càng dễ nghiên cứu. Vận dụng đầy đủ chu trình nghiên cứu vào vấn đề gần gũi, từ câu hỏi đến đánh giá và cải tiến
\item (Sai) Có thể chọn bỏ dữ liệu tùy ý.
\item (Đúng) Có thể đề xuất cải tiến sau khi đánh giá kết quả.
\item (Sai) Câu hỏi càng mơ hồ càng dễ nghiên cứu.
\end{itemchoice}
}
\end{ex}

% ===== Câu 206 =====
% ID: L10C1B1-43-TN
% Mức: NB
\dangbt{Phương pháp thực nghiệm trong nghiên cứu Vật lí}
\begin{ex}
% ID: L10C1B1-43-TN
% Mức: NB
Trong dạng “Phương pháp thực nghiệm trong nghiên cứu Vật lí”, phát biểu nào sau đây đúng?
\choice
{\True Phương pháp thực nghiệm sử dụng thí nghiệm để kiểm tra dự đoán.}
{Thực nghiệm chỉ cần kết luận, không cần dữ liệu.}
{Không cần xác định vấn đề trước khi đo.}
{Một phép đo duy nhất luôn tuyệt đối chính xác.}
\loigiai{
Phát biểu đúng là: Phương pháp thực nghiệm sử dụng thí nghiệm để kiểm tra dự đoán. Tiến trình cơ bản: xác định vấn đề, dự đoán, thiết kế và thực hiện thí nghiệm, xử lí dữ liệu, kết luận và điều chỉnh.
}
\end{ex}

% ===== Câu 207 =====
% ID: L10C1B1-44-DS
% Mức: VD
\dangbt{Vận dụng tiến trình nghiên cứu Vật lí vào tình huống thực tiễn}
\begin{ex}
% ID: L10C1B1-44-DS
% Mức: VD
Xét các phát biểu sau về “Vận dụng tiến trình nghiên cứu Vật lí vào tình huống thực tiễn” (bộ 1).
\choiceTF
{\True Câu hỏi nghiên cứu cần rõ ràng và có thể kiểm tra.}
{Kết luận có thể trái dữ liệu mà không cần giải thích.}
{\True Bảng số liệu giúp tổ chức kết quả đo.}
{Đồ thị không thể hiện quan hệ đại lượng.}
\loigiai{
\begin{itemchoice}
\item (Đúng) Câu hỏi nghiên cứu cần rõ ràng và có thể kiểm tra. (Vì): Câu hỏi nghiên cứu cần rõ ràng và có thể kiểm tra. b) S: Kết luận có thể trái dữ liệu mà không cần giải thích. c) Đ: Bảng số liệu giúp tổ chức kết quả đo. d) S: Đồ thị không thể hiện quan hệ đại lượng. Vận dụng đầy đủ chu trình nghiên cứu vào vấn đề gần gũi, từ câu hỏi đến đánh giá và cải tiến
\item (Sai) Kết luận có thể trái dữ liệu mà không cần giải thích.
\item (Đúng) Bảng số liệu giúp tổ chức kết quả đo.
\item (Sai) Đồ thị không thể hiện quan hệ đại lượng.
\end{itemchoice}
}
\end{ex}

% ===== Câu 208 =====
% ID: L10C1B1-44-TN
% Mức: NB
\dangbt{Phương pháp thực nghiệm trong nghiên cứu Vật lí}
\begin{ex}
% ID: L10C1B1-44-TN
% Mức: NB
Trong dạng “Phương pháp thực nghiệm trong nghiên cứu Vật lí”, phát biểu nào sau đây đúng?
\choice
{\True Kết quả đo cần được xử lí trước khi kết luận.}
{Biến phụ thuộc luôn do người làm tùy ý đặt giá trị.}
{Không cần ghi chép số liệu.}
{Thực nghiệm chỉ cần kết luận, không cần dữ liệu.}
\loigiai{
Phát biểu đúng là: Kết quả đo cần được xử lí trước khi kết luận. Tiến trình cơ bản: xác định vấn đề, dự đoán, thiết kế và thực hiện thí nghiệm, xử lí dữ liệu, kết luận và điều chỉnh.
}
\end{ex}

% ===== Câu 209 =====
% ID: L10C1B1-45-DS
% Mức: VD
\dangbt{Vận dụng tiến trình nghiên cứu Vật lí vào tình huống thực tiễn}
\begin{ex}
% ID: L10C1B1-45-DS
% Mức: VD
Xét các phát biểu sau về “Vận dụng tiến trình nghiên cứu Vật lí vào tình huống thực tiễn” (bộ 5).
\choiceTF
{\True Điều kiện an toàn phải được đưa vào kế hoạch.}
{Không được cải tiến phương án.}
{\True Dự đoán nên dựa trên kiến thức hoặc quan sát ban đầu.}
{Không cần quan tâm an toàn.}
\loigiai{
\begin{itemchoice}
\item (Đúng) Điều kiện an toàn phải được đưa vào kế hoạch. (Vì): iều kiện an toàn phải được đưa vào kế hoạch. b) S: Không được cải tiến phương án. c) Đ: Dự đoán nên dựa trên kiến thức hoặc quan sát ban đầu. d) S: Không cần quan tâm an toàn. Vận dụng đầy đủ chu trình nghiên cứu vào vấn đề gần gũi, từ câu hỏi đến đánh giá và cải tiến
\item (Sai) Không được cải tiến phương án.
\item (Đúng) Dự đoán nên dựa trên kiến thức hoặc quan sát ban đầu.
\item (Sai) Không cần quan tâm an toàn.
\end{itemchoice}
}
\end{ex}

% ===== Câu 210 =====
% ID: L10C1B1-45-TN
% Mức: NB
\dangbt{Phương pháp thực nghiệm trong nghiên cứu Vật lí}
\begin{ex}
% ID: L10C1B1-45-TN
% Mức: NB
Trong dạng “Phương pháp thực nghiệm trong nghiên cứu Vật lí”, phát biểu nào sau đây đúng?
\choice
{\True Các yếu tố khác nên được kiểm soát.}
{Giả thuyết không cần kiểm tra.}
{Có thể thay đổi đồng thời mọi yếu tố mà vẫn xác định chắc chắn nguyên nhân.}
{Biến phụ thuộc luôn do người làm tùy ý đặt giá trị.}
\loigiai{
Phát biểu đúng là: Các yếu tố khác nên được kiểm soát. Tiến trình cơ bản: xác định vấn đề, dự đoán, thiết kế và thực hiện thí nghiệm, xử lí dữ liệu, kết luận và điều chỉnh.
}
\end{ex}

% ===== Câu 211 =====
% ID: L10C1B1-46-DS
% Mức: VDC
\dangbt{Vận dụng tiến trình nghiên cứu Vật lí vào tình huống thực tiễn}
\begin{ex}
% ID: L10C1B1-46-DS
% Mức: VDC
Xét các phát biểu sau về “Vận dụng tiến trình nghiên cứu Vật lí vào tình huống thực tiễn” (bộ 2).
\choiceTF
{\True Kế hoạch cần nêu dụng cụ và cách thu thập dữ liệu.}
{Không cần ghi đơn vị trong bảng số liệu.}
{\True Kết luận phải dựa trên dữ liệu.}
{Sai số không cần được xem xét.}
\loigiai{
\begin{itemchoice}
\item (Đúng) Kế hoạch cần nêu dụng cụ và cách thu thập dữ liệu. (Vì): Kế hoạch cần nêu dụng cụ và cách thu thập dữ liệu. b) S: Không cần ghi đơn vị trong bảng số liệu. c) Đ: Kết luận phải dựa trên dữ liệu. d) S: Sai số không cần được xem xét. Vận dụng đầy đủ chu trình nghiên cứu vào vấn đề gần gũi, từ câu hỏi đến đánh giá và cải tiến
\item (Sai) Không cần ghi đơn vị trong bảng số liệu.
\item (Đúng) Kết luận phải dựa trên dữ liệu.
\item (Sai) Sai số không cần được xem xét.
\end{itemchoice}
}
\end{ex}

% ===== Câu 212 =====
% ID: L10C1B1-46-TN
% Mức: TH
\dangbt{Phương pháp thực nghiệm trong nghiên cứu Vật lí}
\begin{ex}
% ID: L10C1B1-46-TN
% Mức: TH
Trong dạng “Phương pháp thực nghiệm trong nghiên cứu Vật lí”, phát biểu nào sau đây đúng?
\choice
{Có thể thay đổi đồng thời mọi yếu tố mà vẫn xác định chắc chắn nguyên nhân.}
{\True Bước đầu cần xác định vấn đề nghiên cứu.}
{Biến phụ thuộc luôn do người làm tùy ý đặt giá trị.}
{Không cần ghi chép số liệu.}
\loigiai{
Phát biểu đúng là: Bước đầu cần xác định vấn đề nghiên cứu. Tiến trình cơ bản: xác định vấn đề, dự đoán, thiết kế và thực hiện thí nghiệm, xử lí dữ liệu, kết luận và điều chỉnh.
}
\end{ex}

% ===== Câu 213 =====
% ID: L10C1B1-47-DS
% Mức: NB
\dangbt{Xác định đối tượng nghiên cứu và mục tiêu của môn Vật lí}
\begin{ex}
% ID: L10C1B1-47-DS
% Mức: NB
Đối tượng nghiên cứu của Vật lí bao gồm:
\choiceTF
{Vật chất và năng lượng}
{Các hiện tượng tự nhiên}
{\True Các dạng vận động của vật chất và năng lượng}
{\True Các chuyển động cơ học và năng lượng}
\loigiai{
\begin{itemchoice}
\item (Sai) Vật chất và năng lượng. (Vì): Chỉ nêu đối tượng, chưa đề cập đến dạng vận động
\item (Sai) Các hiện tượng tự nhiên. (Vì): Không đầy đủ và không bản chất
\item (Đúng) Các dạng vận động của vật chất và năng lượng. (Vì): ây là định nghĩa chính xác
\item (Đúng) Các chuyển động cơ học và năng lượng. (Vì): Là biểu hiện cụ thể của vận động
\end{itemchoice}
}
\end{ex}

% ===== Câu 214 =====
% ID: L10C1B1-47-TN
% Mức: TH
\dangbt{Phương pháp thực nghiệm trong nghiên cứu Vật lí}
\begin{ex}
% ID: L10C1B1-47-TN
% Mức: TH
Trong dạng “Phương pháp thực nghiệm trong nghiên cứu Vật lí”, phát biểu nào sau đây đúng?
\choice
{Không cần kiểm soát điều kiện.}
{\True Nếu kết quả không phù hợp, cần xem lại giả thuyết hoặc phương án.}
{Giả thuyết không cần kiểm tra.}
{Có thể thay đổi đồng thời mọi yếu tố mà vẫn xác định chắc chắn nguyên nhân.}
\loigiai{
Phát biểu đúng là: Nếu kết quả không phù hợp, cần xem lại giả thuyết hoặc phương án. Tiến trình cơ bản: xác định vấn đề, dự đoán, thiết kế và thực hiện thí nghiệm, xử lí dữ liệu, kết luận và điều chỉnh.
}
\end{ex}

% ===== Câu 215 =====
% ID: L10C1B1-48-DS
% Mức: NB
\dangbt{Xác định đối tượng nghiên cứu và mục tiêu của môn Vật lí}
\begin{ex}
% ID: L10C1B1-48-DS
% Mức: NB
Khi nói về mục tiêu và đối tượng nghiên cứu của Vật lí, hãy chọn tính đúng/sai cho các phát biểu sau.
\choiceTF
{\True Đối tượng nghiên cứu chủ yếu của Vật lí bao gồm các dạng của vật chất và năng lượng.}
{Mục tiêu cốt lõi của Vật lí chỉ dừng lại ở việc mô tả các hiện tượng tự nhiên ở cấp độ vĩ mô mà không giải thích cấp độ vi mô.}
{\True Nghiên cứu tương tác giữa vật chất và năng lượng ở mọi cấp độ là một trong những mục tiêu quan trọng của Vật lí.}
{Môn Vật lí độc lập hoàn toàn và không có mối quan hệ liên môn nào với các ngành khoa học tự nhiên khác.}
\loigiai{
\begin{itemchoice}
\item (Đúng) Đối tượng nghiên cứu chủ yếu của Vật lí bao gồm các dạng của vật chất và năng lượng. (Vì): Mệnh đề một đúng vì đây là định nghĩa về đối tượng nghiên cứu chủ yếu của Vật lí học
\item (Sai) Mục tiêu cốt lõi của Vật lí chỉ dừng lại ở việc mô tả các hiện tượng tự nhiên ở cấp độ vĩ mô mà không giải thích cấp độ vi mô. (Vì): Mệnh đề hai sai vì mục tiêu của Vật lí là khám phá quy luật tổng quát nhất chi phối sự vận động của vật chất và năng lượng ở cả cấp độ vi mô và vĩ mô
\item (Đúng) Nghiên cứu tương tác giữa vật chất và năng lượng ở mọi cấp độ là một trong những mục tiêu quan trọng của Vật lí. (Vì): Mệnh đề ba đúng vì mục tiêu của Vật lí bao gồm nghiên cứu tương tác giữa vật chất và năng lượng ở mọi cấp độ
\item (Sai) Môn Vật lí độc lập hoàn toàn và không có mối quan hệ liên môn nào với các ngành khoa học tự nhiên khác. (Vì): Mệnh đề bốn sai vì Vật lí có quan hệ mật thiết với mọi ngành khoa học tự nhiên và là cơ sở của khoa học tự nhiên
\end{itemchoice}
}
\end{ex}

% ===== Câu 216 =====
% ID: L10C1B1-48-TN
% Mức: TH
\dangbt{Phương pháp thực nghiệm trong nghiên cứu Vật lí}
\begin{ex}
% ID: L10C1B1-48-TN
% Mức: TH
Trong dạng “Phương pháp thực nghiệm trong nghiên cứu Vật lí”, phát biểu nào sau đây đúng?
\choice
{Thí nghiệm không cần lặp lại.}
{\True Thí nghiệm cần có khả năng lặp lại.}
{Kết quả trái dự đoán phải bị bỏ đi.}
{Không cần kiểm soát điều kiện.}
\loigiai{
Phát biểu đúng là: Thí nghiệm cần có khả năng lặp lại. Tiến trình cơ bản: xác định vấn đề, dự đoán, thiết kế và thực hiện thí nghiệm, xử lí dữ liệu, kết luận và điều chỉnh.
}
\end{ex}

% ===== Câu 217 =====
% ID: L10C1B1-49-DS
% Mức: NB
\dangbt{Xác định đối tượng nghiên cứu và mục tiêu của môn Vật lí}
\begin{ex}
% ID: L10C1B1-49-DS
% Mức: NB
Mục tiêu của môn Vật lí gồm:
\choiceTF
{\True Khám phá các quy luật vận động của vật chất và năng lượng}
{Tìm ra phương pháp sản xuất hiện đại}
{\True Tìm hiểu sự tương tác giữa vật chất và năng lượng ở mọi cấp độ}
{Tập trung khảo sát hiện tượng bằng kinh nghiệm}
\loigiai{
\begin{itemchoice}
\item (Đúng) Khám phá các quy luật vận động của vật chất và năng lượng. (Vì): Là mục tiêu cốt lõi
\item (Sai) Tìm ra phương pháp sản xuất hiện đại. (Vì): Không phải mục tiêu của Vật lí
\item (Đúng) Tìm hiểu sự tương tác giữa vật chất và năng lượng ở mọi cấp độ. (Vì): Phản ánh đầy đủ phạm vi nghiên cứu
\item (Sai) Tập trung khảo sát hiện tượng bằng kinh nghiệm. (Vì): Không phải là phương pháp khoa học
\end{itemchoice}
}
\end{ex}

% ===== Câu 218 =====
% ID: L10C1B1-49-TN
% Mức: VD
\dangbt{Phương pháp thực nghiệm trong nghiên cứu Vật lí}
\begin{ex}
% ID: L10C1B1-49-TN
% Mức: VD
Trong dạng “Phương pháp thực nghiệm trong nghiên cứu Vật lí”, phát biểu nào sau đây đúng?
\choice
{Kết quả trái dự đoán phải bị bỏ đi.}
{Không cần kiểm soát điều kiện.}
{\True Sau khi nêu vấn đề có thể đề xuất giả thuyết hoặc dự đoán.}
{Giả thuyết không cần kiểm tra.}
\loigiai{
Phát biểu đúng là: Sau khi nêu vấn đề có thể đề xuất giả thuyết hoặc dự đoán. Tiến trình cơ bản: xác định vấn đề, dự đoán, thiết kế và thực hiện thí nghiệm, xử lí dữ liệu, kết luận và điều chỉnh.
}
\end{ex}

% ===== Câu 219 =====
% ID: L10C1B1-50-DS
% Mức: NB
\dangbt{Đối tượng nghiên cứu và mục tiêu của môn Vật lí}
\begin{ex}
% ID: L10C1B1-50-DS
% Mức: NB
Xét các phát biểu sau về “Đối tượng nghiên cứu và mục tiêu của môn Vật lí” (bộ 3).
\choiceTF
{\True Chuyển động của hành tinh thuộc cấp độ vĩ mô.}
{Vật lí hoàn toàn tách rời công nghệ.}
{\True Nghiên cứu ánh sáng thuộc lĩnh vực Vật lí.}
{Mọi vấn đề sinh học đều chỉ thuộc Vật lí.}
\loigiai{
\begin{itemchoice}
\item (Đúng) Chuyển động của hành tinh thuộc cấp độ vĩ mô. (Vì): Chuyển động của hành tinh thuộc cấp độ vĩ mô. b) S: Vật lí hoàn toàn tách rời công nghệ. c) Đ: Nghiên cứu ánh sáng thuộc lĩnh vực Vật lí. d) S: Mọi vấn đề sinh học đều chỉ thuộc Vật lí. Xác định đúng đối tượng, phạm vi và mục tiêu của Vật lí; phân biệt với đối tượng trung tâm của các môn khoa học khác
\item (Sai) Vật lí hoàn toàn tách rời công nghệ.
\item (Đúng) Nghiên cứu ánh sáng thuộc lĩnh vực Vật lí.
\item (Sai) Mọi vấn đề sinh học đều chỉ thuộc Vật lí.
\end{itemchoice}
}
\end{ex}

% ===== Câu 220 =====
% ID: L10C1B1-50-TN
% Mức: VD
\dangbt{Phương pháp thực nghiệm trong nghiên cứu Vật lí}
\begin{ex}
% ID: L10C1B1-50-TN
% Mức: VD
Trong dạng “Phương pháp thực nghiệm trong nghiên cứu Vật lí”, phát biểu nào sau đây đúng?
\choice
{Một phép đo duy nhất luôn tuyệt đối chính xác.}
{Thí nghiệm không cần lặp lại.}
{\True Biến độc lập là đại lượng người nghiên cứu chủ động thay đổi.}
{Kết quả trái dự đoán phải bị bỏ đi.}
\loigiai{
Phát biểu đúng là: Biến độc lập là đại lượng người nghiên cứu chủ động thay đổi. Tiến trình cơ bản: xác định vấn đề, dự đoán, thiết kế và thực hiện thí nghiệm, xử lí dữ liệu, kết luận và điều chỉnh.
}
\end{ex}

% ===== Câu 221 =====
% ID: L10C1B1-51-DS
% Mức: TH
\dangbt{Đối tượng nghiên cứu và mục tiêu của môn Vật lí}
\begin{ex}
% ID: L10C1B1-51-DS
% Mức: TH
Xét các phát biểu sau về “Đối tượng nghiên cứu và mục tiêu của môn Vật lí” (bộ 4).
\choiceTF
{\True Nghiên cứu dòng điện thuộc lĩnh vực Vật lí.}
{Vật lí không thể dự đoán hiện tượng.}
{\True Kiến thức Vật lí có thể dùng để giải quyết vấn đề thực tiễn.}
{Vật lí chỉ nghiên cứu vật thể nhìn thấy bằng mắt.}
\loigiai{
\begin{itemchoice}
\item (Đúng) Nghiên cứu dòng điện thuộc lĩnh vực Vật lí. (Vì): Nghiên cứu dòng điện thuộc lĩnh vực Vật lí. b) S: Vật lí không thể dự đoán hiện tượng. c) Đ: Kiến thức Vật lí có thể dùng để giải quyết vấn đề thực tiễn. d) S: Vật lí chỉ nghiên cứu vật thể nhìn thấy bằng mắt. Xác định đúng đối tượng, phạm vi và mục tiêu của Vật lí; phân biệt với đối tượng trung tâm của các môn khoa học khác
\item (Sai) Vật lí không thể dự đoán hiện tượng.
\item (Đúng) Kiến thức Vật lí có thể dùng để giải quyết vấn đề thực tiễn.
\item (Sai) Vật lí chỉ nghiên cứu vật thể nhìn thấy bằng mắt.
\end{itemchoice}
}
\end{ex}

% ===== Câu 222 =====
% ID: L10C1B1-51-TN
% Mức: VDC
\dangbt{Phương pháp thực nghiệm trong nghiên cứu Vật lí}
\begin{ex}
% ID: L10C1B1-51-TN
% Mức: VDC
Trong dạng “Phương pháp thực nghiệm trong nghiên cứu Vật lí”, phát biểu nào sau đây đúng?
\choice
{Không cần xác định vấn đề trước khi đo.}
{Một phép đo duy nhất luôn tuyệt đối chính xác.}
{Thí nghiệm không cần lặp lại.}
{\True Phương án thí nghiệm phải xác định đại lượng cần đo.}
\loigiai{
Phát biểu đúng là: Phương án thí nghiệm phải xác định đại lượng cần đo. Tiến trình cơ bản: xác định vấn đề, dự đoán, thiết kế và thực hiện thí nghiệm, xử lí dữ liệu, kết luận và điều chỉnh.
}
\end{ex}

% ===== Câu 223 =====
% ID: L10C1B1-52-DS
% Mức: VD
\dangbt{Đối tượng nghiên cứu và mục tiêu của môn Vật lí}
\begin{ex}
% ID: L10C1B1-52-DS
% Mức: VD
Xét các phát biểu sau về “Đối tượng nghiên cứu và mục tiêu của môn Vật lí” (bộ 1).
\choiceTF
{\True Vật lí nghiên cứu các dạng vận động của vật chất và năng lượng.}
{Vật lí không nghiên cứu năng lượng.}
{\True Electron thuộc đối tượng ở cấp độ vi mô.}
{Sự chuyển động của electron không thuộc Vật lí.}
\loigiai{
\begin{itemchoice}
\item (Đúng) Vật lí nghiên cứu các dạng vận động của vật chất và năng lượng. (Vì): Vật lí nghiên cứu các dạng vận động của vật chất và năng lượng. b) S: Vật lí không nghiên cứu năng lượng. c) Đ: Electron thuộc đối tượng ở cấp độ vi mô. d) S: Sự chuyển động của electron không thuộc Vật lí. Xác định đúng đối tượng, phạm vi và mục tiêu của Vật lí; phân biệt với đối tượng trung tâm của các môn khoa học khác
\item (Sai) Vật lí không nghiên cứu năng lượng.
\item (Đúng) Electron thuộc đối tượng ở cấp độ vi mô.
\item (Sai) Sự chuyển động của electron không thuộc Vật lí.
\end{itemchoice}
}
\end{ex}

% ===== Câu 224 =====
% ID: L10C1B1-52-TN
% Mức: VDC
\dangbt{Phương pháp thực nghiệm trong nghiên cứu Vật lí}
\begin{ex}
% ID: L10C1B1-52-TN
% Mức: VDC
Trong dạng “Phương pháp thực nghiệm trong nghiên cứu Vật lí”, phát biểu nào sau đây đúng?
\choice
{Không cần ghi chép số liệu.}
{Thực nghiệm chỉ cần kết luận, không cần dữ liệu.}
{Không cần xác định vấn đề trước khi đo.}
{\True Biến phụ thuộc là đại lượng được theo dõi theo biến độc lập.}
\loigiai{
Phát biểu đúng là: Biến phụ thuộc là đại lượng được theo dõi theo biến độc lập. Tiến trình cơ bản: xác định vấn đề, dự đoán, thiết kế và thực hiện thí nghiệm, xử lí dữ liệu, kết luận và điều chỉnh.
}
\end{ex}

% ===== Câu 225 =====
% ID: L10C1B1-53-DS
% Mức: VD
\dangbt{Đối tượng nghiên cứu và mục tiêu của môn Vật lí}
\begin{ex}
% ID: L10C1B1-53-DS
% Mức: VD
Xét các phát biểu sau về “Đối tượng nghiên cứu và mục tiêu của môn Vật lí” (bộ 5).
\choiceTF
{\True Nghiên cứu sự truyền nhiệt thuộc lĩnh vực Vật lí.}
{Vật lí không có ứng dụng trong đời sống.}
{\True Mục tiêu của Vật lí là tìm các quy luật tổng quát chi phối vật chất và năng lượng.}
{Mục tiêu của Vật lí chỉ là ghi chép hiện tượng.}
\loigiai{
\begin{itemchoice}
\item (Đúng) Nghiên cứu sự truyền nhiệt thuộc lĩnh vực Vật lí. (Vì): Nghiên cứu sự truyền nhiệt thuộc lĩnh vực Vật lí. b) S: Vật lí không có ứng dụng trong đời sống. c) Đ: Mục tiêu của Vật lí là tìm các quy luật tổng quát chi phối vật chất và năng lượng. d) S: Mục tiêu của Vật lí chỉ là ghi chép hiện tượng. Xác định đúng đối tượng, phạm vi và mục tiêu của Vật lí; phân biệt với đối tượng trung tâm của các môn khoa học khác
\item (Sai) Vật lí không có ứng dụng trong đời sống.
\item (Đúng) Mục tiêu của Vật lí là tìm các quy luật tổng quát chi phối vật chất và năng lượng.
\item (Sai) Mục tiêu của Vật lí chỉ là ghi chép hiện tượng.
\end{itemchoice}
}
\end{ex}

% ===== Câu 226 =====
% ID: L10C1B1-53-TN
% Mức: NB
\dangbt{Quá trình phát triển của Vật lí}
\begin{ex}
% ID: L10C1B1-53-TN
% Mức: NB
Trong dạng “Quá trình phát triển của Vật lí”, phát biểu nào sau đây đúng?
\choice
{\True Thực nghiệm giúp kiểm tra các dự đoán khoa học.}
{Vật lí chỉ phát triển nhờ suy luận chủ quan.}
{Lí thuyết khoa học không cần bằng chứng.}
{Thuyết lượng tử thuộc thời tiền Vật lí.}
\loigiai{
Phát biểu đúng là: Thực nghiệm giúp kiểm tra các dự đoán khoa học. Nhận biết các giai đoạn và đặc trưng phát triển của Vật lí, nhấn mạnh quan hệ giữa lí thuyết, thực nghiệm và thiết bị.
}
\end{ex}

% ===== Câu 227 =====
% ID: L10C1B1-54-DS
% Mức: VDC
\dangbt{Đối tượng nghiên cứu và mục tiêu của môn Vật lí}
\begin{ex}
% ID: L10C1B1-54-DS
% Mức: VDC
Xét các phát biểu sau về “Đối tượng nghiên cứu và mục tiêu của môn Vật lí” (bộ 2).
\choiceTF
{\True Vật lí nghiên cứu cả cấp độ vi mô và vĩ mô.}
{Vật lí chỉ nghiên cứu cấp độ vĩ mô.}
{\True Vật lí là cơ sở của nhiều ngành khoa học tự nhiên.}
{Phản ứng tạo chất mới là đối tượng trung tâm duy nhất của Vật lí.}
\loigiai{
\begin{itemchoice}
\item (Đúng) Vật lí nghiên cứu cả cấp độ vi mô và vĩ mô. (Vì): Vật lí nghiên cứu cả cấp độ vi mô và vĩ mô. b) S: Vật lí chỉ nghiên cứu cấp độ vĩ mô. c) Đ: Vật lí là cơ sở của nhiều ngành khoa học tự nhiên. d) S: Phản ứng tạo chất mới là đối tượng trung tâm duy nhất của Vật lí. Xác định đúng đối tượng, phạm vi và mục tiêu của Vật lí; phân biệt với đối tượng trung tâm của các môn khoa học khác
\item (Sai) Vật lí chỉ nghiên cứu cấp độ vĩ mô.
\item (Đúng) Vật lí là cơ sở của nhiều ngành khoa học tự nhiên.
\item (Sai) Phản ứng tạo chất mới là đối tượng trung tâm duy nhất của Vật lí.
\end{itemchoice}
}
\end{ex}

% ===== Câu 228 =====
% ID: L10C1B1-54-TN
% Mức: NB
\dangbt{Quá trình phát triển của Vật lí}
\begin{ex}
% ID: L10C1B1-54-TN
% Mức: NB
Trong dạng “Quá trình phát triển của Vật lí”, phát biểu nào sau đây đúng?
\choice
{\True Một lí thuyết khoa học cần được đối chiếu với bằng chứng.}
{Mọi quan niệm cổ đều luôn đúng.}
{Thiết bị đo cản trở mọi phát hiện.}
{Vật lí chỉ phát triển nhờ suy luận chủ quan.}
\loigiai{
Phát biểu đúng là: Một lí thuyết khoa học cần được đối chiếu với bằng chứng. Nhận biết các giai đoạn và đặc trưng phát triển của Vật lí, nhấn mạnh quan hệ giữa lí thuyết, thực nghiệm và thiết bị.
}
\end{ex}

% ===== Câu 229 =====
% ID: L10C1B1-55-TN
% Mức: NB
\begin{ex}
% ID: L10C1B1-55-TN
% Mức: NB
Trong dạng “Quá trình phát triển của Vật lí”, phát biểu nào sau đây đúng?
\choice
{\True Thuyết lượng tử thuộc Vật lí hiện đại.}
{Khoa học không bao giờ điều chỉnh lí thuyết.}
{Thí nghiệm không có vai trò kiểm chứng.}
{Mọi quan niệm cổ đều luôn đúng.}
\loigiai{
Phát biểu đúng là: Thuyết lượng tử thuộc Vật lí hiện đại. Nhận biết các giai đoạn và đặc trưng phát triển của Vật lí, nhấn mạnh quan hệ giữa lí thuyết, thực nghiệm và thiết bị.
}
\end{ex}

% ===== Câu 230 =====
% ID: L10C1B1-56-TN
% Mức: TH
\begin{ex}
% ID: L10C1B1-56-TN
% Mức: TH
Trong dạng “Quá trình phát triển của Vật lí”, phát biểu nào sau đây đúng?
\choice
{Thí nghiệm không có vai trò kiểm chứng.}
{\True Galileo góp phần đặt nền móng cho phương pháp thực nghiệm.}
{Mọi quan niệm cổ đều luôn đúng.}
{Thiết bị đo cản trở mọi phát hiện.}
\loigiai{
Phát biểu đúng là: Galileo góp phần đặt nền móng cho phương pháp thực nghiệm. Nhận biết các giai đoạn và đặc trưng phát triển của Vật lí, nhấn mạnh quan hệ giữa lí thuyết, thực nghiệm và thiết bị.
}
\end{ex}

% ===== Câu 231 =====
% ID: L10C1B1-57-TN
% Mức: TH
\begin{ex}
% ID: L10C1B1-57-TN
% Mức: TH
Trong dạng “Quá trình phát triển của Vật lí”, phát biểu nào sau đây đúng?
\choice
{Newton không liên quan cơ học.}
{\True Kết quả mới có thể làm lí thuyết được bổ sung hoặc điều chỉnh.}
{Khoa học không bao giờ điều chỉnh lí thuyết.}
{Thí nghiệm không có vai trò kiểm chứng.}
\loigiai{
Phát biểu đúng là: Kết quả mới có thể làm lí thuyết được bổ sung hoặc điều chỉnh. Nhận biết các giai đoạn và đặc trưng phát triển của Vật lí, nhấn mạnh quan hệ giữa lí thuyết, thực nghiệm và thiết bị.
}
\end{ex}

% ===== Câu 232 =====
% ID: L10C1B1-58-TN
% Mức: TH
\begin{ex}
% ID: L10C1B1-58-TN
% Mức: TH
Trong dạng “Quá trình phát triển của Vật lí”, phát biểu nào sau đây đúng?
\choice
{Quan sát không cung cấp dữ liệu.}
{\True Sự phát triển thiết bị tạo điều kiện cho phát hiện mới.}
{Vật lí hiện đại không nghiên cứu vi mô.}
{Newton không liên quan cơ học.}
\loigiai{
Phát biểu đúng là: Sự phát triển thiết bị tạo điều kiện cho phát hiện mới. Nhận biết các giai đoạn và đặc trưng phát triển của Vật lí, nhấn mạnh quan hệ giữa lí thuyết, thực nghiệm và thiết bị.
}
\end{ex}

% ===== Câu 233 =====
% ID: L10C1B1-59-TN
% Mức: VD
\begin{ex}
% ID: L10C1B1-59-TN
% Mức: VD
Trong dạng “Quá trình phát triển của Vật lí”, phát biểu nào sau đây đúng?
\choice
{Vật lí hiện đại không nghiên cứu vi mô.}
{Newton không liên quan cơ học.}
{\True Vật lí cổ điển phát triển mạnh từ thế kỉ XVII.}
{Khoa học không bao giờ điều chỉnh lí thuyết.}
\loigiai{
Phát biểu đúng là: Vật lí cổ điển phát triển mạnh từ thế kỉ XVII. Nhận biết các giai đoạn và đặc trưng phát triển của Vật lí, nhấn mạnh quan hệ giữa lí thuyết, thực nghiệm và thiết bị.
}
\end{ex}

% ===== Câu 234 =====
% ID: L10C1B1-60-TN
% Mức: VD
\begin{ex}
% ID: L10C1B1-60-TN
% Mức: VD
Trong dạng “Quá trình phát triển của Vật lí”, phát biểu nào sau đây đúng?
\choice
{Thuyết lượng tử thuộc thời tiền Vật lí.}
{Quan sát không cung cấp dữ liệu.}
{\True Quan sát và đo lường thúc đẩy sự phát triển của Vật lí.}
{Vật lí hiện đại không nghiên cứu vi mô.}
\loigiai{
Phát biểu đúng là: Quan sát và đo lường thúc đẩy sự phát triển của Vật lí. Nhận biết các giai đoạn và đặc trưng phát triển của Vật lí, nhấn mạnh quan hệ giữa lí thuyết, thực nghiệm và thiết bị.
}
\end{ex}

% ===== Câu 235 =====
% ID: L10C1B1-61-TN
% Mức: VDC
\begin{ex}
% ID: L10C1B1-61-TN
% Mức: VDC
Trong dạng “Quá trình phát triển của Vật lí”, phát biểu nào sau đây đúng?
\choice
{Lí thuyết khoa học không cần bằng chứng.}
{Thuyết lượng tử thuộc thời tiền Vật lí.}
{Quan sát không cung cấp dữ liệu.}
{\True Vật lí hiện đại nghiên cứu sâu thế giới vi mô.}
\loigiai{
Phát biểu đúng là: Vật lí hiện đại nghiên cứu sâu thế giới vi mô. Nhận biết các giai đoạn và đặc trưng phát triển của Vật lí, nhấn mạnh quan hệ giữa lí thuyết, thực nghiệm và thiết bị.
}
\end{ex}

% ===== Câu 236 =====
% ID: L10C1B1-62-TN
% Mức: VDC
\begin{ex}
% ID: L10C1B1-62-TN
% Mức: VDC
Trong dạng “Quá trình phát triển của Vật lí”, phát biểu nào sau đây đúng?
\choice
{Thiết bị đo cản trở mọi phát hiện.}
{Vật lí chỉ phát triển nhờ suy luận chủ quan.}
{Lí thuyết khoa học không cần bằng chứng.}
{\True Newton có đóng góp quan trọng cho cơ học cổ điển.}
\loigiai{
Phát biểu đúng là: Newton có đóng góp quan trọng cho cơ học cổ điển. Nhận biết các giai đoạn và đặc trưng phát triển của Vật lí, nhấn mạnh quan hệ giữa lí thuyết, thực nghiệm và thiết bị.
}
\end{ex}

% ===== Câu 237 =====
% ID: L10C1B1-63-TN
% Mức: NB
\dangbt{Vai trò của Vật lí đối với khoa học, kĩ thuật, công nghệ và đời sống}
\begin{ex}
% ID: L10C1B1-63-TN
% Mức: NB
Trong dạng “Vai trò của Vật lí đối với khoa học, kĩ thuật, công nghệ và đời sống”, phát biểu nào sau đây đúng?
\choice
{\True Vật lí cung cấp cơ sở cho nhiều công nghệ hiện đại.}
{Vật lí không liên quan đến y học.}
{Năng lượng gió không liên quan kiến thức Vật lí.}
{Bán dẫn là thành tựu không liên quan Vật lí.}
\loigiai{
Phát biểu đúng là: Vật lí cung cấp cơ sở cho nhiều công nghệ hiện đại. Phân tích được vai trò của Vật lí qua ví dụ cụ thể và đánh giá cả lợi ích lẫn hệ quả của ứng dụng công nghệ.
}
\end{ex}

% ===== Câu 238 =====
% ID: L10C1B1-64-TN
% Mức: NB
\begin{ex}
% ID: L10C1B1-64-TN
% Mức: NB
Trong dạng “Vai trò của Vật lí đối với khoa học, kĩ thuật, công nghệ và đời sống”, phát biểu nào sau đây đúng?
\choice
{\True Vật lí hỗ trợ chế tạo thiết bị đo chính xác.}
{Thiết bị đo không cần nguyên lí Vật lí.}
{Siêu âm không sử dụng sóng.}
{Vật lí không liên quan đến y học.}
\loigiai{
Phát biểu đúng là: Vật lí hỗ trợ chế tạo thiết bị đo chính xác. Phân tích được vai trò của Vật lí qua ví dụ cụ thể và đánh giá cả lợi ích lẫn hệ quả của ứng dụng công nghệ.
}
\end{ex}

% ===== Câu 239 =====
% ID: L10C1B1-65-TN
% Mức: NB
\begin{ex}
% ID: L10C1B1-65-TN
% Mức: NB
Trong dạng “Vai trò của Vật lí đối với khoa học, kĩ thuật, công nghệ và đời sống”, phát biểu nào sau đây đúng?
\choice
{\True Kính thiên văn giúp mở rộng nghiên cứu vũ trụ.}
{Kính thiên văn dùng để nghiên cứu vi khuẩn.}
{Internet không sử dụng thành tựu Vật lí.}
{Thiết bị đo không cần nguyên lí Vật lí.}
\loigiai{
Phát biểu đúng là: Kính thiên văn giúp mở rộng nghiên cứu vũ trụ. Phân tích được vai trò của Vật lí qua ví dụ cụ thể và đánh giá cả lợi ích lẫn hệ quả của ứng dụng công nghệ.
}
\end{ex}

% ===== Câu 240 =====
% ID: L10C1B1-66-TN
% Mức: TH
\begin{ex}
% ID: L10C1B1-66-TN
% Mức: TH
Trong dạng “Vai trò của Vật lí đối với khoa học, kĩ thuật, công nghệ và đời sống”, phát biểu nào sau đây đúng?
\choice
{Internet không sử dụng thành tựu Vật lí.}
{\True Máy chụp X-quang là một ứng dụng của Vật lí.}
{Thiết bị đo không cần nguyên lí Vật lí.}
{Siêu âm không sử dụng sóng.}
\loigiai{
Phát biểu đúng là: Máy chụp X-quang là một ứng dụng của Vật lí. Phân tích được vai trò của Vật lí qua ví dụ cụ thể và đánh giá cả lợi ích lẫn hệ quả của ứng dụng công nghệ.
}
\end{ex}

% ===== Câu 241 =====
% ID: L10C1B1-67-TN
% Mức: TH
\begin{ex}
% ID: L10C1B1-67-TN
% Mức: TH
Trong dạng “Vai trò của Vật lí đối với khoa học, kĩ thuật, công nghệ và đời sống”, phát biểu nào sau đây đúng?
\choice
{Vật lí chỉ có vai trò trong phòng thí nghiệm.}
{\True Ứng dụng Vật lí có thể vừa mang lợi ích vừa gây tác động môi trường.}
{Kính thiên văn dùng để nghiên cứu vi khuẩn.}
{Internet không sử dụng thành tựu Vật lí.}
\loigiai{
Phát biểu đúng là: Ứng dụng Vật lí có thể vừa mang lợi ích vừa gây tác động môi trường. Phân tích được vai trò của Vật lí qua ví dụ cụ thể và đánh giá cả lợi ích lẫn hệ quả của ứng dụng công nghệ.
}
\end{ex}

% ===== Câu 242 =====
% ID: L10C1B1-68-TN
% Mức: TH
\begin{ex}
% ID: L10C1B1-68-TN
% Mức: TH
Trong dạng “Vai trò của Vật lí đối với khoa học, kĩ thuật, công nghệ và đời sống”, phát biểu nào sau đây đúng?
\choice
{Tự động hóa không cần cảm biến.}
{\True Vật lí góp phần tự động hóa sản xuất.}
{Mọi ứng dụng công nghệ đều không có rủi ro.}
{Vật lí chỉ có vai trò trong phòng thí nghiệm.}
\loigiai{
Phát biểu đúng là: Vật lí góp phần tự động hóa sản xuất. Phân tích được vai trò của Vật lí qua ví dụ cụ thể và đánh giá cả lợi ích lẫn hệ quả của ứng dụng công nghệ.
}
\end{ex}

% ===== Câu 243 =====
% ID: L10C1B1-69-TN
% Mức: VD
\begin{ex}
% ID: L10C1B1-69-TN
% Mức: VD
Trong dạng “Vai trò của Vật lí đối với khoa học, kĩ thuật, công nghệ và đời sống”, phát biểu nào sau đây đúng?
\choice
{Mọi ứng dụng công nghệ đều không có rủi ro.}
{Vật lí chỉ có vai trò trong phòng thí nghiệm.}
{\True Thông tin liên lạc phát triển dựa trên điện từ học và quang học.}
{Kính thiên văn dùng để nghiên cứu vi khuẩn.}
\loigiai{
Phát biểu đúng là: Thông tin liên lạc phát triển dựa trên điện từ học và quang học. Phân tích được vai trò của Vật lí qua ví dụ cụ thể và đánh giá cả lợi ích lẫn hệ quả của ứng dụng công nghệ.
}
\end{ex}

% ===== Câu 244 =====
% ID: L10C1B1-70-TN
% Mức: VD
\begin{ex}
% ID: L10C1B1-70-TN
% Mức: VD
Trong dạng “Vai trò của Vật lí đối với khoa học, kĩ thuật, công nghệ và đời sống”, phát biểu nào sau đây đúng?
\choice
{Bán dẫn là thành tựu không liên quan Vật lí.}
{Tự động hóa không cần cảm biến.}
{\True Công nghệ bán dẫn dựa trên hiểu biết Vật lí chất rắn.}
{Mọi ứng dụng công nghệ đều không có rủi ro.}
\loigiai{
Phát biểu đúng là: Công nghệ bán dẫn dựa trên hiểu biết Vật lí chất rắn. Phân tích được vai trò của Vật lí qua ví dụ cụ thể và đánh giá cả lợi ích lẫn hệ quả của ứng dụng công nghệ.
}
\end{ex}

% ===== Câu 245 =====
% ID: L10C1B1-71-TN
% Mức: VDC
\begin{ex}
% ID: L10C1B1-71-TN
% Mức: VDC
Trong dạng “Vai trò của Vật lí đối với khoa học, kĩ thuật, công nghệ và đời sống”, phát biểu nào sau đây đúng?
\choice
{Năng lượng gió không liên quan kiến thức Vật lí.}
{Bán dẫn là thành tựu không liên quan Vật lí.}
{Tự động hóa không cần cảm biến.}
{\True Kiến thức Vật lí góp phần phát triển năng lượng tái tạo.}
\loigiai{
Phát biểu đúng là: Kiến thức Vật lí góp phần phát triển năng lượng tái tạo. Phân tích được vai trò của Vật lí qua ví dụ cụ thể và đánh giá cả lợi ích lẫn hệ quả của ứng dụng công nghệ.
}
\end{ex}

% ===== Câu 246 =====
% ID: L10C1B1-72-TN
% Mức: VDC
\begin{ex}
% ID: L10C1B1-72-TN
% Mức: VDC
Trong dạng “Vai trò của Vật lí đối với khoa học, kĩ thuật, công nghệ và đời sống”, phát biểu nào sau đây đúng?
\choice
{Siêu âm không sử dụng sóng.}
{Vật lí không liên quan đến y học.}
{Năng lượng gió không liên quan kiến thức Vật lí.}
{\True Siêu âm y học có cơ sở từ sự truyền sóng.}
\loigiai{
Phát biểu đúng là: Siêu âm y học có cơ sở từ sự truyền sóng. Phân tích được vai trò của Vật lí qua ví dụ cụ thể và đánh giá cả lợi ích lẫn hệ quả của ứng dụng công nghệ.
}
\end{ex}

% ===== Câu 247 =====
% ID: L10C1B1-73-TN
% Mức: NB
\dangbt{Vận dụng tiến trình của phương pháp mô hình}
\begin{ex}
% ID: L10C1B1-73-TN
% Mức: NB
Sự ra đời của vật lí hiện đại vào đầu thế kỉ XX chủ yếu xuất phát từ việc vật lí cổ điển gặp phải giới hạn nào trong việc giải thích các hiện tượng tự nhiên?
\choice
{\True Khả năng giải thích các hiện tượng vĩ mô}
{Khả năng giải thích các hiện tượng vi mô}
{Khả năng giải thích các hiện tượng ở tốc độ ánh sáng}
{Khả năng giải thích các hiện tượng ở nhiệt độ thấp}
\loigiai{
Vật lí cổ điển, phát triển mạnh mẽ từ thế kỉ XVII đến cuối thế kỉ XIX, đã rất thành công trong việc mô tả các hiện tượng vĩ mô và chuyển động với tốc độ thông thường. Tuy nhiên, vào cuối thế kỉ XIX và đầu thế kỉ XX, các nhà khoa học bắt đầu phát hiện ra những hiện tượng mà vật lí cổ điển không thể giải thích được hoặc giải thích không chính xác. Cụ thể, vật lí cổ điển gặp khó khăn khi nghiên cứu:
 - Các hiện tượng ở tốc độ rất lớn, gần tốc độ ánh sáng, dẫn đến sự ra đời của Thuyết tương đối hẹp và rộng.
 - Các hiện tượng ở kích thước rất nhỏ, cấp độ nguyên tử và hạ nguyên tử, dẫn đến sự phát triển của Cơ học lượng tử.
 
 Các phương án khác không đúng:
 - B: Vật lí cổ điển đã có những phương pháp thực nghiệm rất tinh vi và là nền tảng cho nhiều khám phá quan trọng.
 - C: Vật lí cổ điển đã sử dụng toán học (đặc biệt là giải tích) một cách mạnh mẽ để xây dựng các lí thuyết và mô hình.
 - D: Các định luật của vật lí cổ điển (như định luật vạn vật hấp dẫn của Newton) đã giải thích rất thành công chuyển động của các vật thể trong hệ Mặt Trời.
}
\end{ex}

% ===== Câu 248 =====
% ID: L10C1B1-74-TN
% Mức: NB
\begin{ex}
% ID: L10C1B1-74-TN
% Mức: NB
Chat GPT ra đời trong cuộc cách mạng công nghiệp lần thứ
\choice
{\True Tư}
{Ba}
{Hai}
{nhất}
\loigiai{
Chat GPT là ứng dụng trí tuệ nhân tạo, thuộc nhóm thành tựu tiêu biểu của cuộc cách mạng công nghiệp lần thứ tư.
}
\end{ex}

% ===== Câu 249 =====
% ID: L10C1B1-75-TN
% Mức: NB
\begin{ex}
% ID: L10C1B1-75-TN
% Mức: NB
Đặc trưng của cuộc cách mạng công nghiệp lần thứ hai là
\choice
{tự động hóa các quá trình sản xuất}
{\True sử dụng các thiết bị điện trong mọi lĩnh vực của đời sống}
{thay thế sức lực cơ bắp bằng máy móc}
{sử dụng trí tuện nhân tạo, robot và internet toàn cầu}
\loigiai{
Cuộc cách mạng công nghiệp lần thứ hai có đặc trưng nổi bật là sử dụng điện năng và các thiết bị điện trong nhiều lĩnh vực.
}
\end{ex}

% ===== Câu 250 =====
% ID: L10C1B1-76-TN
% Mức: NB
\begin{ex}
% ID: L10C1B1-76-TN
% Mức: NB
Sự kiện Vật lí nào sau đây xảy ra vào năm $1831$ ?
\choice
{Galilei làm thí nghiệm tại tháp nghiêng Pisa}
{Einstein xây dựng thuyết tương đối}
{Joule tìm ra các định luật nhiệt động lực học}
{\True Faraday tìm ra hiện tượng cảm ứng điện từ}
\loigiai{
Năm $1831$, Faraday phát hiện hiện tượng cảm ứng điện từ. Đây là cơ sở quan trọng cho sự ra đời của máy phát điện và động cơ điện.
}
\end{ex}

% ===== Câu 251 =====
% ID: L10C1B1-77-TN
% Mức: NB
\begin{ex}
% ID: L10C1B1-77-TN
% Mức: NB
Sự kiện Vật lí nào sau đây xảy ra vào năm $1785$ ?
\choice
{Einstein xây dựng thuyết tương đối}
{\True Joule tìm ra các định luật nhiệt động lực học}
{Faraday tìm ra hiện tượng cảm ứng điện từ}
{Galilei làm thí nghiệm tại tháp nghiêng Pisa}
\loigiai{
Theo đáp án đã đánh dấu trong đề, sự kiện được chọn là Joule tìm ra các định luật nhiệt động lực học.
}
\end{ex}

% ===== Câu 252 =====
% ID: L10C1B1-78-TN
% Mức: NB
\begin{ex}
% ID: L10C1B1-78-TN
% Mức: NB
Quá trình phát triển của vật lí được chia thành bao nhiêu giai đoạn ?
\choice
{5}
{\True 4}
{2}
{3}
\loigiai{
Quá trình phát triển của Vật lí thường được chia thành bốn giai đoạn chính, từ thời tiền Vật lí đến Vật lí hiện đại.
}
\end{ex}

% ===== Câu 253 =====
% ID: L10C1B1-79-TN
% Mức: NB
\begin{ex}
% ID: L10C1B1-79-TN
% Mức: NB
Thành tựu nghiên cứu nào sau đây của Vật lí được coi là có vai trò quan trọng trong việc mở đầu cho cuộc cách mạng công nghệ lần thứ hai ?
\choice
{Nghiên cứu về thuyết tương đối}
{Nghiên cứu về lực vạn vật hấp dẫn}
{\True Nghiên cứu về cảm ứng điện từ}
{Nghiên cứu về nhiệt động lực học}
\loigiai{
Hiện tượng cảm ứng điện từ là cơ sở cho máy phát điện và nhiều thiết bị điện, góp phần quan trọng mở đầu cho thời kì sử dụng điện năng rộng rãi.
}
\end{ex}

% ===== Câu 254 =====
% ID: L10C1B1-80-TN
% Mức: NB
\begin{ex}
% ID: L10C1B1-80-TN
% Mức: NB
Cấp độ vĩ mô là:
\choice
{cấp độ dùng để mô phỏng vật chất nhỏ bé}
{cấp độ to, nhỏ tùy thuộc vào quy mô được khảo sát}
{\True cấp độ dùng để mô phỏng tầm rộng lớn hay rất lớn của vật chất}
{cấp độ tinh vi khi khảo sát một hiện tượng vật lí}
\loigiai{
Vĩ mô là cấp độ khảo sát các hệ vật thể lớn như xe, tàu, hành tinh,...
}
\end{ex}

% ===== Câu 255 =====
% ID: L10C1B1-81-TN
% Mức: NB
\begin{ex}
% ID: L10C1B1-81-TN
% Mức: NB
Thiết bị nào sau đây có ứng dụng kiến thức về nhiệt là chủ yếu ?
\choice
{Cân điện tử}
{Ti vi}
{\True Nhiệt kế}
{Điện thoại}
\loigiai{
Nhiệt kế dùng để đo nhiệt độ, hoạt động dựa trên các kiến thức về nhiệt học như sự nở vì nhiệt hoặc sự thay đổi tính chất vật lí theo nhiệt độ.
}
\end{ex}

% ===== Câu 256 =====
% ID: L10C1B1-82-TN
% Mức: NB
\begin{ex}
% ID: L10C1B1-82-TN
% Mức: NB
Thành tựu vật lí nào sau đây không thuộc cuộc cách mạng khoa học lần thứ tư ?
\choice
{Ô tô không người lái}
{Điện thoại}
{Rôbốt}
{\True Động cơ hơi nước}
\loigiai{
Động cơ hơi nước là thành tựu tiêu biểu của cuộc cách mạng công nghiệp lần thứ nhất, không thuộc cuộc cách mạng khoa học lần thứ tư.
}
\end{ex}

% ===== Câu 257 =====
% ID: L10C1B1-83-TN
% Mức: NB
\begin{ex}
% ID: L10C1B1-83-TN
% Mức: NB
Trong các lĩnh vực dưới đây, lĩnh vực nào KHÔNG thuộc phạm vi nghiên cứu của Vật lí?
\choice
{Nghiên cứu về sự hình thành và phát triển của các loài sinh vật}
{\True Nghiên cứu về chuyển động của các vật thể}
{Nghiên cứu về các dạng năng lượng và sự biến đổi giữa chúng}
{Nghiên cứu về cấu tạo của vật chất và tương tác giữa các hạt cơ bản}
\loigiai{
Vật lí nghiên cứu về các hiện tượng tự nhiên liên quan đến vật chất, năng lượng, không gian và thời gian. Sự chuyển động của các hành tinh, nguyên tắc hoạt động của động cơ và sự lan truyền sóng âm đều thuộc phạm vi nghiên cứu của Vật lí. Cơ chế di truyền và biến dị là lĩnh vực nghiên cứu của Sinh học.
}
\end{ex}

% ===== Câu 258 =====
% ID: L10C1B1-84-TN
% Mức: NB
\begin{ex}
% ID: L10C1B1-84-TN
% Mức: NB
Kiến thức về từ trường Trái Đất được dùng để giải thích đặc điểm nào của loài chim di trú ?
\choice
{Làm tổ}
{Sinh sản}
{Kiếm ăn}
{\True Xác định hướng bay}
\loigiai{
Một số loài chim di trú có khả năng định hướng nhờ từ trường Trái Đất. Vì vậy kiến thức về từ trường được dùng để giải thích khả năng xác định hướng bay của chúng.
}
\end{ex}

% ===== Câu 259 =====
% ID: L10C1B1-85-TN
% Mức: NB
\begin{ex}
% ID: L10C1B1-85-TN
% Mức: NB
Cấp độ vi mô là:
\choice
{\True cấp độ dùng để mô phỏng vật chất nhỏ bé}
{cấp độ to, nhỏ tùy thuộc vào quy mô được khảo sát}
{cấp độ dùng để mô phỏng tầm rộng lớn hay rất lớn của vật chất}
{cấp độ tinh vi khi khảo sát một hiện tượng vật lí}
\loigiai{
Vi mô là cấp độ nghiên cứu những hệ thống rất nhỏ như nguyên tử, phân tử, hạt cơ bản.
}
\end{ex}

% ===== Câu 260 =====
% ID: L10C1B1-86-TN
% Mức: NB
\begin{ex}
% ID: L10C1B1-86-TN
% Mức: NB
Đặc trưng của cuộc cách mạng công nghiệp lần thứ nhất là
\choice
{\True thay thế sức lực cơ bắp bằng máy móc}
{sử dụng các thiết bị điện trong mọi lĩnh vực của đời sống}
{sử dụng trí tuện nhân tạo, robot và internet toàn cầu}
{tự động hóa các quá trình sản xuất}
\loigiai{
Cuộc cách mạng công nghiệp lần thứ nhất gắn với cơ khí hóa sản xuất, đặc biệt là việc dùng máy móc thay thế sức lực cơ bắp.
}
\end{ex}

% ===== Câu 261 =====
% ID: L10C1B1-87-TN
% Mức: NB
\begin{ex}
% ID: L10C1B1-87-TN
% Mức: NB
Mục tiêu nào sau đây thể hiện rõ nhất vai trò của Vật lí trong việc giải thích và dự đoán các hiện tượng tự nhiên?
\choice
{\True Giải thích và dự đoán các hiện tượng tự nhiên}
{Nghiên cứu các định luật vật lí}
{Phát triển các công nghệ mới}
{Giáo dục và đào tạo con người}
\loigiai{
Vật lí là một ngành khoa học cơ bản, có đối tượng nghiên cứu là vật chất, năng lượng và sự tương tác giữa chúng. Mục tiêu cốt lõi của Vật lí là tìm hiểu các quy luật tự nhiên, từ đó xây dựng các mô hình và định luật để giải thích các hiện tượng đã quan sát được và dự đoán các hiện tượng mới sẽ xảy ra. Phương án $A$ mô tả chính xác mục tiêu này, nhấn mạnh vào việc giải thích và dự đoán các quy luật tự nhiên.
 Phương án $B$ là một ứng dụng quan trọng và là hệ quả của việc nghiên cứu Vật lí, nhưng không phải là mục tiêu cốt lõi trong việc giải thích và dự đoán hiện tượng tự nhiên.
 Các phương án $C$ và $D$ thuộc về các lĩnh vực khoa học khác $(Sinh học và Khoa học xã hội)$, không phải là đối tượng nghiên cứu chính của Vật lí.
 Vậy, đáp án đúng là A.
}
\end{ex}

% ===== Câu 262 =====
% ID: L10C1B1-88-TN
% Mức: NB
\begin{ex}
% ID: L10C1B1-88-TN
% Mức: NB
Động cơ điện ra đời trong cuộc cách mạng công nghiệp lần thứ mấy ?
\choice
{Lần thứ ba}
{\True Lần thứ hai}
{Lần thứ nhất}
{Lần thứ tư}
\loigiai{
Động cơ điện gắn với việc ứng dụng điện năng rộng rãi trong sản xuất và đời sống, là đặc trưng của cuộc cách mạng công nghiệp lần thứ hai.
}
\end{ex}

% ===== Câu 263 =====
% ID: L10C1B1-89-TN
% Mức: NB
\begin{ex}
% ID: L10C1B1-89-TN
% Mức: NB
Phương pháp thực nghiệm gồm các bước nào sau đây ?
\choice
{Quan sát, suy luận, kết luận}
{Xác định đối tượng nghiên cứu, quan sát thu thập thông tin, đưa ra dự đoán, kết luận}
{\True Xác định vấn đề nghiên cứu, quan sát thu thập thông tin, đưa ra dự đoán, thí nghiệm kiểm tra, kết luận}
{Xác định đối tượng nghiên cứu, xây dựng mô hình, kiểm tra mô hình, điều chỉnh mô hình, kết luận}
\loigiai{
Phương pháp thực nghiệm gồm các bước cơ bản: xác định vấn đề nghiên cứu, quan sát và thu thập thông tin, đưa ra dự đoán, tiến hành thí nghiệm kiểm tra, sau đó rút ra kết luận.
}
\end{ex}

% ===== Câu 264 =====
% ID: L10C1B1-90-TN
% Mức: NB
\begin{ex}
% ID: L10C1B1-90-TN
% Mức: NB
Lĩnh vực nào sau đây KHÔNG phải là trọng tâm nghiên cứu chính của vật lí cổ điển (trước thế kỉ XX)?
\choice
{Cơ học}
{\True Thuyết tương đối}
{Nhiệt động lực học}
{Điện từ học}
\loigiai{
Vật lí cổ điển, phát triển mạnh mẽ từ thế kỉ XVII đến cuối thế kỉ XIX, chủ yếu tập trung vào các lĩnh vực như Cơ học Newton, Nhiệt học, Điện từ học Maxwell và Quang học. Các lĩnh vực như vật lí hạt nhân và vật lí lượng tử, nghiên cứu cấu trúc và tương tác của hạt hạ nguyên tử, chỉ thực sự phát triển và trở thành trọng tâm nghiên cứu của vật lí học từ đầu thế kỉ XX trở đi, đánh dấu sự ra đời của vật lí hiện đại.
}
\end{ex}

% ===== Câu 265 =====
% ID: L10C1B1-91-TN
% Mức: NB
\begin{ex}
% ID: L10C1B1-91-TN
% Mức: NB
Bước nào sau đây không thuộc bốn bước cơ bản của tiến trình phương pháp mô hình trong nghiên cứu Vật lí?
\choice
{Xác định đối tượng cần được mô hình hóa}
{Xây dựng mô hình (giả thuyết)}
{Kiểm tra sự phù hợp của các mô hình với thực tế}
{\True Tiến hành thí nghiệm thả rơi quả cầu tại tháp nghiêng Pisa}
\loigiai{
Tiến hành thí nghiệm thả rơi quả cầu tại tháp nghiêng Pisa là một thí nghiệm thực tế cụ thể thuộc tiến trình phương pháp thực nghiệm của Galilei, không phải là một bước khái quát trong sơ đồ tiến trình phương pháp mô hình.
}
\end{ex}

% ===== Câu 266 =====
% ID: L10C1B1-92-TN
% Mức: NB
\begin{ex}
% ID: L10C1B1-92-TN
% Mức: NB
Sự kiện Vật lí nào sau đây xảy ra vào năm $1687$ ?
\choice
{\True Newton công bố các nguyên lí toán của triết học tự nhiên}
{Joule tìm ra các định luật nhiệt động lực học}
{Einstein xây dựng thuyết tương đối}
{Galilei làm thí nghiệm tại tháp nghiêng Pisa}
\loigiai{
Năm $1687$, Newton công bố tác phẩm về các nguyên lí toán học của triết học tự nhiên, đặt nền móng quan trọng cho cơ học cổ điển.
}
\end{ex}

% ===== Câu 267 =====
% ID: L10C1B1-93-TN
% Mức: NB
\begin{ex}
% ID: L10C1B1-93-TN
% Mức: NB
Ý nào dưới đây không phải là vai trò của khoa học tự nhiên trong đời sống ?
\choice
{\True Định hướng tư tưởng, phát triển hệ thống chính trị}
{Bảo vệ môi trường, ứng phó với biến đổi khí hậu}
{Bảo vệ sức khỏe và cuộc sống của con người}
{Mở rộng sản xuất, phát triển kinh tế}
\loigiai{
Khoa học tự nhiên có vai trò trong bảo vệ môi trường, sức khỏe, đời sống và phát triển sản xuất. Việc định hướng tư tưởng, phát triển hệ thống chính trị không phải vai trò trực tiếp của khoa học tự nhiên.
}
\end{ex}

% ===== Câu 268 =====
% ID: L10C1B1-94-TN
% Mức: NB
\begin{ex}
% ID: L10C1B1-94-TN
% Mức: NB
Cơ chế của các phản ứng hóa học được giải thích dựa trên kiến thức thuộc lĩnh vực nào của Vật lí ?
\choice
{Nhiệt học}
{\True Lượng tử}
{Cơ học}
{Quang học}
\loigiai{
Cơ chế phản ứng hóa học liên quan đến cấu trúc nguyên tử, phân tử và sự phân bố electron. Những nội dung này được giải thích sâu bằng kiến thức vật lí lượng tử.
}
\end{ex}

% ===== Câu 269 =====
% ID: L10C1B1-95-TN
% Mức: NB
\begin{ex}
% ID: L10C1B1-95-TN
% Mức: NB
Kết luận nào sau đây là sai khi nói về ảnh hưởng của vật lí đến một số lĩnh vực trong đời sống và kĩ thuật ?
\choice
{Vật lí có ảnh hưởng mạnh mẽ và có tác dụng làm thay đổi mọi lĩnh vực hoạt động của con người}
{\True Vật lí đem lại cho con người những lợi ích tuyệt vời và không gây ra một ảnh hưởng xấu nào}
{Kiến thức vật lí trong các phân ngành được áp dụng kết hợp để tạo ra kết quả tối ưu}
{Vật lí là cơ sở của khoa học tự nhiên và công nghệ}
\loigiai{
Vật lí đem lại nhiều lợi ích cho con người, nhưng một số ứng dụng nếu sử dụng không đúng cách cũng có thể gây tác động xấu. Vì vậy khẳng định vật lí không gây ra bất kì ảnh hưởng xấu nào là sai.
}
\end{ex}

% ===== Câu 270 =====
% ID: L10C1B1-96-TN
% Mức: NB
\begin{ex}
% ID: L10C1B1-96-TN
% Mức: NB
Biểu hiện nào sau đây không phải là biểu hiện của phát triển năng lực vật lí ?
\choice
{Có được kiến thức, kĩ năng cơ bản về vật lí}
{\True Nhận biết được hạn chế của bản thân để tìm cách khắc phục}
{Nhận biết được năng lực, sở trường của bản thân, định hướng nghề nghiệp}
{Vận dụng được kiến thức, kĩ năng để khám phá, giải quyết các vấn đề có liên quan trong học tập cũng như trong cuộc sống}
\loigiai{
Năng lực vật lí thể hiện qua việc nhận thức kiến thức vật lí, tìm hiểu thế giới tự nhiên dưới góc độ vật lí và vận dụng kiến thức vật lí vào thực tiễn. Việc nhận biết hạn chế của bản thân là biểu hiện chung về năng lực tự chủ, không phải biểu hiện đặc trưng của năng lực vật lí.
}
\end{ex}

% ===== Câu 271 =====
% ID: L10C1B1-97-TN
% Mức: NB
\begin{ex}
% ID: L10C1B1-97-TN
% Mức: NB
Phát biểu nào sau đây mô tả đúng nhất sự thay đổi cơ bản trong phương pháp nghiên cứu tự nhiên, đánh dấu sự hình thành vật lí học theo phương pháp khoa học hiện đại so với thời kì cổ đại?
\choice
{\True Tập trung vào việc quan sát và mô tả các hiện tượng tự nhiên một cách chi tiết hơn}
{Sử dụng các công cụ và thiết bị đo lường chính xác hơn để thu thập dữ liệu}
{Kết hợp chặt chẽ giữa quan sát, thực nghiệm và sử dụng toán học để mô tả và dự đoán các hiện tượng}
{Chuyển từ giải thích dựa trên thần thoại sang giải thích dựa trên logic và suy luận}
\loigiai{
Vật lí học hiện đại, bắt đầu từ thời kì Phục hưng và Cách mạng khoa học, đã chuyển mình từ việc chỉ dựa vào suy luận triết học và quan sát đơn thuần sang một phương pháp khoa học chặt chẽ hơn. Phương pháp này nhấn mạnh vai trò của thực nghiệm để kiểm chứng các giả thuyết và lí thuyết, đồng thời sử dụng toán học để mô tả định lượng các hiện tượng. Do đó, sự kết hợp giữa quan sát, thực nghiệm và kiểm chứng lí thuyết là thay đổi cơ bản nhất.
}
\end{ex}

% ===== Câu 272 =====
% ID: L10C1B1-98-TN
% Mức: NB
\begin{ex}
% ID: L10C1B1-98-TN
% Mức: NB
trí tuệ nhân tạo ( AI ) ra đời trong cuộc cách mạng công nghiệp lần thứ
\choice
{Ba}
{Hai}
{nhất}
{\True Tư}
\loigiai{
Trí tuệ nhân tạo là một trong những công nghệ tiêu biểu của cuộc cách mạng công nghiệp lần thứ tư.
}
\end{ex}

% ===== Câu 273 =====
% ID: L10C1B1-99-TN
% Mức: NB
\begin{ex}
% ID: L10C1B1-99-TN
% Mức: NB
Kết luận nào sau đây về ô tô điện là chưa đúng ?
\choice
{\True Hoạt động bằng nhiên liệu}
{Hoạt động bằng năng lượng Mặt Trời}
{Hoạt động bằng pin acquy}
{Thân thiện với môi trường}
\loigiai{
Ô tô điện chủ yếu hoạt động nhờ điện năng được tích trữ trong pin hoặc acquy, không hoạt động trực tiếp bằng nhiên liệu như xe dùng động cơ đốt trong.
}
\end{ex}

% ===== Câu 274 =====
% ID: AIQ_AUTO_01
% Mức: NB
\dangbt{Nhận diện và sắp xếp các bước của phương pháp thực nghiệm}
\begin{bt}
% ID: AIQ_AUTO_01
% Mức: NB
Trình bày các bước của phương pháp thực nghiệm trong nghiên cứu Vật lí.
\loigiai{
Các bước của phương pháp thực nghiệm: 1. Xác định vấn đề cần nghiên cứu. 2. Quan sát, thu thập thông tin. 3. Đưa ra dự đoán. 4. Thí nghiệm kiểm tra dự đoán. 5. Kết luận.
}
\end{bt}

% ===== Câu 275 =====
% ID: AIQ_AUTO_02
% Mức: NB
\begin{ex}
% ID: AIQ_AUTO_02
% Mức: NB
Có bao nhiêu bước cơ bản trong tiến trình tìm hiểu thế giới tự nhiên dưới góc độ vật lí theo phương pháp thực nghiệm?
\shortans{5}
\loigiai{
Có 5 bước cơ bản: Xác định vấn đề, Quan sát, Dự đoán, Thí nghiệm, Kết luận.
}
\end{ex}

% ===== Câu 276 =====
% ID: AIQ_AUTO_03
% Mức: NB
\dangbt{Phân biệt các loại mô hình trong nghiên cứu Vật lí}
\begin{ex}
% ID: AIQ_AUTO_03
% Mức: NB
Công thức tính vận tốc $v = \frac{s}{t}$ là ví dụ của loại mô hình nào?
\choice
{Mô hình vật chất}
{\True Mô hình toán học}
{Mô hình lí thuyết}
{Mô hình đồ thị}
\loigiai{
Công thức tính vận tốc sử dụng phương trình toán học để mô tả đại lượng vật lí nên là mô hình toán học.
}
\end{ex}

% ===== Câu 277 =====
% ID: AIQ_AUTO_04
% Mức: NB
\begin{ex}
% ID: AIQ_AUTO_04
% Mức: NB
Quả địa cầu là một ví dụ về mô hình gì?
\shortans{Mô hình vật chất}
\loigiai{
Quả địa cầu là vật thể thu nhỏ mô phỏng Trái Đất nên thuộc mô hình vật chất.
}
\end{ex}

% ===== Câu 278 =====
% ID: AIQ_AUTO_05
% Mức: NB
\dangbt{Phân tích vai trò và ảnh hưởng hai mặt của Vật lí đối với đời sống, khoa học và công nghệ}
\begin{ex}
% ID: AIQ_AUTO_05
% Mức: NB
Thành tựu nào sau đây của Vật lí có ảnh hưởng tiêu cực đến môi trường nếu không được kiểm soát tốt?
\choice
{Chế tạo pin mặt trời}
{\True Phát triển năng lượng hạt nhân}
{Chế tạo đèn LED}
{Phát triển xe điện}
\loigiai{
Năng lượng hạt nhân nếu không được kiểm soát an toàn có thể gây rò rỉ phóng xạ, ảnh hưởng tiêu cực đến môi trường.
}
\end{ex}

% ===== Câu 279 =====
% ID: AIQ_AUTO_06
% Mức: NB
\begin{ex}
% ID: AIQ_AUTO_06
% Mức: NB
Đánh giá các nhận định sau về ảnh hưởng của Vật lí:
\choiceTF
{Vật lí chỉ mang lại lợi ích cho con người mà không có tác động tiêu cực nào.}
{\True Việc ứng dụng Vật lí vào công nghệ thông tin giúp con người liên lạc dễ dàng hơn.}
{\True Khí thải từ các động cơ nhiệt do Vật lí chế tạo có thể gây ô nhiễm môi trường.}
{Vật lí không có vai trò gì trong sự phát triển của y học hiện đại.}
\loigiai{
\begin{itemchoice}
\item (Sai) Vật lí chỉ mang lại lợi ích cho con người mà không có tác động tiêu cực nào. (Vì): Vật lí cũng có thể mang lại tác động tiêu cực nếu bị lạm dụng (ví dụ: vũ khí hạt nhân). B. Đúng. C. Đúng. D. Sai - Vật lí đóng vai trò quan trọng trong y học (X-quang, MRI
\item (Đúng) Việc ứng dụng Vật lí vào công nghệ thông tin giúp con người liên lạc dễ dàng hơn.
\item (Đúng) Khí thải từ các động cơ nhiệt do Vật lí chế tạo có thể gây ô nhiễm môi trường.
\item (Sai) Vật lí không có vai trò gì trong sự phát triển của y học hiện đại.
\end{itemchoice}
}
\end{ex}

% ===== Câu 280 =====
% ID: AIQ_AUTO_07
% Mức: NB
\dangbt{Phương pháp thực nghiệm trong nghiên cứu Vật lí}
\begin{ex}
% ID: AIQ_AUTO_07
% Mức: NB
Thí nghiệm kiểm tra dự đoán có vai trò gì trong phương pháp thực nghiệm?
\choice
{Xác định vấn đề cần nghiên cứu}
{Hình thành giả thuyết}
{\True Chứng minh hoặc bác bỏ dự đoán}
{Rút ra kết luận cuối cùng mà không cần phân tích}
\loigiai{
Thí nghiệm giúp kiểm chứng xem dự đoán đưa ra có đúng với thực tế hay không, từ đó chứng minh hoặc bác bỏ nó.
}
\end{ex}

% ===== Câu 281 =====
% ID: AIQ_AUTO_08
% Mức: NB
\begin{ex}
% ID: AIQ_AUTO_08
% Mức: NB
Nhận định nào sau đây đúng về phương pháp thực nghiệm?
\choiceTF
{Phương pháp thực nghiệm luôn luôn bắt đầu bằng việc tiến hành thí nghiệm.}
{\True Nếu kết quả thí nghiệm không phù hợp với dự đoán, ta cần điều chỉnh dự đoán hoặc vấn đề nghiên cứu.}
{Phương pháp thực nghiệm chỉ áp dụng được cho các hiện tượng cơ học.}
{\True Galileo Galilei là người có đóng góp to lớn trong việc xây dựng phương pháp thực nghiệm.}
\loigiai{
\begin{itemchoice}
\item (Sai) Phương pháp thực nghiệm luôn luôn bắt đầu bằng việc tiến hành thí nghiệm. (Vì): Bắt đầu bằng xác định vấn đề. B. Đúng. C. Sai - Áp dụng cho nhiều lĩnh vực Vật lí. D. Đúng
\item (Đúng) Nếu kết quả thí nghiệm không phù hợp với dự đoán, ta cần điều chỉnh dự đoán hoặc vấn đề nghiên cứu.
\item (Sai) Phương pháp thực nghiệm chỉ áp dụng được cho các hiện tượng cơ học.
\item (Đúng) Galileo Galilei là người có đóng góp to lớn trong việc xây dựng phương pháp thực nghiệm.
\end{itemchoice}
}
\end{ex}
